\documentclass[lang=cn,11pt]{elegantbook}

\title{道德与法治}
\subtitle{OCR整理稿}

\author{夏同}
\institute{华南理工大学}
\date{\today}
\version{1.0}
\bioinfo{说明}{由道法.pdf OCR整理}

\extrainfo{OCR初稿，仍需人工校对}

\setcounter{tocdepth}{3}

\logo{logo-blue.png}
\cover{cover.jpg}

\usepackage{array}
\usepackage{mathtools}
\IfFileExists{esint.sty}{
  \usepackage{esint}
}{
  \providecommand{\oiint}{\iint}
}
\ExplSyntaxOn
\tl_gset:Nn \CJKttdefault { ebtt }
\ExplSyntaxOff
\IfFontExistsTF{FangSong}{
  \setCJKmonofont[BoldFont={FangSong},ItalicFont={FangSong},BoldItalicFont={FangSong}]{FangSong}
  \setCJKfamilyfont{ebfs}[BoldFont={FangSong},ItalicFont={FangSong},BoldItalicFont={FangSong}]{FangSong}
  \renewcommand*{\fangsong}{\CJKfamily{ebfs}}
}{}
\IfFontExistsTF{KaiTi}{
  \setCJKfamilyfont{ebkai}[BoldFont={KaiTi},ItalicFont={KaiTi},BoldItalicFont={KaiTi}]{KaiTi}
  \renewcommand*{\kaishu}{\CJKfamily{ebkai}}
}{}
\IfFontExistsTF{Segoe UI Symbol}{
  \newfontfamily\daofasymbolfont{Segoe UI Symbol}
  \newcommand{\daofasym}[1]{{\daofasymbolfont #1}}
}{
  \newcommand{\daofasym}[1]{#1}
}
\newcommand{\ccr}[1]{\makecell{{\color{#1}\rule{1cm}{1cm}}}}
\newenvironment{shortexenum}
  {\par\noindent\setlength{\tabcolsep}{0pt}\renewcommand{\arraystretch}{1.15}%
   \begin{tabular}{@{}>{\raggedright\arraybackslash}p{0.485\linewidth}@{\hspace{0.03\linewidth}}>{\raggedright\arraybackslash}p{0.485\linewidth}@{}}}
  {\end{tabular}\par}
\newcommand{\shortitem}[2]{#1.\enspace #2}

\begin{document}


\section*{PDF第1页}
\addcontentsline{toc}{section}{PDF第1页}

\begin{note}[手写批注]
页面顶部有手写批注，疑似补充核心素养或法治观念相关内容；（此处有手写笔记但无法识别）。\par
\end{note}

\subsection*{七上}
\addcontentsline{toc}{subsection}{七上}

\textbf{关键词：中学时代}\par

1.中学时代对个人成长的意义？(p4)\par

\daofasym{①}中学时代是人生发展的一个新阶段为我们的一生奠定重要基础：\par

\daofasym{②}中学时代见证着一个人从少年到青年的生命进阶】为我们的人生长卷打上丰富而厚实的\par

底色。\par

\textbf{2.中学生活带给我们哪些成长的礼物？}\par

\daofasym{①}中学生活为我们的发展提供了多种机会；\par

\daofasym{②}新的目标和要求激发着我们的潜能，激励着我们不断实现自我超越；\par

\daofasym{③}在新的环境中重新塑造一个我\par

\textbf{关键词：梦想}\par

1.梦想的作用？（青少年为什么要有梦想？）（p9）\par

编织人生梦想，是青少年时期的重要生命主题。（原因)\par

\daofasym{①}它能不断激发我们生命的热情和勇气，让生活更有色彩。有梦想，就有希望；\par

\daofasym{②}有了梦想，人类社会才能不断地进步和发展；\par

2少年应怎样确立自己的梦想？（少年梦想的特点）(p10)\par

\daofasym{①}少年的梦是天真无邪，美丽可爱的。\daofasym{②}少年的梦要与个人的人生目标紧密相连；\daofasym{③}少年\par

的梦想要与时代的脉搏紧密相连，与中国梦密不可分。\par

3.怎样让梦想成为现实？（我们怎样去实现自己的梦想？）（p11-13)\par

\daofasym{①}实现梦想，不应止于心动，更应在于行动。努力，是梦想与现实之间的桥梁\par

\daofasym{②}努力需要立志： \daofasym{③}努力，需要坚持： \daofasym{④}努力也需要方法。\par

（分值大的时候需要解释）\par

4、中国梦：中华民族的伟大复兴。中国梦的基本内涵：实现国家富强、民族振兴、人民幸\par


\section*{PDF第2页}
\addcontentsline{toc}{section}{PDF第2页}

福。实现途径：必须走中国道路，必须弘扬中国精神，必须凝聚中国力量。\par

同长翼过统，红识更度快.们题非士核\par

\textbf{关键词；学习为53 幼身学化会中、长展}\par

1. 如何正确认识初中阶段的学习？(p17)\par

\daofasym{①}（内容上）初中阶段的学习，包括知识的获取，（能力的培养以及如何做人；\par

范围上)(学习不仅仅局限在学校，我们所看，所听，我们所触，我们所做，都可以是\par

学习。\par

3（表现上）学习不仅表现为接受和掌握，而且表现为探究、发现、体验和感悟；\par

\daofasym{④}（态度上）学习需要自觉、主动的态度：\par

\daofasym{③}（过程上）学习伴随着我们的成长\par

\daofasym{③}又时间上）学习没有终点。我们终生都在学习，终生都需要学习。我们要树立终身学习\par

的理念。\par

\daofasym{①}存实践中学习。我们始终不能停止学习的步伐，在生活和工作中学习，主动服务社会。\par

我们要重视实践，积极参加各类社会实践活动，努力做到知行合一。必须树立终身学习的\par

理念，养成主动学习的习惯，增强自我更新、学以致用的能力。\par

2.学习的重要意义？为什么要学习？（p18-19）\par

\daofasym{①}学习，不仅让我们能够生存，而且可以让我们有更充实的生活：\par

\daofasym{②}学习点亮我们心中的明灯，激发前进的动力；\par

\daofasym{③}在学习中，我们分享生命经验，获得成长，同时也可以帮助他人，服务社会，为幸福生\par

活奠基。\par

\daofasym{④}学习是我们的权利，也是我们的责任和义务应成为我们的一种生活方式，\par

【注意】这里可以链接到教育的作用：\daofasym{①}一个民族创新能力的提高离不开创新人才的培养。百年大\par

计,教育为本。\par

\daofasym{②}教育是民族振兴、社会进步的基石（非常非常非常喜欢考选择）,是提高国民素质、培养创新型人才、\par

促进人的全面发展的根本途径。\par

\daofasym{③}教育为人们的幸福生活奠基，教育寄托着亿万家庭对美好生活的期盼。\par


\section*{PDF第3页}
\addcontentsline{toc}{section}{PDF第3页}

3. 体味学习\par

学习中有快乐，学习也有辛苦，学习是苦乐交织的过程。（启示：珍惜受教育的机会，享\par

受学习的苦与乐。）\par

4.如何学会学习、享受学习？（p22-23)\par

\daofasym{①}学会学习，需要发现并保持对学习的兴趣\par

\daofasym{②}学会学习，需要掌握科学的学习方法（做笔记(）预习人合理安排学习时间等）\par

\daofasym{③}学会学习，还意味着要善于运用不同的学习方式。例如自主学习合作学习探究学习\par

\textbf{关键词：自己}\par

1.我们为什么要正确认识自己？（正确认识自己的意义？）（P27)\par

\daofasym{①}正确认识自己，可以促进自我发展；\par

\daofasym{②}正确认识自己，可以促进与他人的交往。（分值多加解释）\par

2.认识自己的途径和方法（如何正确认识自己？）(P28-29)\par

他明药\par

\daofasym{①}从生理心理、社会等方面认识自己，(全面客观地认识自己\par

\daofasym{③}通过他人评价认识自己。（更为客观、完整、清晰的认识）\par

3.我们应如何正确对待他人的评价？（P30）\par

\daofasym{①}我们要重视他人的态度与评价，但也要客观冷静分析，既不能盲丛，也不能忽视。用理\par

性的心态面对他人的评价，是走向成熟的表现。\daofasym{②}正确对待他人评价的方法：用心聆听\par

勇于面对，平静拒绝。\par

\textbf{4.如何接纳自己？}\par

\daofasym{①}接纳自己，需要接纳自己的全部；\par

\daofasym{②}接纳自己，需要乐观的心态，勇气和智慧。\par


\section*{PDF第4页}
\addcontentsline{toc}{section}{PDF第4页}

\textbf{5.如何欣赏自己？}\par

\daofasym{①}欣赏自己的独特、优点、努力和奉献；\par

\daofasym{②}欣赏自己不是骄傲自大、也不是目中无人，是成长道路上面对压力与挫折的自我鼓励与\par

自我奋进：\par

\daofasym{③}欣赏自己要善于向他人学习、与他人合作，成为一个内涵更加丰富的我，学会欣赏他人。\par

6.我们如何做更好的自己？(P34-35)\par

\daofasym{①}做更好的自己，就要扬长避短； \daofasym{②}做更好的自己，需要主动改正缺点：\par

\daofasym{③}做更好的自己，需要不断激发自己的潜能。（怎样激发潜能？P35）\par

\daofasym{④}做更好的自己，是在和他人共同生活的过程中不断成长的，更在为他人，为社会带来福\par

社的过程中实现的。\par

\textbf{关键词：友谊}\par

1、朋友的重要作用？(P41-43）【注意：友谊不能起决定作用】\par

\daofasym{①}朋友对一个人影响很大，我们的言谈举止、兴趣爱好甚至性格等都或多或少受朋友的影\par

响。\par

\daofasym{②}月朋友见证了我们一起走过的成长历程，我们需要真诚友善的朋友。\par

\daofasym{③}朋友丰富了我们的生活经验，友谊让我们体悟到生命的美好。（多做情景题理由）\par

2 友谊的特质?(P44-46)\par

\daofasym{①}友谊是一种亲密的关系；\daofasym{②}友谊是平等的，双向的；\daofasym{③}友谊是一种心灵的相遇。\par

3.我们要澄清哪些对友谊的认识？（如何正确看待友谊？）（P46-48）\par

\daofasym{①}友谊不是一成不变的；\daofasym{②}竞争并不必然伤害友谊，关键是我们对待友谊的态度。在竞争\par

中做到坦然接受并欣赏朋友的成就，做到自我反省和激励。我们要相信自己的独特，也要\par

相信朋友的成功并不意味着自己的失败。\par

\daofasym{③}友谊不能没有原则。友谊需要信任和忠诚，这并不等于不加分辨地为朋友做任何事。当\par

朋友误入歧途，不予规劝甚至推波助澜，反而会伤害友谊。坚持原则是对友谊的负责。\par


\section*{PDF第5页}
\addcontentsline{toc}{section}{PDF第5页}

4.如何建立友谊?,(P49-50)\par

\daofasym{①}建立友谊，需要开放自己。敞开心扉，主动表达\daofasym{②}建立友谊，需要持续的行动\par

5. 如何呵护友谊？(P51-53)\par

\daofasym{①}呵护友谊，需要用心关怀对方\}体会朋友的需要，表达关心与支持。\par

\daofasym{②}呵护友谊，需要学会尊重对方。把握好彼此界限与分寸，需要坦诚相待，但不意味着毫\par

无保留；\par

\daofasym{③}呵护友谊，需要学会正确处理冲突。相互协商，寻找彼此能接受的解决方式。\par

（保持冷静坦诚交流及时处理承担责任换位思考）\par

\daofasym{④}呵护友谊，需要学会正确对待交友中受到的伤害。我们可以选择宽容对方或结束这段友\par

谊。两者都需要勇气与智慧。\par

\textbf{关键词：网上交往}\par

1.网络交往有什么特点？（P55）\par

网络上的交往具有虚拟、平等、自主等特点。\par

2.说说网络交往的利和弊。为什么说网络是把双刃剑？（如何认识网络交往？）（重点掌握）\par

(p55-58)\par

(1)利：网络超越时空限制，开辟了人际交往的新通道，让我们有机会结交新的伙伴，拓展\par

交往圈。\par

网上交往可以满足我们的一些心理需要。\par

(2)弊：\daofasym{①}网络有时却关闭了与他人沟通的心灵之门：\par

虚拟的交往难以触摸到生活中的真实。\par

\daofasym{②}沉迷于网络交友会影响正常的学习和生活。\par

3.怎样慎重结交网友？（P56-58）\par

（1）在网络上交友需要考虑对学习和生活的影响，学会理性辨别、慎重选择；\par

（2）虚拟世界的交往带有很多的不确定性，我们要有一定的自我保护意识；（拒绝见网友\par


\section*{PDF第6页}
\addcontentsline{toc}{section}{PDF第6页}

的理由)\par

（3）将网上的朋友转化为现实中的朋友，需要慎重：\par

（4）在网上交往过程中，我们要遵守法律法规。\par

（5）我们要学会在现实中与同伴交往。增加真实而贴切的感受，为友谊奠定可靠的基础。\par

\textbf{4.网上交友，如何自我保护？}\par

\daofasym{①}对陌生人的邀请，不轻易接受；\par

\daofasym{②}对个人的家庭住址、经济状况、联系方式等，要注意保密；\par

\daofasym{③}遇到问题，要学会求助。\par

\textbf{拓展：我们应该怎样正确使用手机网络？}\par

\daofasym{①}自觉规范上网行为，遵守网络规则，有网络道德，做文明上网人；\par

\daofasym{②}对网上信息要有分辨能力，避免轻信盲从，提高抗诱惑能力。\par

\daofasym{③}提高防范意识，学会自我保护，不轻易泄露个人资料，不随意答应网友的要求。\par

\daofasym{①}学会自我管控，抵制不良诱惑，不沉溺于网络游戏\par

\textbf{关键词：师生情}\par

1.怎样理解教师的职业（教师有何作用？）P60-61\par

\daofasym{①}教师作为教育工作者，是人类文明的主要传承者之一；\par

\daofasym{②}教师是一种专门职业。教师是履行教育教学职责的专业人员，承担教书育人的使命。\par

\daofasym{③}教师是我们知识学习的指导者和精神成长的引路人。\par

2.新时代对老师的新要求？（四有好老师）P62\par

今天的教师要努力成为有理想信念、有道德情操、有扎实学识、有仁爱之心的好教师。\par

3.老师的风格不同产生的原因P63\par

由于年龄、学识、阅历、性格、情感与思维方式等差异，每位老师解决问题的方法和表达\par

方式不同，由此呈现出不同的风格。\par

4.如何正确对待不同风格的老师？P64\par


\section*{PDF第7页}
\addcontentsline{toc}{section}{PDF第7页}

\daofasym{①}承认老师之间的差异，接纳每位老师的不同\par

\daofasym{②}无论什么风格的老师都应该受到尊重\par

\daofasym{③}走进老师，更深入了解老师，发现不同风格老师的优点\par

\daofasym{④}了解老师教育行为的目的，与老师主动交往\par

5.怎样建立教学相长的师生关系？（我们怎样与老师交流？）P66\par

（1）面对老师的引领和指导，主动参与、勤学好问的态度有助于我们与老师相互交流。\daofasym{①}\par

真诚、恰当地向老师表达自己的观点和见解；\daofasym{②}与老师分享自己的学习感受、学习成果\par

（2）学会正确对待老师的表扬和批评。\par

6.如何正确对待老师的表扬和批评？P67\par

\daofasym{①}老师的表扬意味着肯定、鼓励和期待，激励我们更好地学习和发展；\par

\daofasym{②}老师的批评意味着关心、提醒和劝诫，可以帮助我们反省自己，改进不足；\par

\daofasym{③}正确对待老师的批评，我们要把注意力放在老师批评的内容和用意上，理解老师的良苦\par

用心。（情境题理由）\par

7.师生交往的良好状态是怎样的？（教学相长，亦师亦友）P68\par

学生乐于学习，老师寓教于乐，师生之间彼此尊重、相互关心、携手共进，是师生交往的\par

良好状态。\par

8.如何建立良好的师生关系？P68-69\par

\daofasym{①}彼此尊重，是我们与老师建立良好关系的开始，尊重老师的人格尊严、个性差异和劳动\par

成果等；\par

（尊重老师的原因见课本三点）\par

\daofasym{②}平等相待、相互促进，和老师成为朋友；\daofasym{③}主动关心体谅老师、理解老师；\par

\daofasym{④}保持冷静，积极与老师沟通，用恰当的方式化解师生矛盾。（题号5和8都是师生交往\par

综合运用)\par

9.与老师发生矛盾该怎么处理？（方法与技能P69）\par


\section*{PDF第8页}
\addcontentsline{toc}{section}{PDF第8页}

\daofasym{①}自我反思，冷静、客观地分析原因；\par

\daofasym{②}相信善意，多些宽容和理解；\par

\daofasym{③}坦诚相待，注意沟通方式，主动找老师交换意见；\par

\daofasym{④}求同存异，主动关心老师。\par

\textbf{10.我们为什么要尊重老师？}\par

\daofasym{①}尊重老师，是中华民族的优良传统，是我们作为“晚辈”的基本道德修养；\par

\daofasym{②}老师是我们知识学习的指导者，也是我们精神成长的引路人。\par

\daofasym{③}老师关心爱护我们，我们也要尊重老师。\par

\textbf{关键词：家/亲情}\par

1.家庭的含义：家庭是由婚姻关系、血缘关系或收养关系结合成的亲属生活组织。\par

家对我们有什么意义？P72-73\par

\daofasym{①}家是我们身心的寄居之所。我们的生命是父母给予的，我们的成长也离不开家庭的哺育\par

和支持；\par

家是我们心灵的港湾。家里有亲人，家中有亲情。亲情激励我们奋斗拼搏\par

2./中国人心中的“家”P74\par

家是代代传承、血脉相连的生活共同体，是甜蜜、温暖、轻松的避风港，\par

4.我们为什么要孝亲敬长？P74（相关法律：宪法、民法典）\par

\daofasym{①}在中国的家庭文化中，“孝”是重要的精神内涵;\par

\daofasym{②}孝亲敬长是中华民族的传统美德，也是每个中国公民的法律义务，（我国宪法规定：成\par

年子女有赡养扶助父母的义务；《民法典》规定成年子女对父母有养、扶助和保护的义\par

务)\par

\daofasym{③}我们的生命是父母给予的，我们的成长离不开家庭的哺育和支持。\par

5.如何孝亲敬长？（家庭美德和家风要关注）P74-75\par


\section*{PDF第9页}
\addcontentsline{toc}{section}{PDF第9页}

\daofasym{①}尽孝在当下。从现在开始，应该用行动表达孝敬之心。\par

\daofasym{②}尊敬双亲长辈，听取他们的意见和教导；\par

\daofasym{③}与双亲长辈保特亲近、融洽的关系，倾听他们的心声；经常和他们保持联系。\par

\daofasym{④}知恩、感恩，用行动表达感恩之情，如认真学习，心理想着父母体谅他们的辛劳，平日\par

里承担力所能及的家务劳动等。\par

6.为什么会产生“爱的冲突”？p79-80\par

\daofasym{①}独立与依赖并存，一方面，我们希望父母像对待成人一样尊重我们，不要过多干涉我们\par

的生活；另一方面，面对父母的信任和放手，有时我们又觉得失落和不安，期望得到更多\par

的关注和呵护\par

\daofasym{②}我们开始审视父母给予我们的爱，开始质疑父母，甚至挑战父母的权威与经验；\par

\daofasym{③}我们与父母在心智、学识、经历等方面差异较大，对问题的理解感受等方面必然存在差\par

异冲突难以避免。\par

7.处理不好亲子冲突的危害？逆反心理的危害？P81\par

亲子冲突就会伤害双方的感情，影响家庭的和睦，不利于双方的身心健康。\par

8.好何化解亲子冲突？P81-82\par

\daofasym{①}既需要父母做出榜样，也需要我们自己努力；\par

\daofasym{②}爱是需要呵护的，多与父母沟通，掌握互动沟通的技巧和应对冲突的智慧；\par

\daofasym{③}试着接纳父母的做法，理解父母行为中蕴含的爱；\par

\daofasym{④}尝试让父母了解我们的变化和需要，用他们能接受的方式表达我们的爱。\par

\daofasym{③}掌握和父母沟通的技巧（P82方法与技能）注意事实。把握时机。留意态度。选择方\par

式。考虑环境。\par

10.现代家庭发生了哪些变化？P83-84\par

\daofasym{①}现代家庭的结构、规模、观念等都发生了不同程度的变化； 了解核心家庭、主干家庭\par

等家庭结构。注意：核心家庭是两代人，主干家庭是三代人）\daofasym{②}家庭成员的交流、沟通方\par


\section*{PDF第10页}
\addcontentsline{toc}{section}{PDF第10页}

式发生了较大的变化；\par

3家庭氛围越来越平等、民主。\par

\daofasym{④}现代家庭生活的重要内容变化。关心世界和国家大事，探讨社会和人生问题）学习现代\par

科学文化知识，创建学习型家庭，参与社区活动，这些已成为现代家庭生活的重要内容。\par

(重精神）\par

12. 家庭中为什么会存在矛盾？（注意不是亲子矛盾）P85-86\par

\daofasym{①}在一个家庭中，祖辈、父辈、子辈之间有着不同的价值观念和生活方式，这些差异可能\par

带来家庭成员间的矛盾和冲突，影响家庭和睦。\daofasym{②}家庭成员的增加或减少，会带来人际关\par

系的变化，家庭成员的工作、身体和情绪状况等变化，也可能带来家庭氛围的变化。\par

13 如何创建和谐家庭？P85-86\par

\daofasym{①}家和万事兴。家庭成员之间相互理解信体谅积包容，可以增进理解，化解矛盾和\par

冲突;\par

\daofasym{②}家务我分担。积极参加家务劳动，养成劳动习惯，不断提高自我管理能力，增强家庭责\par

任意识。为和谐家庭作贡献。\par

\daofasym{③}以良好的心态面对家庭发展中的变化和问题。让亲情更浓，让我们的家庭更和睦。\par

苹女花化质文们了美恋\par

14、如何看待家规、家风、家 化自信\par

压然V精种\par

（1）作用：\daofasym{①}家规、家训、家风是中华民族传统文化的载体）也是中华传统美德的表现形\par

式。\daofasym{②}优秀的家规、家训、家风是建设和谐家庭（美好社会的重要精神支撑。\daofasym{③}传承优良\par

的家规、家风、家训是我们坚定文化自寝建设文化强国的需要。\par

（2）正确对待：对于家规、家训、家风，我们要结合时代发展，与时俱进，不断丰富其内\par

涵，取其精华，去其糟粕，践行社会主义核心价值观\par

（3）举例子：好的家风家训：一粥一饭，当思来之不易；玉不琢，不成器；诚实守信，堇\par

俭节约，孝亲敬长等\par

\textbf{关键词：生命}\par


\section*{PDF第11页}
\addcontentsline{toc}{section}{PDF第11页}

.生命有何特点？P89-90\par

生命来之不易，生命是独特的，生命是不可逆的;生命也是短暂的\par

2.向死而生的意义？（如何正确看待死亡？）P91\par

\daofasym{①}我们每个人都难以抗拒生命发展的自然规律。\daofasym{②}死亡是人生不可避免的归宿，它让我\par

们感激生命的获得。\daofasym{③}从容面对生命的不可预知，更加热爱生命，热爱生活。把有限的生\par

命投入到无限的奋斗和奉献中去。\par

3.如何理解生命有接续（生命的接续有什么意义？）P92\par

\daofasym{①}一代又一代的个体生命实现了人类生命的接续。在人类生命的接续中，我们总能为自己\par

的生命找到一个位置，担当一份使命。\daofasym{②}生命的接续，使得每个人的生命不仅仅是“我”\par

的生命，还是“我们”的生命。我们能更好地认识和面对自己的生命。\daofasym{③}在生命的接续中,\par

人类生命不断发展，人类的精神文明也不断积累和丰富。在精神上不断继承和创造人类的\par

文明成果\par

3.为什么要敬畏生命\}怎样理解生命至上？P94-96\par

\daofasym{①}生命是脆弱的、艰难的；生命是坚强的、有力量的，生命是崇高的、神圣的；我们对生\par

命要有一种敬畏的情怀。\daofasym{②}对生命怀有敬畏之心时，我们就会珍视它，我们的生命都是宝\par

贵的，生命价值高于一切：\par

\daofasym{③}生命至上，并不意味着只看到自己生命的重要性，我们也必须承认别人的生命同样重要\par

5.如何敬畏生命？P96-98\par

\daofasym{①}敬畏生命，让我们从对自己生命的珍惜走向对他人生命的关怀；\par

\daofasym{②}不漠视自己的生命，也不漠视他人的生命，尊重、关注、关怀和善待身边的每一个人；\par

\daofasym{③}对生命的敬畏并不是谁的命令，而是内心的自愿选择。\par

6.如何守护生命？（两方面：爱护身体和养护精神）P100\par

(1）爱护身体\par

\daofasym{①}守护生命首先要关注自己的身体。关心身体的状况，养成健康的生活方式；\par


\section*{PDF第12页}
\addcontentsline{toc}{section}{PDF第12页}

\daofasym{②}关注自己的内在感受，学会爱惜自己的身体；\par

\daofasym{③}增强安全意识、自我保护意识，提高安全防范能力，掌握一些基本的自救自护方法\par

（2）养护精神（经典诵读等活动意义见课本）\par

\daofasym{①}我们的精神发育，需要物质的支持，但不完全受物质生活条件和外部环境的制约。守住\par

自己的心灵，努力追求真善美；\par

\daofasym{②}自觉接受中华优秀传统文化的滋养\par

\daofasym{③}守护精神家园，在个人精神世界的充盈中发扬民族精神。\par

补充：精神发育和物质生活的关系？ \daofasym{①}精神发育需要物质的支持，但不完全\par

受物质生活条件、外部环境的制约。\daofasym{②}在物质贫乏、外部环境艰苦的情况下，只要\par

我们守住自己的心灵，仍然可以看到真、善、美。\daofasym{③}过度注重物质容易使我们丧失\par

对真、善、美的体验，丢失精神世界的财富。\par

【题头小修：第8题挫折的影响】\par

8.挫折的影响（如何正确认识挫折？）P107\par

消极影响：面对挫折，我们可能会产生负面情绪，但如果一味沉浸在负面情绪中，我们就\par

容易消沉，甚至做出不恰当的行为。我们需要及时调整自己，正确看待挫折。\par

积极影响：生活中的挫折是我们生命成长的一部分。得意时，挫折会使我们更清醒，避免\par

盲目乐观、精神懈怠；失意时，挫折会使我们获得更加丰富的生活经验。（补充一点见课\par

本)\par

9.如何正确面对挫折，发掘生命的力量，增强生命的韧性？P107-109\par

\daofasym{①}面对挫折，我们需要发现、挖掘自己的生命力量。每个人的生命都蕴含一定的承受力、\par

自我调节和自我修复的能力；\daofasym{②}培养自己面对困难的勇气和坚强的意志；\par

\daofasym{③}借助外力，学会与他人建立联系，向他人寻求帮助，获得他人的支持和鼓励。\daofasym{④}坚持目\par

标，不断努力。+具体方法（增强生命韧性的方法有哪些？）欣赏、培养幽默感；，和自己\par

信任的人谈一谈；考虑并接受最糟糕的结果，关心，帮助他人；培养某方面的兴趣。\par

10.怎样的一生是值得过的？（什么是有意义的生命？）P113\par

Z1\par


\section*{PDF第13页}
\addcontentsline{toc}{section}{PDF第13页}

\daofasym{①}能够活出自己的人生，实现自我价值： （自己）\par

\daofasym{②}当别人需要帮助时付出自己的爱心与承担责任； (他人）\par

\daofasym{③}将个人的理想与国家发展、民族复兴人类命运结合起来。 （人类）\par

11.发现我的生命（如何发现和创造生命的意义？）P113-115\par

\daofasym{①}探索生命意义，是人类生命的原动力之一。（原因）\par

\daofasym{②}生命是独特的，生命意义是具体的。每个人的生活不尽相同，我们都是在自己的生活经\par

历中一点一点地建构自己，形成人们所说的“我的人生”\par

\daofasym{③}生命的意义需要自己发现和创造。（这句话要记牢）\par

13.如何让生命充盈？P116\par

\daofasym{①}热爱学习，乐于实践（在探索中扩展生活的阅历（让生命充满色彩与活力\par

\daofasym{②}敞开自己的胸怀，不断尝试与他人、与社会、与自然建立联系，生命中的道德体验就会\par

不断丰富，对生命的感受力、理解力就会不断增强。\par

\daofasym{③}具体方法：阅读好书；专注地做自己热爱和有兴趣的事情，参加公益活动，帮助更多的\par

人等。(p116 探究与分享)\par

14.如何传递生命的温暖？（生命怎样相互关切？）P117\par

\daofasym{①}人在相互关切中感受温暖，传递温暖，拒绝冷漠。\daofasym{②}我们不仅要关注自身的发展，而且\par

要关切他人的生命，设身处地地思考并善待他人。\daofasym{③}让我们用真诚、热情、给予去感动、\par

改变他人／消融冷漠，共同营造一个互信、友善、和谐的社会。\par

15.怎样在平凡中创造伟大？P118-120\par

\daofasym{①}每个人的生命都有自己独特的使命。伟大在于创造和贡献人生的真正价值在于对社会\par

的责任与贡献。一个人的伟大，在于他能够运用自身品德、才智和劳动，创造出比自己是\par

有限的生命更长久的、不平凡的社会价值；\par

\daofasym{②}面对生活的艰难考验，不放弃、不懈怠，为家庭的美好和社会的发展贡献自己的力量，；\par

用认真(勤劳善良坚持(责任勇敢书写自己的生命价值：\par


\section*{PDF第14页}
\addcontentsline{toc}{section}{PDF第14页}

\daofasym{③}生命虽然平凡，但也能时时创造伟大。当我们将个体生命和他人的、集体的、民族的、\par

国家的甚至人类的命运联系在一起时，生命便会平凡中闪耀出伟大。\par

16.如何活出生命的精彩？（三点：拒绝贫乏充盈生命：拒绝冷漠温暖他人；平凡中创造伟大）\par

(1)拒绝贫乏，充盈生命；\par

(2)拒绝冷漠，关爱他人；\par

(3)在平凡中创造伟大。（综合13、14、15题中要点展开）\par

\textbf{拓展：作为中学生，怎样提升自己的生命的价值，创造有意义的人生？}\par

\daofasym{①}树立远大理想和的人生目标，确定人生的意义\par

\daofasym{②}在劳动、学习和奉献中创造价值，实现自身价值。\par

\daofasym{③}积极投身于社会实践活动中，增强社会责任感和服务社会的精神\par

\daofasym{①}从现在做起，从日常生活小事做起，为他人、社会和国家多做贡献。\par

\daofasym{③}珍爱生命、敬畏生命，守护生命，增强生命的韧性，活出生命的精彩。\par

\subsection*{七下}
\addcontentsline{toc}{subsection}{七下}

\subsection*{第一单元 青春时光（自信，自强，自尊，自立）}
\addcontentsline{toc}{subsection}{第一单元 青春时光（自信，自强，自尊，自立）}

\textbf{1.进入青春期，我们的身体变化主要表现在哪些方面？}\par

青春期我们的身体变化表现在三个方面：身体外形的变化，内部器官的完善，性机能的成熟。\par

\textbf{2.青春期的生理变化带给我们的积极影响是什么？}\par

青春期的生理变化带给我们旺盛的生命力，使我们的身体充满能量。我们拥有充沛的精力、敏捷的思维，对成长\par

满强烈渴望，感觉生活拥有无限可能。\par

\textbf{3.青春期的我们身体发育情况为什么各不相同？怎样正确看待这些不同？}\par

(1)原因：受遗传、营养、锻炼等因素的影响，我们身体的发育情况各不相同：有的长得快，有的长得慢；有的先长胖\par


\section*{PDF第15页}
\addcontentsline{toc}{section}{PDF第15页}

\begin{note}[手写批注]
页面右侧和正文周围有多处手写批注；（此处有手写笔记但无法识别）。\par
\end{note}

有的先长高。\par

(2)态度\daofasym{①}我们要正视身体的变化，欣然接受青春花蕾的绽放。\daofasym{②}不因自己的生理变化而自卑，是我们对自己的尊\par

重；不嘲弄同伴的生理变化，是我们对同伴的尊重\daofasym{③}在追求形体、仪表等外在美的同时，也要提高品德和文化修养\par

体现青春的内在美\par

\textbf{4.进入青春期，我们的矛盾心理是怎样产生的？主要表现有哪些？}\par

(1)产生：伴随着生理发育，我们的认知能力得到发展，自我意识不断增强，情感世界愈加丰富...这些变化让我们感到\par

新奇，也使我们产生矛盾和困惑。\par

(2)主要表现：\daofasym{①}反抗与依赖\daofasym{②}闭锁与开放\daofasym{③}勇敢与怯懦。\par

\textbf{5.青春期矛盾心理的积极影响有哪些？}\par

(1)有时让我们烦恼，但也为我们的成长提供了契机。\par

(2)积极面对和正确处理这些心理矛盾，我们才能健康成长\par

\textbf{6.怎样调节青春期的矛盾心理？}\par

(1)参加集体活动，在集体的温暖中放松自己。\par

(2)通过求助他人，学习化解烦恼的方法。\par

(3)通过培养兴趣爱好转移注意力，接纳和调适青春期的矛盾心理。\par

(4)还可以学习自我调节，成为自己的“心理保健医生”。\par

\textbf{7.怎样正确认识思维的独立（思维的独立要求我们如何做）？}\par

思维的独立并不等同于一味追求独特（而是意味着不人云亦云，有自己独到的见解，同时接纳他人合理、正确的意\par

见。\par

8.如何认识思维的批判性？为什么？怎么做\par

是什么一一在我们的成长中，与思维的独立性相伴随的是思维的批判性。思维的批判性（表现在对事情有自己的看法，\par

并且敢于表达不同观点，敢于对不合理的事情说“不”，敢牙向权威挑战“学贵有疑，在学习中需要批判的精神和\par

勇气.\par

为什么一一作用(意义)：\daofasym{①}有助于我们发现问题、提出问题，并从不同角度去思考问题，探索解决方案；\daofasym{②}能调动我们\par


\section*{PDF第16页}
\addcontentsline{toc}{section}{PDF第16页}

的经验，激发我们新的学习动机，促使我们解决问题，改进现状。\par

怎么做一—：要有质疑的勇气，有表达自己观点、提出合理化建议的能力；\daofasym{②}要考虑他人的感受，知道怎样的批判\par

容易被人接受，更有利于解决问题\par

【注意】批判既包括对外部事物的质疑、批评，又包括对自己的反思。技巧：\daofasym{①}批判只针对事情本身，而不攻击他人。\par

\daofasym{②}批判要具有一定的建设性。\par

\textbf{9.青少年为什么要开发创造潜力？}\par

为什么一一(1)青春凝聚着动人的活力，蕴含着伟大的创造力，为我们的成长带来无限可能\par

(2)青春的我们思想活跃，感情奔放，朝气蓬勃，充满对未来的美好憧憬，拥有改变自己（改变世界的创造潜力\par

(3)青春韶华/我们要争当勤奋学习、自觉劳动勇于创造的好少年，这是时代的要求。\par

\textbf{怎么做一我们应怎样进行青春的创造？}\par

(1)敢于打破常规，追求创新生活的丰富多彩，开创前人未走之路。\par

(2)关注他人与社会，看重创造的意义和价值，做一名对国家和社会有用的创造者。\par

(3)创造离不开实践。青春的创造意味着用自己的智慧和双手去尝试、探索、实践，通过劳动改变自己，影响世界。\par

10.如何认识性别刻板印象2\par

(1)含义：人们对男性或女性角色特征的固有印象表明人们对性别角色的期望和看法\par

(2)影响：在社会中人们对性别的认识通常会受到性别刻板印象的影响。有时我们也会这样认识自己认识身边的同学\par

朋友。可以帮助我们了解自己与异性的不同特点，学会如何塑造自我形象（如何与异性相处，但是性别刻板印象也可能\par

在某种程度上影响我们自身潜能的发挥，我们在接受自己生理性别的同时，不要过于受性别刻板印象的影响。\par

\textbf{11.男生女生为什么要欣赏对方的性别优势？对我们有何启示？}\par

(1)原因：男生女生各自拥有自身的性别优势。欣赏对方的优势，有助于我们不断完善自己。\par

(2)启示：我们不仅要认识自己的优势，而且要发现对方的优势，不应因自己某一方面的优势而自傲，也不应因自己某\par

方面的欠缺而自卑，相互取长补短，相互帮助，相互学习，让自己变得更加优秀。\par

\textbf{12.怎样理解青春期的心理萌动？}\par

(1)含义：步入青春期，一股从未有过的心潮悄然涌动，带给我们一种特殊的情感体验，这是青春期的心理萌动。\par


\section*{PDF第17页}
\addcontentsline{toc}{section}{PDF第17页}

【漏识补录：第13题异性交往做法】\par

怎么做：异性之间的友谊，可能让人敏感、遭到质疑，但只要我们内心坦荡、言谈得当、举止得体，这份友谊就会成为我们青春美好的见证。\par

(2)表现：在异性面前，我们有自我表现的欲望，更加在意自己的形象，渴望得到异性的肯定和接受。\par

\textbf{13与异性交往（相处）的作用有哪些？如何交往？}\par

为什么一一(1)有助于我们了解异性的思维方式、情感特征。\par

(2)我们与异性交往，成为好朋友能从对方身上看到和学习某些优秀品质。\par

(3)与异性交往是我们成长的一个重要方面，也是对我们的考验。\par

为我们青春美好的见证。\par

\textbf{14.面对青春期的心理萌动，我们该如何认识和面对异性情感？}\par

(1)认识：\daofasym{①}相遇青春，我们心中开始萌发一些对异性朦胧的情感，这是青春成长中的正常现象。\daofasym{②}在与异性交往的过程\par

中，我们会因为对异性的欣赏、对美好的向往而愉悦，也容易把这种欣赏和向往理解为爱情。其实，这并不是真正的爱\par

情。\par

(2)正确态度：面对生活中可能出现的朦胧的情，我们应该慎重对待，理智处理。\par

\textbf{【注意】什么是真正的爱情？}\par

爱情需要具有爱他人的能力，这种能力包含自我的成熟、，道德的完善，也包含对家庭的责任。爱情需要彼此深入了\par

解，需要一定的物质基础和共同的生活理想，是强烈、稳定、专一的感情。真正的爱情包含尊重、责任、珍惜、平等和\par

自律。\par

\textbf{15.青春的探索为什么需要自信（自信的作用有哪些）？}\par

自信让我们充满激情。有了自信/，我们才能怀着坚定的信心和希望，开始伟大而光荣的事业。自信的人有勇气交往\par

与表达，有信心尝试与坚持，能够展现优势与才华，激发潜能与活力，获得更多的实践机会与创造可能\par

16.青春期我们要自强的原因及做法。\par

(1)原因：\daofasym{①}自强可以让我们更自信。我们的每一点进步，都是成长的足迹，印记着我们克服惰性、抵制诱惑、战胜自我\par

的努力。\daofasym{②}自强，让青春奋进的步伐永不停息。\par

(2)做法：\daofasym{①}相信自己，勇敢尝试，不断进步，才能体验成功带来的自信。\daofasym{②}不断克服自己的弱点，战胜自己、超越自己\par

是自强的重要内容。自强，要靠坚强的意志、进取的精神和持久坚持。\par


\section*{PDF第18页}
\addcontentsline{toc}{section}{PDF第18页}

【漏识补录：第17题题头】\par

\textbf{17.什么叫“行己有耻”？我们该怎样做到“行己有耻”？}\par

(1)含义：行己有耻，凡自己认为可耻的就不去做。\par

(2)做法：\daofasym{①}要知廉耻，懂荣辱；有所为，有所不为。\daofasym{②}需要我们有知耻之心，不断提高辨别“耻"的能力。在行动之前,\par

审查愿望；在行动之中，监督调节；在行动之后，反思效果与影响。\daofasym{③}要求我们树立底线意识，触碰道德底线的事情不\par

做，违反法律的事情坚决不做。\daofasym{③}需要我们磨砺意志，拒绝不良诱惑，不断增强自控力。\par

18.“止于至善”，是人的一种精神境界，怎样做到“止于至善”\par

(1)我们应该有良好的格调；有我们的“至善"追求。\par

(2)每个人都可以从点滴小事做起。“勿以恶小而为之，勿以善小而不为。”积少成多，积善成德。\par

(3)在生活中寻找“贤"将他们作为榜样。\par

(4)要求我们养成自我省察的习惯。检视自身的不足，不盲目自责，积极调整自己，通过自省和慎独，端正自己的行为。\par

\textbf{19.榜样的作用有哪些？}\par

榜样不仅是一面镜子，而且是一面旗帜。好的榜样昭示着做人、做事的基本态度，激发我们对人生道路和人生理\par

的思考，给予我们自我完善的力量。善于寻找好的榜样、向榜样学习、汲取榜样的力量，我们的社会、我们的国家才\par

变得更加美好。\par

\subsection*{第一单元：做情绪情感的主人}
\addcontentsline{toc}{subsection}{第一单元：做情绪情感的主人}

\textbf{20.什么是情绪？人的基本情绪有哪些？}\par

(1)情绪是一种综合表现出来的生理和心理状态，包括内在体验、外显表情和生理激活三种成分。\par

(2)基本情绪喜、怒、哀、惧\par

\textbf{21.影响情绪的因素有哪些？}\par

\daofasym{①}个人的生理周期\daofasym{②}对某件事情的预期\daofasym{③}舆论氛围\daofasym{④}自然环境\par

\textbf{22.情绪的作用有哪些？}\par

(1)各种各样的情绪丰富了我们的生活，影响着我们的观念和行动。\par

(2)积极作用:它可能激励我们克服困难、努力向上,\par

(3)消极作用：也可能让我们因为某个小小的挫败而止步不前。\par

页 5 18\par


\section*{PDF第19页}
\addcontentsline{toc}{section}{PDF第19页}

【漏识补录：第25题调节情绪的方法】\par

\textbf{25.为什么要调节情绪？调节情绪的方法有哪些？}\par

调节情绪的方法：\daofasym{①}改变认知评价；\daofasym{②}转移注意；\daofasym{③}合理宣泄；\daofasym{④}放松训练。\par

\textbf{23.青春期的情绪特点有哪些？产生的原因是什么？}\par

双妇以主观因差\par

(1)特点：@反应强烈：\daofasym{②}波动与固执：\daofasym{①}细腻性：\daofasym{①}闭锁性之\daofasym{③}表现性\par

(2)产生原因：进入青春期，随着身体发育和生活经验不断丰富，我们的情绪也发生着变化，表现出青春期的情绪特点。\par

\textbf{24.情绪的表达有何特点？对我们有何启示？}\par

(1)特点：情绪表达具有相通性或相互感染性，即使没有语言的交流）个大的表情、声调、姿态和动作所表达的情绪\par

也会影响周围的人。情绪表达不仅影响自己的身心健康，而且关乎人际交往。\par

(2)启示：在人际交往中，既要了解自己的情绪，接受它们，同时也要学会以恰当方式表达情绪。\par

\textbf{为什么要调节情绪？调节的方法？}\par

(1)保持积极的心态，享受喜悦和快乐，可以让我们的青春生活更加美好。\par

(2)适度负面情绪，可帮助我们适应突发事件，持续地处于负面情绪状态，则危害我们身心健康。调节情绪的方法有哪些？\par

\textbf{26.情感与情绪的关系是怎样的？}\par

(1)联系：\daofasym{①}情感与情绪紧密相关，伴随着情绪反应逐渐积累和发展。\daofasym{②}我们对某些人或者事物的情绪，随着时间的推移\par

形成比较稳定的倾向，就可能产生某种情感。\par

(2)区别\daofasym{①}情绪是短暂的、不稳定的，会随着情境的改变而变化；\daofasym{②}情感则是我们在生活中不断强化、逐渐积累的，相\par

对稳定\par

\textbf{2情感的种类有哪些？}\par

基础性情感如安全感。 高级情感：如道德感。\par

正面的体验爱的情感负面的体验：如恐惧感）两方混杂的体验：如敬畏感。\par

\textbf{28.情感的作用有哪些？}\par

(1)在社会生活中情感是人最基本的精神需求\par

(2)情感反映着我们对人和对事的态度、观念，影响我们的判断和选择，驱使我们做出行动。\par

(3)情感与我(们想象力、(创造力相关。丰富、深刻情感有助于我们更全面地观察事物，探索未知。\par

(4)情感伴随着我们的生活经历不断积累、发展，这正是我们生命成长的体现。狭隘的生活经验容易导致偏执的情感。在\par


\section*{PDF第20页}
\addcontentsline{toc}{section}{PDF第20页}

【漏识补录：第33题联结度理解】\par

（2）联结度理解：\daofasym{①}集体的联结度越高，个人感知到的集体温暖就越多。\par

\daofasym{②}集体的联结度通常与以下因素密切相关：成员间相互关联的程度，集体对成员的重要性，成员间相互交流的频率，成员对共同目标的共识程度，成员间的默契程度，集体存在的时间长短。\par

1\par

生活经验的不断扩展中，我们的情感才可能更加丰富、深刻，我们的情怀才可能更加宽广、博大。\par

29.美好情感的作用。\par

生活中美好的人和事物，让我们身心愉悦，逐渐丰富我们对生活、对人生的美好情感。这些情感表达着我们的愿望\par

促进我们的精神发展。\par

\textbf{30.怎样获得或体味到美好的情感？}\par

(1)情由心生，美好的情感是在人的社会交往、互动中自然引发的，不能强迫。\par

(2)我们可以通过阅读、与人交往、参与有意义的社会活动等方式获得美好的情感。\par

(3)学会承受一些负负面感受，善于将负面情感转变为成长的助力，也可以让我们从中获得美好的情感体验，不断成长。（4\par

创造正面的情感体验完成一项自己喜欢的活动;帮助他人；)走进博物馆或大自然;欣赏艺术作品。\par

\textbf{31.负面情感的意义是什么？}\par

(1)生活中某些负面的情感体验尽管不那么美好，但对于我们的成长也有意义。例如，羞耻感、焦虑感和挫败感等会带来\par

不舒服、不愉快的负负面感受。但是，体验负负面感受未必是件坏事，它可以丰富我们的人生阅历，使我们的生命变得更力\par

饱满丰盈。\par

(2)学会承受一些负面感受，善于将负面情感转变为成长的助力，也可以让我们从中获得美好的情感体验，不断成长。\par

32.怎样传递情感正能量？ (1)用自己的热情和行动来影响环境。\par

(2)我们的情感需要表达、回应，需要共鸣。在与他人的情感交流中，我们可以传递美好的情感，传递生命的正能量。\par

(3)在生活中不断创造美好的情感体验，在传递情感的过程中不断获得新的感受，使我们的生命更有力量，周围的世界\par

因为我们的积极情感多一份美好。\par

\subsection*{第三单元在集体中成长}
\addcontentsline{toc}{subsection}{第三单元在集体中成长}

\textbf{33.什么是集体？如何理解集体的联结度？}\par

(1)含义：集体是人们联合起来的有组织的整体。\par

间相互关联的程度，集体对成员的重要性，成员间相互交流的频率，成员对共同目标的共识程度，成员间的默契程度，\par

集体存在的时间长短。通常情况下，以上因素赋值越大，集体的联结度就越高。\par


\section*{PDF第21页}
\addcontentsline{toc}{section}{PDF第21页}

\textbf{34.怎样感受集体的温暖（我们为什么要过集体生活）？}\par

(1)我们每个人都有过集体生活的情感需要。\par

(2)在集体中，我们希望被认可和接纳，得到尊重和理解，获得安全感和归属感。作为集体的一员，我们也散发着自己的\par

光和热，彼此传递关爱和温暖。\par

(3)当集体取得成绩、受到表彰或奖励时，我们可以体验到集体荣誉感。这种荣誉感令我们骄傲、自豪，给我们温暖和力\par

量，激励我们不断前进。集体荣誉是集体成员共同奋斗的结果，是我们共同的荣誉。\par

\textbf{35.集体力量的来源是什么？}\par

集体的力量来源于成员共同的目标和团结协作\par

\textbf{36.个人力量与集体力量直接的关系是怎样的？}\par

个人对集体的影响：\par

\daofasym{①}个人的力量是分散的,但在集体中汇聚,就会变得强大；\daofasym{②}集个人的力量是有限的,但通过优化组合可以实现优势互\par

补,产生强大的合力；\daofasym{③}借助这种合力,我们得以完成许多单凭一己之力无法完成的事情；\par

集体对个人的影响：\par

\daofasym{③}集体的力量是强大的,在某种程度上可以影响甚至改变一个人；\daofasym{③}个人在集体生活中会不自觉地产生与集体要求相\par

致的态度和行为：\daofasym{③}集体有助王我们获得安全感和自信心,也有助于我们学习他人的经验,扩大视野,获得成长。\par

\textbf{37.为什么说集体生活可以涵养品格？}\par

(1)集体生活可以培养我们负责任的态度和能力。\par

(2)集体生活可以培养人际交往的基本态度和能力\par

\textbf{38.怎样在集体中涵养品格？}\par

(1)在集体中，每个人有不同的角色，承担不同的职责；我们在认真做事的过程中体现自己的价值，体验责任感，做有担\par

当的人。\par

(2)在集体中，每个人来自不同的家庭，有不同的生活经历和性格特点，我们在交往中要学会彼此接纳、尊重、理解和包\par

容，学会友好相处。\par

\textbf{39.怎样在集体中发展个性？}\par


\section*{PDF第22页}
\addcontentsline{toc}{section}{PDF第22页}

(1)集体生活为我们搭建起与他人、与周围世界交往的平台。在这个平台上，我们展示自己的个性，发展自己的个性，不\par

断认识和完善自我。\par

(2)每个人的个性特点不同。包容他人的不同，学习他人的优点，有助于我们完善个性。\par

(3)积极参与共同活动，把握机遇，自主发展，使自己的个性不断丰富。\par

\textbf{40.个人意愿与集体规则的关系是什么？}\par

(1)在集体中，每个人都有自己的意愿，集体又必须有一些共同的规则。\par

(2)当集体规则与我们的个人意愿一致，并且能够保障个人利益时，我们更乐于积极遵守和维护。\par

(3)我们有时候会感受到集体规则与我们的某些个性化需要之间存在矛盾甚至冲突。\par

\textbf{41.个人意愿与集体规则存在冲突的原因？}\par

集体规则与我们的某些个性化需要之间存在矛盾甚至冲突，这可能基于一方有不正当或不合理的要求，也可能是\par

个人和集体的需要不同。\par

\textbf{42.当个人意愿与集体规则发生矛盾和冲突时，该怎么做？}\par

(1)面对冲突，理解集体要求的合理性，反思个人意愿的合理性和实现的可能性，我们就可能找到解决冲突的平衡点。\par

(2)对于集体要求中存在的不合理因素，我们可以通过恰当的方式，提出积极的改进建议。\par

(3)当个人利益与集体利益发生冲突时，应把集体利益放在个人利益之上，坚持集体主义\par

43.怎样才能让集体的和声更美（更和谐）\par

(1)需要每个人尽力做好自己，遵守规则，以保持和声的和谐之美。\par

(2)对于集体要求中存在的不合理因素，我们要通过恰当的方式表达自己的意见，提出积极的改进建议。\par

(3)当个人利益与集体利益发生冲突时，应把集体利益放在个人利益之上，坚持集体主义。\par

(4)承认个人利益的合理性、保护个人正当利益，反对只顾自己、不顾他人的极端个人主义。\par

(5)在集体生活中要学会处理与他人的各种关系。\par

(6)无论个人之间有多大的矛盾和冲突，我们都应心中有集体，识大体、顾大局，不得因个人之间的矛盾做有损集体利益\par

的事情。\par

44.（重点知识）---个人利益与集体利益的关系\par


\section*{PDF第23页}
\addcontentsline{toc}{section}{PDF第23页}

个人利益和集体利益在本质上是一致的。（1）当个人利益与集体利益发生冲突时，\par

应把集体利益放在个人利益之上，坚持集体主义。（2）坚持集体利益，是在承认个\par

人利益的合理性、保护个人正当利益的前提下，反对只顾自己，不顾他人的极端个\par

人主义。\par

45.在不同的集体中，我们扮演不同的角色，承担不同的责任。\par

我们同时属于多个集体，每个集体都有自己的旋律。在不同的集体中，我们扮演不同的角色，承担不同的责任。\par

\textbf{46.个人节奏与集体旋律的关系是怎样的？}\par

(1)每个人都有自己的生活节奏，当自己的节奏与集体的旋律和谐时，我们就可以顺利地融入集体。\par

(2)当自己的节奏和集体的旋律存在差异时，为了保持旋律和谐，我们需要调整自己的节奏。\par

(3)当我们面对不同集体中的角色无法统一节奏时，角色之间的冲突就可能给我们带来烦恼。\par

\textbf{47.怎样排解由于角色冲突带来的烦恼（方法是什么)？}\par

(1)在排解角色冲突带来的烦恼时，我们通常会考虑自己更关注哪个集体，或在其中的角色和责任的重要性，也会考虑自\par

己的兴趣、爱好以及任务的紧迫程度等。\par

(2)当遇到班级、学校等不同集体之间的矛盾时，我们应从整体利益出发，自觉地让局部利益服从整体利益，个人利益服\par

从集体利益。\par

(3)我们在不断地调整自己的节奏中学习过共同生活，在解决不同集体的角色冲突中学习过集体生活，让自己更好地融入\par

集体，感受集体生活带给我们成长的快乐。\par

\textbf{48.怎样看待集体中的小群体？}\par

(1)形成：在集体生活中，一些志趣相投、个性相似，或者生活背景类似的同学，往往自觉或不自觉地形成小群体\par

(2)积极作用：\par

A.在小群体中，彼此相互接纳，相互欣赏，找到自己的位置，体会到归属感和安全感。\par

B.在小群体中，我们与同伴更容易相互理解、沟通，在与同伴的交往中学习交往，在与同伴的互学共进中增长才干。\par

C.在日常生活中，有些小群体往往以“集体”的面目出现，当小群体的节奏融入集体生活的旋律时，小群体成员就\par

能感受到集体生活的美好，更愿意积极参与集体的建设\par


\section*{PDF第24页}
\addcontentsline{toc}{section}{PDF第24页}

【漏识补录：第51题美好集体的作用】\par

在美好集体中，每个人都能在其中获得丰富的精神养料，拥有充实的精神生活，感受集体的关爱和吸引，凝聚拼搏向上的力量，坚定自己的生活信念。\par

（3）消极作用：\par

A.当小群体不能很好地融入集体生活时，其成员就会产生与小群体之外的其他同学的矛盾和冲突，甚至与集体的\par

同要求产生矛盾和冲突\par

B.小群体内成员之间的友谊如果沾染上江湖义气，这样的小群体往往会将自身利益置于集体利益之上，沦为小团\par

主义。\par

(4)如何面对(或应对)：面对这样的矛盾和冲突，我们需要“心怀一把尺子"，不为成见所“扰”，不为人言所“惑"\par

明辨是非，坚持正确的行为，坚持集体主义，反对小团体主义。\par

(5)意义：在集体生活中，我们面对矛盾和冲突，解决问题的过程，是我们学习过集体生活的经历，也是促进集体和谐\par

展、健康成长的过程。\par

49.集体愿景的含义：美好集体拥有共同的梦想，向往美好的前景，承担共同的使命，认同正确的价值观，形成一致的\par

标和追求，这就是集体的愿景。\par

50.集体愿景的作用：愿景是集体的精神动力之源，是推动集体发展的内驱力。共同的愿景引领集体成员团结一致，开\par

进取。\par

51美好集体的作用：美好集体是我们共同学习、共同生活的精神家园，引领我们成长。在美好集体中，每个人都能在！\par

52.美好集体的特点\par

(1)美好集体是民主的、公正的 (2)美好集体是充满关怀与友爱的。\par

(3)美好集体是善于合作的 (4)美好集体是充满活力的\par

重点知识：集体中的争与合作（辨析题）\par

（1）竞争是集体的动力，也是集体活力的重要表现，他是以承认和尊重为前提的。\par

（2）美好集体是善于合作的。合作意味着每个人都发挥自己最大的作用，同时又避免\par

人英雄主义。每个人都努力实现自己的价值，而不是消极依赖或者袖手旁观。集体成员讠\par

过合作交流互鉴，合作学习，共同提高。\par

（3）在竞争中合作，在合作中竞争会使我们的集体更强大，使我们每个人更快地进步。\par

页 5 24\par


\section*{PDF第25页}
\addcontentsline{toc}{section}{PDF第25页}

因此，在集体生活中，我们既需要合作，也需要竞争，互相学习，你追我赶。竞争是以承\par

认、尊重为前提，集体成员之间交流互鉴，合作学习，共同提高。\par

\textbf{53.如何在集体共建中尽责？}\par

(1)积极参与集体共同愿景和目标的制定，并予以理解和认同。\par

(2)积极参与共同商定集体的规则与制度内容，协商确定对组织领导者的品格与才能要求，通过民主程序选举乐于服务集\par

体、有责任心的组织领导者，在尊重不同意见的基础上努力达成共识。\par

(3)发扬自治”精神，主动参与集体建设，积极参加集体活动，自觉维护集体荣誉。\par

(4)共同创造良好的集体氛围,良好的人际关系、健康的舆论氛围、积极的精神气息，离不开每个人的努力\par

54.如何建设美好（班）集体？（集体需要每个人怎样去担当）\par

(1）集体建设需要每个人的智慧与力量。\par

(2)为集体出力，需要每个人从实际情况出发，各尽其能，发挥所长。集体的事务，无论大小都要认真对待，努力做好。\par

(3)集体荣誉是我们共同的利益和荣誉，需要我们悉心呵护。\par

(4)在集体中承担责任。\par

(5)在集体生活中，学会纳他人，理解和包容他人。\par

\textbf{55勇于承担责任的意义（为什么）？}\par

(1）承担责任既是个人有所成就的基础，也是集体发展的必要前提\par

(2)在集体生活中学会承担责任，是自我磨砺的过程。这个过程有助于我们学会正确地做事（提高能力/获得他人的认可\par

与尊重，扩大自我成长的空间。\par

(3）勇于担责可以为自己赢得信任，被赋予更大的责任，从而拥有更多发展的机会。\par

\textbf{56.我们应如何与集体共成长？}\par

（1）在共建中尽责。在集体建设中，我们要从实际情况出发，各尽所能，发挥特长，主动参与集体建设，积极参加集体\par

活动，自觉维护集体荣誉。共同确定集体的愿景和目标，共同商定集体的规则和制度内容，共同创造良好的集体氛围。\par

（2）在担当中成长。在集体生活中学会承担责任，在担当中为自己赢得更多的发展机会和更大的自我成长空间；学会接\par

纳他人，理解和包容他人：；学会关爱他人，互相帮助；学会参与，学会担当，从而为自己步入社会奠定坚实的基础。\par


\section*{PDF第26页}
\addcontentsline{toc}{section}{PDF第26页}

简单总结：个人成长与集体成长的关系\par

一方面，集体离不并个人。集体的荣誉是集体成员共同奋斗的结果，是集体共\par

同的荣誉。集体的建设需要每个人从实际出发，各尽其能，发挥所长。另一方\par

面，集体成就个人。集体给我们温暖和力量，为我们搭建成长的平台。在班复\par

体这个平台上，发展个性，提升能力，体现了自己的价值。所以，个人与集作\par

共成长，相互成就，我们要在积极承担集体责任申获得更夫发展。\par

\subsection*{第四单元走进法治天地}
\addcontentsline{toc}{subsection}{第四单元走进法治天地}

\textbf{57.怎样理解生活与法律息息相关？}\par

(1)法律就在我们身边，就在我们的生活中。我们在生活中形成的各种社会关系，以及由此产生的矛盾和纠纷，不仅需\par

依靠道德、亲情、友情来协调，而且需要法律来调整。每一部法律都是应生活的需要而制定和颁布，又对生活加以规\par

和调整。\par

(2)法律已经深深地嵌入我们的生活之中，渗透到社会的方方面面。作为社会关系的调节器，法律不仅服务于人们的当\par

生活，而且指导着人们未来的生活。\par

((3)法律与我们每个人如影随形，相伴一生。我们）生都享有法律规定的各项权利，同时必须履行法律规定的各项义\par

法律规定的权利和义务为我们每个人提供了自由生存和发展的空间。\par

58法律的本质？法律是统治阶级意志的体现，是用来统治国家、管理社会的工具，也是调整社会关系、判断是非\par

直、处理矛盾和纠纷的标尺。\par

\textbf{59.法治的含义和重要性是什么？}\par

(1)含义:就是依法对国家和社会事务进行治理，强调依法治国、法律至上，要求任何组织和个人都要服从法律，遵守\par

律，依法办事。\par

(2)重要性：\daofasym{①}法治是人们共同的生活愿景，也是国家治理现代化的重要标志。\daofasym{②}法治助推中国梦的实现，是实现政治\par

明、社会公平、民心稳定、国家长治久安的必由之路。\par

\textbf{60／建设社会主义法治国家的总目标和意义是什么？}\par

党的十八届四中全会提出了全面推进依法治国的总目标，即建设中国特色社会主义法治体系人建设社会主义法治\par

7-6\par


\section*{PDF第27页}
\addcontentsline{toc}{section}{PDF第27页}

【漏识补录：第65题法律普遍约束力】\par

（1）在法治社会里，公民在法律面前一律平等，任何人都没有超越法律的特权。\par

（2）每个公民都平等地受到法律的保护，平等地享有权利和履行义务。\par

家\par

意义:法治助推中国梦实现,是实现政治清明、社会公平、民心稳定、国家长治久安必由之路。\par

\textbf{61.道德、法律等行为规范的作用是什么？}\par

它们共同约束人们的行为，调整社会关系，维护社会秩序。\par

62.法律的特征。 ()是由国家制定或认可的。 (2)法律是由国家强制力保证实施的。（最本质的特征）\par

(3)法律对全体社会成员具有普遍约束力。\par

\textbf{63.法律区别于道德等行为规范的最主要特征是什么？}\par

法律的实施以强大的国家力量作后盾，而其他行为规范主要依靠社会舆论信念习俗、教育和道德、纪律等力量\par

保证实施。这是法律区别于其他行为规范的最主要特征。\par

\textbf{64.国家强制力是指什么？}\par

国家强制力主要包括军队、警察、法庭、监狱等。它的运用必须以合法为前提。\par

\textbf{65.如何理解“法律对全体社会成员具有普遍约束力”？}\par

等地享有权利和履行务；(3)任何人不论职务高低、功劳大小，只要触犯国家法律，都必须承担相应的法律责任。\par

66.法律的作用，\par

(1)法律规范全体社会成员的行为，保护着我们的生活，为我们的成长和发展创造安全、健康、有序的社会环境。\par

(2)法律规定我们享有的权利，应该履行的义务。法律也为我们评判、预测自己和他人的行为提供了准绳，指引、教育人\par

向善。\par

3)法律通过解决纠纷和制裁违法犯罪，惩恶扬善、伸张正义，维护我们的合法权益。\par

\textbf{67.未成年人为什么需要给予特殊保护？}\par

(1)未成年人身心发育尚不成熟，自我保护能力较弱，辨别是非能力和自我控制能力不强，容易受到不良因素的影响和不\par

法侵害，需要给予特殊的保护。\par

(2)未成年人的生存和发展事关人类的未来，对他们给予特殊关爱和保护，已经成为人类的共识。保护未成年人的合法权\par

益，是人类文明和社会进步的应有之义。\par


\section*{PDF第28页}
\addcontentsline{toc}{section}{PDF第28页}

\textbf{68.对未成年人实施特殊保护的法律有哪些？}\par

我国宪法和婚姻法、义务教育法等法律，都对保护未成年人作出了特别规定；我国还颁布了未成年人保护法、预防\par

未成年人犯罪法等专门法律，给予未成年人特殊关爱和保护。\par

\textbf{69保护未成年人合法权益的四道防线是什么？}\par

政府保护、社会保护、学校保护、家庭保护、司法保护、网络保护六位一体，构筑起保障未成年人合法权益的六\par

防线，形成全社会关心、保护未成年人的有效机制和良好风尚。\par

\textbf{70.家庭保护、（学校保护（社会保护（司法保护的含义是什么？}\par

(1)家庭保护是指父母或其他监护人对未成年人进行的保护，包括生活上的关心照顾和思想上的教育培养。\par

(2)学校保护是指学校、幼儿园和其他教育机构对未成年人实施的保护。\par

(3)社会保护是指国家、社会团体、企业事业组织以及其他组织和个人对未成年人实施的保护。（社会组织等）\par

(4)司法保护是指国家司法机关(包括公安机关、人民检察院、人民法院以及司法行政部门在内广义上的司法机关)依法\par

依法履行职责，对未成年人实施专门保护，是维护未成年人合法权益重要保障。（未成年人的违法犯罪，遗产继承、抚养权\par

问题）\par

（5网络保护是指国家、社会、学校和家庭对未成年人的网络生活实施的专门保护。（手机、网络）\par

（6）政府保护是指各级政府及其相关部门在各自职责范围内对未成年人实施的保护（有关的行政机关，比如教育局）\par

\textbf{71未成年人在享受特殊保护的同时，要注意什么？}\par

宪法和法律赋予未成年人受到特殊保护的权利，未成年人要珍惜自己的权利，依法行使自己的权利，同时要尊重和\par

维护他人的权利，自觉履行公民应尽的义务。\par

72.依法办事的重要性。\par

国无法不治，无民法不立，法律保障人们的幸福生活。法律保障功能的实现靠我们每一个人对法律的尊崇和遵守。\par

73.依法办事的要求。\par

(1)遵守各种法律法规。\daofasym{①}遇到问题需要解决，应当通过法治方式，表达自身合法的诉求和愿望。\par

\daofasym{②}在实现利益的过程中，还要自觉维护他人和集体的合法权益。\par

(2)要养成尊法学法守法用法的习惯，逐步成长为社会主义法治的忠实崇尚者、自觉遵守者、坚定捍卫者。\par


\section*{PDF第29页}
\addcontentsline{toc}{section}{PDF第29页}

\textbf{74.为什么要树立法治意识？}\par

(1)建设法治中国是中国人民的共同事业，人民既是法治的践行者，又是法治的受益者。\par

(2)人民权益要靠法律保障，法律权威要靠人民维护。\par

(3)推动全社会树立法治意识，增强全社会厉行法治的积极性和主动性，对于全面推进依法治国、建设社会主义法治国家\par

具有重要的意义。\par

75.怎样树立法律意识？ 法件不做，法降男水不做\par

(1)对法律发自内心的认可、崇尚、遵守和服从。\par

(2)不否认道德的重要性。人们道德水平的提高，有利于增强尊法守法的意识和自觉性；有利于促进法治生活方式的形成。\par

(3)青少年要增强法治意识，依法办事，成为法治中国建设的受益者、参与者、推动者。\par

\textbf{76.作为国家未来的建设者，青少年应为法治中国建设做出哪些努力？}\par

（1）在学习法律方面，认真学习法律知识，了解法律的基本内容，提高法律意识\par

（2）在遵守法律方面：法律要求的必须去做，法律提倡的积极去做，法律禁止的坚决不做。依法行使自己权利的同时，\par

尊重和保护他人和集体的合法权利，自觉履行法定义务\par

（3）在依法维权方面：当自己的合法权益受到侵害时，依法维护自己的合法权益，并积极运用法律帮助他人。\par

\textbf{77.如何保护未成年人？}\par

政府保护、社会保护、学校保护、家庭保护、司法保护、网络保护六位一体，构筑起保障未成年人合法权益的六道\par

防线，形成全社会关心、保护未成年人的有效机制和良好风尚。\par

人保纤(6第-7降也\par

（2）国家、社会、学校和家庭应当教育和帮助未成年人维护自身的合法权益，增强自我保护意识和自我保护能力\par

（3）依法行使自己权利的同时，尊重和保护他人和集体的合法权利，自觉履行法定义务。\par

\textbf{78、未成年人违法犯罪的原因有哪些？}\par

（1）未成年人自身：法律意识淡薄，道德水平不高，辨别是非能力和自我控制能力不强，缺乏责任意识，交友不慎等。\par

（2）外在角度：家庭、学校教育不到位，监管不到位。\par

79.国家出台“双减”政策的重要性。\par

【答案】有利于减轻学生的学习压力，促进学生全面发展：有利于规范校外培训机构的发展】营造良好的教育环境；有\par


\section*{PDF第30页}
\addcontentsline{toc}{section}{PDF第30页}

\begin{note}[手写批注]
页面顶部有手写批注；（此处有手写笔记但无法识别）。\par
\end{note}

利于促进教育均衡发展，维护和促进社会公平正义\par

117、谈谈为减轻学生的作业负担和校外培训负担，我们应该如何做？（分角度作答)\par

国家：强化管理，优化政策环境，引导校外培训机构规范发展：\par

教育机构：遵守法律规定和政策引导，规范办学，合理竞争；\par

学校：优化教学，合理安排课程，促进学生全面发展；\par

家长：主动与学校和老师沟通，构建和谐家校关系；更新教育理念，理性规划孩子未来发展方向；倾听孩子心声，扎\par

孩子合理利用课余时间；\par

学生：认真学习，充分利用好课堂时间，提高学习效率；加强自律，养成良好的学习和劳动习惯，不断促进自身发展\par

\subsection*{八上}
\addcontentsline{toc}{subsection}{八上}

\subsection*{第一单元 走进社会生活}
\addcontentsline{toc}{subsection}{第一单元 走进社会生活}

\subsection*{第一课 丰富的社会生活}
\addcontentsline{toc}{subsection}{第一课 丰富的社会生活}

1.当今社会生活的特点：绚丽多彩。\par

2.了解社会生活的方式\daofasym{★}\daofasym{★}\daofasym{★}\daofasym{★}\daofasym{★}\par

\daofasym{①}观看新闻；\daofasym{②}阅读报刊；\daofasym{③}参观访问；开展社会调查活动：\daofasym{③}参加社会公益活动等。\par

3.参与社会生活的意义\daofasym{★}\daofasym{★}\par

\daofasym{①}对社会生活的感受越来越丰富，认识越来越深刻：\par

\daofasym{②}会更加关注社区治理，并献计献策；会更加关心国家发展，或为之自豪，或准备为之分忧。\par

4.个人与社会的关系\par

(心为人离不开社会\par

\daofasym{①}义的身份是在社会关系中确立的。在不同的社会关系中，我们具有不同的身份。\par

\daofasym{②}个人的生存和发展离不开社会提供的物质支持和精神滋养。\par

\daofasym{③}人的成长是不断社会化的过程\par

（2）社会离不开个人\par

\daofasym{①}个人是社会的有机组成部分。\daofasym{②}社会中的物质资料和精神资料都是人们劳动实践产生的。\par

5.社会关系从建立的基础可分为：血缘关系、地缘关系（邻居等）和业缘关系（同学、同事、师生关系等）\par

\textbf{6.如何理解“人的成长是不断社会化的过程”？}\par

（1）含义：指一个人从最初的自然的生物个体转化为社会人的过程。\par

（2）途径：通过父母的抚育、同伴的帮助、老师的教诲和社会的关爱等。\par

（3）表现：我们的知识不断丰富，能力不断提高，规则意识不断增强，价值观念日渐养成。\par

（4）结果：我们逐渐成长为一名合格的社会成员。\par

7.亲社会行为\par

（）含义：指人们在社会交往中有利于他人和社会的行为（2）表现：谦让、分享、帮助他人、关心社会发展等\par

(3) 意义\daofasym{★}\daofasym{★}\daofasym{★}\daofasym{★}\daofasym{★}\par

\daofasym{①}利于我们养成良好的行为习惯，塑造健康的人格，形成正确的价值观念，获得他人和社会的接纳与认可。\par


\section*{PDF第31页}
\addcontentsline{toc}{section}{PDF第31页}

\begin{note}[手写批注]
页面右侧有手写批注，可辨识片段：打造清朗的网络空间；其余（此处有手写笔记但无法识别）。\par
\end{note}

\daofasym{②}积极融入社会，倾力奉献社会，有利于实现自己的人生价值。\par

\daofasym{③}青少年处于走向社会化的关键时期，我们应树立积极的生活态度，关注社会，了解社会，服务社会，养成亲社会行为。\par

(4) 要求\daofasym{★}\daofasym{★}\daofasym{★}\daofasym{★}\daofasym{★}\par

\daofasym{①}亲社会行为在人际交往和社会实践中养成。\daofasym{②}我们要主动了解社会，关注社会发展变化，积极投身于社会实践。\par

\daofasym{③}在社会生活中，我们要守社会规则和习俗，热心帮助他人。\par

\subsection*{第二课网络生活新空间}
\addcontentsline{toc}{subsection}{第二课网络生活新空间}

1. 网络的影响 \daofasym{☆}\daofasym{★}\par

网络是一把双刃有利有：（积极影响）\daofasym{①}网络丰富日常生活。信息传递和交流变得方便迅捷；打破了时空限制，\par

促进了人际交往；让生活变得更加便利和丰富多彩。（利）\par

\daofasym{②}网络推动社会进步。网络为经济发展注入新的活力；网络促进民主政治的进步；网络为文化传播和科技创新搭建新平\par

台。 (利)\par

（消极影响）网络信息良莠不齐；沉迷于网络会影响学习、工作和生活：个人隐私容易被侵犯。（）\par

2. 如何理性参与网络生活？ \daofasym{★}\daofasym{★}\daofasym{★}\daofasym{★}\daofasym{★}\par

\daofasym{①}要提高媒介素养，积极利用互联网获取新知、促进沟通、完善自我，培养面对多种信息时的选择、理解、质疑和评价\par

等能力\daofasym{②}要学会“信息节食”，浏览、寻找与学习和工作有关的信息。\daofasym{③}学会辨析网络信息，让谣言止于智者，自觉抵\par

制暴力、色情、恐怖等不良信息。\par

（注意：这个可以与上的网络交友结合）\par

3. 如何传播网络正能量？ \daofasym{★}\daofasym{★}\daofasym{★}\daofasym{★}\daofasym{★}\par

\daofasym{①}充分利用网络平台为社会发展建言献策，为决策科学化、民主化贡献自己的力量。\par

\daofasym{②}积极践行社会主义核心价值观，不断提高网络媒介素养，共同培育积极健康、向上向善的网络文化。\par

\textbf{4、如何建立安全、健康、充满正能量的网络空间？}\par

（1）国家：修改和完善有关互联网的法律法规，加大和完善执法力度，建立健全网络监管、预警机制。\par

（2）网络运营商（企业）：要依法经营，切实加强对网络平台和网络从业人员的监督和管理。\par

（3）青少年：\daofasym{①}要提高媒介素养，积极利用互联网获取新知、促进沟通、完善自我，培养面对多种信息时的选择、理解、\par

质疑和评价等能力\daofasym{②}要学会“信息节食”，浏览、寻找与学习和工作有关的信息。\daofasym{③}学会辨析网络信息，让谣言止于智\par

\daofasym{④}自觉遵守道德和法律，做一名负责任的网络参与者。\daofasym{③}学会在网络上传播正能量。\par

\textbf{5、（补充知识点）网络谣言的危害？}\par

\daofasym{①}扰乱人们的思想和行为\daofasym{③}破坏互联网的形象和公信力，影响互联网的健康发展\daofasym{③}污染网络环境，扰乱网络秩序\daofasym{①}引发\par

社会动荡，危害公共安全，损害公共利益，影响社会稳定等\par

\textbf{6、如何辨别网络谣言？}\par

\daofasym{①}注意信息出处\daofasym{②}关注官方信息\daofasym{③}对信息进行多方验证\daofasym{④}向他人求助等\par

\subsection*{第三课·社会生活离不开规则}
\addcontentsline{toc}{subsection}{第三课·社会生活离不开规则}

3.1维护秩序\par

1.社会秩序\par

（1）含义：社会秩序是社会生活的一种有序化状态。\par

（2）种类：社会管理秩序、生产秩序、交通秩序、公共场所秩序等。\par

（3）重要性：\daofasym{①}社会正常运行需要秩序，这样才能避免混乱、减少障碍、化解矛盾，提高社会运行效率，降低社会管\par

理成本。\daofasym{②}社会秩序是人民安居乐业的保障。在有序、整洁、安全的社会环境中，我们才能享有人身自由、财产安全和\par

公平的发展机遇等。\par

2．社会规则\par

（1）形成：社会规则是人们为了维护有秩序的社会环境，在逐渐达成默契与共识的基础上形成的。\par

（2）种类：道德、纪律、法律等。\par

(3）社会规则如何维护社会秩序？\daofasym{★}\daofasym{★}\daofasym{★}\daofasym{★}\daofasym{★}\par

\daofasym{①}社会规则明确社会秩序的内容。规则告诉我们如何作为社会一员享有权利，承担责任，处理好与他人、社会的关系。\par


\section*{PDF第32页}
\addcontentsline{toc}{section}{PDF第32页}

\daofasym{②}社会规则保障社会秩序的实现。社会规则明确了当有人破坏秩序时该如何处罚，从而保障社会的良性运行。\par

3.2 遵守规则\par

1.自由与规则的关系\daofasym{★}\daofasym{★}\daofasym{★}\daofasym{★}\daofasym{★}\par

\daofasym{①}社会规则划定了自由的边界，自由不是随心所欲，它受道德、纪律、法律等社会规则的约束。我国宪法规定，公民在\par

使自由和权利的时候，不得损害国家的、社会的、集体的利益和其他公民的合法的自由和权利。\par

\daofasym{②}社会规则是人们享有自由的保障。人们建立规则的目的是保证每个人不越过自由的边界，促进社会有序运行。违法规\par

扰乱积序的行为应当受到相应的惩罚。\par

2. 如何遵守规则？ \daofasym{★}\daofasym{★}\daofasym{★}\daofasym{★}\par

\daofasym{①}避守社会规则需要他律（监督、提醒、奖惩等外在约束）和自律（严于律已，自我反省，克服已经发现的不良行为\par

\daofasym{②}遵守社会规则，需要我们发自内心地敬畏规则，将规则作为自己行动的准则，内化于心、外化于行。\par

\textbf{3.如何维护与改进规则？}\par

\daofasym{①}我们要坚定维护规则。一方面要从自己做起,自觉遵守规则另一方面要在保护自身安全的前提下,提醒、监督、帮助\par

人遵守规则。\daofasym{②}我们要积极改进规则。我们要积极参与规则的改进和完善，善于与他人沟通交流、寻求共识，积极为亲\par

则的形成建言献策，使之更加符合人民的利益和社会发展的要求。\par

4. 为什么要改进规则？ 规则不是一成不变的。\daofasym{★}\daofasym{★}\daofasym{★}\par

\daofasym{①}随着社会生活的发展和社会生活的变迁，一些原来没有的规则，需要制定；\par

\daofasym{②}一些原有的规则失去了存在的合理性，需要废除；\par

\daofasym{③}一些原有的规则不能完全适应实际生活的变化，需要加以调整和完善。\par

\subsection*{第四课社会生活讲道德}
\addcontentsline{toc}{subsection}{第四课社会生活讲道德}

4. 1 尊重他人\par

(注意：(尊重是交往的起点。 选择题)\par

1.尊重的表现：尊重他人的人格、权利等。\par

2. 尊重他人的重要性\daofasym{★}\daofasym{★}\daofasym{★}\daofasym{★}\daofasym{★}\par

\daofasym{①}尊重他人是一个人内在修养的外在表现。\daofasym{②}每个人都是有尊严的个体,都希望得到他人和社会的尊重。\par

\daofasym{③}尊重使社会生活和谐融治。尊重是维系良好人际关系的前提，是文明社会的重要特征。\par

3.尊重他人的方法\daofasym{★}\daofasym{★}\daofasym{★}\daofasym{★}\daofasym{★}\par

\daofasym{①}积极关注、重视他人。考虑他人感受，认真对待他人，而不冷落、忽视他人\daofasym{②}平等对待他人。我们每个人在人格和\par

律地位上都是平等的。我们要发自内心的尊重他人人格，对所有人的都一视同仁。\daofasym{③}学会换位思考。\daofasym{④}学会欣赏他人\par

4.2 以礼待人\par

1.文明有礼的表现：主要表现在语言文明、仪表端庄、举止文明。如一个人的尊重、谦让、与人为善等良好品质。\par

2.文明有礼的意义\daofasym{★}\daofasym{★}\daofasym{★}\daofasym{★}\daofasym{★}\par

\daofasym{①}文明有礼是一个人立身处世的前提。（选择题）\daofasym{②}文明有礼促进社会和谐，\daofasym{③}文明有礼体现国\par

形象。\par

3.怎么做文明有礼的人？\daofasym{★}\daofasym{★}\daofasym{★}\daofasym{★}\daofasym{★}\par

\daofasym{①}态度谦和用语文与\daofasym{②}仪表整洁、举止端庄。\daofasym{③}需要在社会生活中不断学习、观察、思考和践行。\par

4.3诚实守信\par

1.为什么要讲诚信?梦*\daofasym{★}\daofasym{★}\par

\daofasym{①}诚信是社会主义核心价值观在公民个人层面的一个价值准则，是一种道德规范和品质，是中华民族的传统美德。诚\par

也是一项民法原则。\par

7.1\par


\section*{PDF第33页}
\addcontentsline{toc}{section}{PDF第33页}

【结构整理：违法行为的分类及异同点表】\par

表头：类别｜违反的法律｜对社会危害程度｜承担的责任｜举例\par

民事违法行为｜民事法律规范（如《民法典》）｜较轻｜民事责任（如停止侵害、返还财产、恢复原状、赔偿损失等）｜欠债不还、违反合同、侵犯肖像权、著作权、隐私权等。\par

行政违法行为｜行政法律规范（如《行政处罚法》《治安管理处罚法》《义务教育法》等）｜较轻｜行政制裁，包括行政处分和行政处罚（警告、罚款、没收违法所得、行政拘留等）｜扰乱社会治安、谎报险情、破坏铁路封闭网、殴打他人等。\par

说明：民事违法行为和行政违法行为都属于一般违法行为。\par

\daofasym{②}诚信是一个人安身立命之本，是我们融入社会的“通行证”。（对个人)\par

\daofasym{③}诚信是企业的无形资产。一个企业只有坚持诚信经营、诚信办事，才能塑造良好的形象和信誉，赢得客户；才能带来\par

持久的效益，长盛不衰。(对企业)\par

\daofasym{④}诚信促进社会文明、国家兴旺。国无信则衰，社会成员之间以诚相待、以信为本，能够增进社会互信，减少社会矛盾，\par

净化社会风气，促进社会和谐：能够降低社会交往和市场交易成本，积累社会资本：能够提高国家的形象和声誉，增强\par

国家的文化软实力。（对国家和社会）\par

2. 如何践行诚信\daofasym{★}\daofasym{★}\daofasym{★}\par

\daofasym{①}树立诚信意识。我们要真诚待人，信守承诺，说老实话，办老实事，做老实人，勇于承认过错，主动承担责任，争取\par

他人的谅解。\par

\daofasym{②}运用诚信智慧。当尊重他人隐私与对人诚实发生冲突时，我们应遵循伦理原则和法律要求，权衡利弊，既恪守诚实的\par

要求，又尊重他人隐私。\par

\daofasym{③}珍惜个人的诚信记录。大力弘扬诚信文化，以诚实守信为荣，以见利忘义为耻，营造社会诚信环境。\par

\daofasym{④}我们要弘扬诚信文化，提高全社会信用水平，营造社会诚信环境，努力促进社会发展和文明进步。\par

3、（辨析题）小案认为，生活中难免会说一些“善意的谎言”，诚实和保护隐私可以并存。请评析观点？\par

答案：正确。社会生活是复杂的，我们有时会面临艰难抉择。运用诚信智慧。当尊重他人隐私与对人诚实发生冲突时，\par

我们应遵循伦理原则和法律要求，权衡利弊，既恪守诚实的要求，又尊重他人隐私。\par

4、请你为构建诚信社会提出合理化建议？（分主体一一国家、社会、个人、企业等）\par

(1)国家：\daofasym{①}建立健全相关法律法规，加大执法力度，严惩不讲诚信的行为：\daofasym{②}加强法治建设，推进个人诚信体系和社会\par

信用体系建设，完善个人守信激励和失信惩戒机制；\daofasym{③}加强道德建设和法治教育提高公民的诚信意识、增强公民的法\par

治观念。\par

(2)社会：大力弘扬诚信文化，共同营造“以诚实守信为荣、以见利忘义为耻”的良好社会风尚，提高全社会信用水平，\par

营造社会诚信环境，努力促进社会发展和文明进步。 UR多性\par

到而违法1为 \daofasym{②}迅以·法伴污规特记\par

\subsection*{第五课做守法的公民同处与}
\addcontentsline{toc}{subsection}{第五课做守法的公民同处与}

5.1法不可违\par

1.法律的内涵：法律是全体社会成员都要共同遵守的行为规范。法律是评价人们的行为是否合法有效的准绳。法律是最刚\par

性的社会规则，不违法是人们行为的底线。\par

2.什么是违法行为？违法行为是指出于过错违反法律、法规的规定，危害社会的行为。\par

3.违法行为的分类及异同点\daofasym{★}\daofasym{★}\daofasym{★}\daofasym{★}\daofasym{★}\par

根据违反法 对社会危害\par

违反的法律 承担的责任 举例\par

律的类别 程度\par

民事责任（如停止欠债不还，违反合同，侵犯肖\par

民事违法行民事法律规范 侵害，返还财产，像权、著作权、隐私权等。（侵\par

较轻\par

（如《民法典》） 恢复原状等)赔犯他人的民事权利，包括人身\par

权和财产权)\par

行政制裁：行政处记 陈联见\par

分和 记大过\par

一般违法行\par

为 行政法律规范 行政处罚（罚\par

行政违法行（如《行政处罚法》 扰乱社会治安（说报险情、\par

为 《治安管理处罚法》 较轻款，没收违法所破坏铁路封闭网、殴打他人）\par

《义务教育法》《物\par

权法》） 得，行政拘\par

留）。\par


\section*{PDF第34页}
\addcontentsline{toc}{section}{PDF第34页}

【结构整理：违法行为的分类表续】\par

刑事违法行为（犯罪）｜刑事法律规范（如《刑法》）｜严重｜刑罚处罚｜故意杀人、抢劫、贩毒、醉驾等。\par

【结构整理：一般违法行为与刑事违法行为的异同点】\par

相同点：\daofasym{①}都有社会危害性；\daofasym{②}都是违法行为；\daofasym{③}都要承担法律责任。\par

不同点：一般违法行为包括民事违法行为和行政违法行为，它们对社会的危害相对轻微；刑事违法行为是违法行为中最严重的一种，就是我们常说的犯罪，社会危害性大，情节严重。\par

严重违法行刑事违法行刑事法律规范 故意杀人，抢劫，贩毒，醉驾\par

严重 刑罚处罚\par

为 (犯罪) 为 (如《刑法》)\par

一般违法行为与刑事违法行为的异同点：\daofasym{★}\daofasym{★}\daofasym{★}\par

（1）相同点：\daofasym{①}都有社会危害性；\daofasym{②}都是违法行为；\daofasym{③}都要承担法律责任。\par

（2）不同点：一般违法行为包括民事违法行为和行政违法行为，它们对社会的危害相对轻微；而刑事违法行为是违\par

为中最严重的一种，就是我们常说的犯罪，社会危害性大，情节严重。\par

4.身边的违法行为有哪些?\daofasym{★}\daofasym{★}\daofasym{★}\par

(1)侵犯他人民事权利或者没有依法履行合同的义务，都是典型的民事违法行为。\par

(2)最常见的行政违反行为是违反治安管理的行为。谎报险情、破坏铁路封闭网、段打他人等行为都是违反治安管理\par

法的行政违法行为。\par

5.遵章守法是社会和谐的保证。生活在法治社会里，我们应该怎么做？\daofasym{★}\daofasym{★}\par

\daofasym{①}我们要认识一般违法行为的危害，自觉依法规范自己的行为。\par

\daofasym{②}在社会生活中，要分清是非，增强守法观念，严格遵守治安管理的法律规定。\par

\daofasym{③}在社会交往中，要依法从事民事活动，积极防范民事侵权行为和合同违约行为，既要维护自己的权益，又要尊重他人\par

的权益，促进社会健康和谐发展。\par

5.2 预防犯罪\par

1. 刑法的作用\par

刑法是惩治犯罪，保护国家和人民利益的有力武器。它明确规定了什么行为是犯罪，以及对犯罪应当判处什么样的刑\par

2.犯罪（1）含义：犯罪是具有严重社会危害性、触犯了刑法、应当受到刑罚处罚的行为。\par

（2） 基本特征\daofasym{★}\daofasym{★}\daofasym{★}\daofasym{★}\daofasym{★}\par

（严重社会危害性（最本质特征)刑事违法性（法律标志应受刑罚处罚性（必然法律后果）\par

（3）后果：法网恢恢，疏而不漏，任何人的犯罪行为都应当受到法律的制裁，犯罪的法律后果是刑罚。\par

3.刑罚的种类\par

(1）主刑:（管制（3 月-2 年）6拘役（1月-6 月有期徒刑（6月-15年）无期徒刑、死刑;\par

（2）附加刑罚金、剥夺政治权利、没收财产、驱逐出境（适用于外国人）\par

4.怎样加强罪的自我防范？（个人层面）\daofasym{★}\daofasym{★}\daofasym{★}\daofasym{★}\daofasym{★}\par

\daofasym{①}要珍惜美好生活，认清犯罪危害，远离犯罪。\par

\daofasym{②}杜绝不良行为，从小事做起，避免沾染不良习气，自觉遵纪守法，防患于未然。\par

\daofasym{③}增强法治观念，依法自律，做一个自觉守法的人。礼身他人\par

5.如何预防未成年人犯罪？（开放题）\daofasym{★}\daofasym{★}\daofasym{★}\daofasym{★}\daofasym{★}\par

个大\daofasym{①}把4面第4点知识点\daofasym{①}\daofasym{②}\daofasym{③}合为一点；\par

家庭：\daofasym{②}家长应尽到监护职责，加强对未成年人的教育和管理；\par

学校：\daofasym{③}应加强对未成年人的道德、法制教育，对他们的不良行为及时发现，及时教育引导：\par

社会：\daofasym{④}加大法制宣传，积极营造良好的社会文化环境。\par

国家有关部门：\daofasym{③}要完善相关法律法规，对于侵犯未成年人的违法行为，有关部门要加强监管和严格执法、严惩。\par

6.不良行为和犯罪之间没有不可逾越的鸿沟，有了不良行为如果不加以改正，任其发展，就有可能滑进犯罪的深渊。\par

5.3善用法律防微杜渐\par

\textbf{1.学会用法的必要性是什么？}\par


\section*{PDF第35页}
\addcontentsline{toc}{section}{PDF第35页}

答：在社会生活中，我们要学会用法律与人打交道。在遇到法律问题或者权益受到侵害时，要及时寻求法律救助，依靠法\par

律维护自己的合法权益。\par

2. 导求法律救助的途径有哪些?\daofasym{★}\daofasym{★}\daofasym{★}\daofasym{★}\daofasym{★}(2 种)\par

\daofasym{①}我们可以通过法律服务机构来维护合法权益，如法律服务所、律师事务所、公证处、法律援助中心等。\par

\daofasym{②}受到非法侵害，可以寻求国家的法律救济。我们可以依法到公安机关、人民法院或人民检察院中的任何一个机关控\par

告、举报，必要时可以真接向人民法院起诉。\par

3. 依法维权方式**4用年\par

（1)非诉讼方式：越扇、(调解、仲裁等。 (2) 诉讼方式\par

4.诉讼（1）含义：诉讼是人民法院在诉讼当事人参与下，依照法定程序解决纠纷和冲突的活动。\par

（2）地位：诉讼是处理纠纷和应对侵害最正规、最权威、最有效的手段，是维护合法权益的最后保障。（选择。另外\par

\textbf{一定要记得是向法院提起诉讼？}\par

（3）分类：民事诉讼、行政诉讼、刑事诉讼\par

（4）适用条件：如果受到非法侵害，通过非诉讼手段（借助第三方）不能解决问题，就要通过使用诉讼手段，通过打官司\par

讨回公道。\par

5.为什么要同违法犯罪作斗争？\daofasym{★}\daofasym{★}\daofasym{★}\daofasym{★}\daofasym{★}\par

同违法犯罪作，是包括我们青少年在内的全体公民义不容辞的责任。当国家利益、公共利益、本人或他人的权益受\par

到不法侵害时，我们要敢于并善于依法维护正当权益。\par

6. 青少年见义智为的原因\daofasym{★}\daofasym{★}\daofasym{★}\daofasym{★}\daofasym{★}\par

未成年体力不具优势心智尚未成熟，如果鲁莽行事，自己极易受到伤害，也不利于制止违法犯罪。\par

7.青少年怎么圆选法犯罪作斗争？\daofasym{★}\daofasym{★}\daofasym{★}\daofasym{★}\daofasym{★}\par

（1）同违法犯罪作斗争，是包括我们青少年在内的全体公民义不容辞的责任。当国家利益、公共利益、本人或他人的权\par

益受到不法侵害时，我们要敢于并善于依法维护正当权益。\par

（2）未成年体力不具优势，心智尚未成敦，如果鲁茬行事，自己极易受到伤害，也不利于制止违法犯罪。\par

我们要善于斗争，在保全自己、减少伤害的前提下，巧妙地借助他人或社会的力量，采取机智灵活的方式，同违法犯罪\par

作斗争。\par

（3）回违法犯罪作斗争的常见方法：\daofasym{①}及时拨打110报整电话或争取其他成人的帮助\daofasym{②}虚张声势，与违法犯罪分子巧妙\par

周旋。\daofasym{③}记住违法犯罪分子的体貌特征。\daofasym{④}了解违法犯罪分子的去向。\daofasym{③}保护好作案现场。\par

(4)我们要积极弘扬社会主义法治精神，形成守法光荣、违法可耻的观念，做到自觉守法、遇事找法、解决问题靠法，努\par

力成为一名社会主文法的忠实崇尚者、自觉遵守者和坚定捍卫者。\par

8、\par

\subsection*{第三单元勇担社会责任}
\addcontentsline{toc}{subsection}{第三单元勇担社会责任}

\subsection*{第六课责任与角色同在}
\addcontentsline{toc}{subsection}{第六课责任与角色同在}

6. 1 我对谁负责谁对我负责\par

1. 责任\par

（1）含义：责任是一个人分内应该做的事\par

（2）来源：来自对他人的承诺职业要求道德规范法律规定等。\par

（3）贵任与角色关系\par

\daofasym{①}在社会生活的舞台上，每个人都扮演着不同的角色。随着时代发展和所处环境的变化，我们会不断变换自己的角色，\par

调节角色行为。\par

\daofasym{②}每一种角色都意味着承担相应的责任。只有人人认识到自己扮演的角色、承担应尽的责任，才能构建各尽其能、各得\par

其所而又和谐相处的社会。\par

2.如何理解“我对谁负责，谁对我负责”？\daofasym{★}\daofasym{★}\daofasym{★}\daofasym{★}\daofasym{★}\par

\daofasym{①}作为社会的一员，我们每个人首先要对自己负责。（我对我负责）\par

\daofasym{②}很多人在为我们的成长和幸福生活承担着责任。（他人对我负责）\par

\daofasym{③}我们应该学会感恩，主动关心、帮助、和服务他人，在辛苦付出中体验对他人负责的快乐和幸福。（我对他人负责）\par

3. 承担责任的重要性2\daofasym{★}\daofasym{★}\daofasym{★}\daofasym{★}\daofasym{★}\par


\section*{PDF第36页}
\addcontentsline{toc}{section}{PDF第36页}

\daofasym{①}只有对自己负责的人，才能使自己的潜能得到充分挖掘和发挥，才有资格、有能力、有信心承担起时代和国家所赔\par

的使命。\par

\daofasym{②}很多人在为我们的成长和幸福生活承担着责任。如果没有他们的付出，没有他们履行各自承担的责任，我们的生活\par

设想。\par

\daofasym{③}只有人人具有责任心，自觉履行应尽的责任，我们才能共享更加幸福美好的生活。\par

6. 2 做负贵任的人\par

1.承担责任会付出哪些代价？又有哪些回报？\daofasym{★}\daofasym{★}\daofasym{★}\daofasym{★}\daofasym{★}\par

(1）代你。时间精力金钱,)做得不好受到责备，甚至处罚。\par

（2）回报：承担责在在在伴随着获得回报的权利。既包括物质方面，又包括精神方面。更重要的是精神方面的回报，\par

良好的自我感觉、获得新的知识或技能、赢得他人的尊重和赞许等。\par

\textbf{2.如何选择自己承担的贡任？}\par

要有勇气凭借自己的经验和智慧，对承担责任的代价和回报做出正确评估，做出合理的选择。一旦做出选择，就应该\par

无反顾地担当起应负的责任。\par

3. 如何对待生活中不可推卸的责任？\daofasym{★}\daofasym{★}\daofasym{★}\daofasym{★}\daofasym{★}\par

\daofasym{①}虽然有些应该做的事情不是我们自愿选择的，但是我们仍然应该自觉承担相应的责任。\par

\daofasym{②}我们往往无法改变自己在生活和工作中的位置，但我们可以改变自己对待应该做的事情的态度。\par

\daofasym{③}只要我们把它们当做一种不可推卸的责任担在肩头，不抱怨、不怠，全身心地投入，同样能够把事情做得很出色\par

4.履行责任\par

（1）表现：在偏远山村支教者、不留姓名救人者、保护珍贵野生动物者等。\par

（2）意义：敢于承担责任，敢于担当，我们的生活才更加安全，更加温暖，更加充满阳光希望。\par

5.怎样做负责任的人？\daofasym{★}\daofasym{★}\daofasym{★}\daofasym{★}\daofasym{★}\par

@还言代价与回报2\par

\daofasym{②}做到我承担我无悔。有些事情不是我们自愿选择的但是我们仍然应该自觉承担相应的责在。\par

\daofasym{③}努力提开自身质，增强履行责任的能力，勇于承担责任，在奉献社会过程实现自己的人生价值。\par

\subsection*{第七课 积极奉献社会}
\addcontentsline{toc}{subsection}{第七课 积极奉献社会}

7.1 关爱他人\par

1. 关爱\par

（1）含义：关爱就是关心爱护。\par

（2）表现：舒个人都被他人关爱着，也都在关爱着他人。关爱无时不在，无处不在。\par

2.关爱他人的意义\daofasym{★}\daofasym{★}\daofasym{★}\daofasym{★}\daofasym{★}\par

\daofasym{①}关爱传递着美好情感，给大带来温暖和希望，是维系友好关系的桥梁（对他人）\par

\daofasym{③}关受是社会和谐稳定的润滑剂和正能量，关爱使人民在交往过程中互谅互让，相互尊敬，与人为善，增进信任，\par

主形成良妊的大际氛围。促进社会文明进步。（对社会\par

\daofasym{③}关爱他人，收获幸福。关爱他人也是关爱和善待自己。（对自己）\par

3.怎样关爱他人？\daofasym{★}\daofasym{★}\daofasym{★}\daofasym{☆}\daofasym{★}\par

\daofasym{①}关爱他人，要心怀善意。当他人遇到困难的时候，我们应该在道义上给予支持，物质上给予帮助，精神上给予关\par

\daofasym{②}关爱他人，要尽已所能。关爱不分大小，贵在有爱心。一个人的能力有大小，但只要尽已所能为他人排忧解难、\par

社会，就是一个友善和值得称赞的人。\par

\daofasym{③}关受他人，要讲究策略。A、考虑他人的内心感受，不伤害他人的自尊心。 善于作出明智的判断，增强防范意\par

自我保护意识，任保护自己不受伤害的前携下采取果敢和理智的行动。\par

7.2服务社会\par

1.服务赶会对个人成长的意义\daofasym{★}\daofasym{★}\daofasym{★}\daofasym{★}\daofasym{★}\par

\daofasym{①}服务在会体现人生价值〉A.一个人的价值应该看他贡献什么，而不应当看他得到什么。人生的真正价值在于对社\par

责任和贡献。B.人人都有责任回报社会C只有积极为社会作贡献，才能得到人们的尊重和认可，实现我们自身的价\par

\daofasym{②}服务社会能够便进我们全面发展。在服务社会的过程中，我们的视野不断拓展，知识不断丰高，观察、分析、解\par

题的能力以及人标交往能力不断提升，道德境界不断提高。\par

\textbf{2.怎样服务社会奉献社会？}\par

\daofasym{①}需要我忙否年担当责任。需要我们积极参与社会公益活动。从实际出发，讲求实际效果。\daofasym{③}需要我们热爱劳动，\par

岗敬业。我们要努力学习，增强劳动观念，焙养敬业精神，精益求精，全力以越\par

3.社会公益活动\par

页 5 36\par


\section*{PDF第37页}
\addcontentsline{toc}{section}{PDF第37页}

（1）具体形式：环境保护、社区服务、在社区、公园、车站等公共场所纠正不文明行为、到科技馆、博物馆做志愿者。\par

(2）要求：从实际出发，讲求实际效果。\par

4.中国共产主义青年团\par

（1）性质：中国共产主义青年团是中国共产党领导的先进青年的群团组织，是广大青年在实践中学习中国特色社会\par

主义和共产主义的学校，是中国共产党的助手和后备军。\par

（2）作用：中国共产主义青年团带领青年在经济社会发展中发挥生力军和突击队作用。\par

（3）任务：在新时代，共青团要组织青年参加改革开放和社会主义现代化建设的实践，为发展社会生产力、提高人民生\par

活水平，为实现“两个一百年”奋斗目标建功立业。\par

\subsection*{第八课 国家利益至上}
\addcontentsline{toc}{subsection}{第八课 国家利益至上}

8.1国家好 大家才会好\par

1.国家利益\par

（1）含义：是一个主权国家在国际社会中生存需求和发展需求的总和，包括人口、领土、主权和政权等\par

（2）重要性：关系到民族生存、国家兴亡。\par

（3）范围：涉及政治、经济、文化、社会、军事等领域，包括安全利益、政治利益、经济利益、文化利\par

益等。\par

2.国家的核心利益的范围（选择题核心利益有6个）\par

国家主权国家安到领土完整国家统）、宪法确定的国家政治制度和社会大局稳定、经济社会可持续\par

的基本保障。\par

3.如何理解“国家利益是人民利益的集中表现”？（国家利益与人民利益的关系）\daofasym{★}\daofasym{★}\daofasym{★}\daofasym{★}\daofasym{★}\par

\daofasym{①}在我们国家，国家利益反映广大人民的共同需求，是人民利益的集中表现，两者相辅相成。\par

\daofasym{②}人民利益只有上升、集中到国家利益，运用国家的工具，才能得到真正的维护。\par

\daofasym{③}国家利益只有反映人民利益，依靠人民艰苦奋斗，才能得到真正实现。\par

在当代中国，国家利益与人民利益是高度统一的。将国家和人民视为一个命运共同体，将国家利益和人\par

民利益紧密联系在一起。\par

4.实现中华民族伟大复兴最鲜明的特点：就是将国家和人民视为一个命运共同体，将国家利益和人民利益\par

紧密联系在一起。\par

最长达较根于刊二一)刊道\par

8.2坚持国家利益至上\par

1.国家利益与集体利益、个人利益的关系\daofasym{★}\daofasym{★}\daofasym{★}\daofasym{★}\daofasym{★}=如何正确认识国家利益和个人利益的关系\par

\daofasym{①}国家利益与集体利益、个人利益既有区别，又有联系。\par

\daofasym{②}国家利益是整体利益，集体利益、个人利益是局部利益。整体利益照顾到了大多数人民的长远需求，符\par

合大多数群体的利益。（区别)\par

\daofasym{③}国家利益与个人利益有时会有冲突。但是从根本上来说，国家利益与集体利益、个人利益是一致的。（联\par

系)\par

2.如何坚持国家利益垒上？\daofasym{★}\daofasym{★}\daofasym{★}\daofasym{★}\daofasym{★}\par

（1）树立维护国家利益意识（思想上）\par

\daofasym{①}要心怀爱国之情，牢固树立国家利益至上观念，以热爱祖国为荣，以危害祖国为耻。\par

\daofasym{②}要树立和增强危机意识和防范意识。\par

\daofasym{③}要增强维护国家利益的贵任感和使命感。珍惜发展机遇努力学习提高素质，用理性、务实、文明的\par

心态，合法有序地表达爱国情感，维护国家利益。\par

（2）捍卫国家利益（行动上）=如何正确处理国家利益和个人利益的关系\par

\daofasym{①}无论何时何地，我们都应当着眼长远，顾全大局，以国家利益为重，把国家利益放在第一位\par

\daofasym{②}为了国家利益，有时不仅需要放弃个人利益，甚至要献出白已的生命。\par

\daofasym{③}我们要始终把国家利益放在第一位，捍卫国家尊严，坚决同一切损害国家利益的行为作斗争。在日常生\par

活中，我们要自觉逆守道德与法律，积极维护国家团结稳定的局面。\par


\section*{PDF第38页}
\addcontentsline{toc}{section}{PDF第38页}

\subsection*{第九课 树立总体国家安全观}
\addcontentsline{toc}{subsection}{第九课 树立总体国家安全观}

物 9.1认识总体国家安全观\par

1.国家安全：是指国家政权、主权、统一和领土完整，经济社会可持续发展和国家其他重大利益相对处于没有危险\par

受国内外威胁的状态，以及保障可持续安全状态的能力。\par

※2.国家安全对国家的意义\par

\daofasym{①}国家安全是实现国家利益贵想本的保障。P96\par

\daofasym{②}国家安全是国家生存串发展的重要保障。国家安全有保障，经济社会才能不断发展，祖国才能更加繁荣富强。\par

\daofasym{③}国家安全是人民幸福安康的前提。\par

3.为什么要树立总体国家安全观？（我国目前的安全形势是怎么样的）\par

今天，我国国家按全的内涵和外延比历史上任何时期都更丰富，时空领域比历史上任何时期都更宽广，内外因素比！\par

上任何时候都更复杂。\par

\textbf{4.如何构建国家安全体系？}\par

\daofasym{①}以大民安垒为宗旨，以政治安全为根本，以经济安全为基础，以军事、文化、\par

会安全为保障，以促进国际安全为依托，走出一条中国特色的安全道路。（非常\par

常重要，会考选择）\par

\daofasym{②}面对国家安全形势的新变化，需要“十个重视”，如，既重视外部安全，又重视内部安全；既重视发展问题，又\par

安全问题；既重视自身安全，又重视共同安全等。\par

\daofasym{③}我国要构建集政治安全、国土安全、军事安全、资源安全、核安全等于一体的国家安全体系。\par

5、为什么要维护国家安全？（答案第二题+第三题）\par

9.2维护国家安全\par

\textbf{※1、我国武装力量属于谁？任务是什么？}\par

\daofasym{①}性质：中华人民共和国的武装力量属于人民。\par

\daofasym{②}任务：巩固国防，抵抗侵略，保卫祖国，保卫人民的和平劳动，参加建设事业，努力为人民服务。\par

\textbf{※2、如何全面推进国防和军队现代化？}\par

\daofasym{①}必须全面贯彻新时代党的强军思想，贯彻新形势下军事战略方针。\par

\daofasym{②}坚持党对人民军队的绝对领导。（根本原则，根本制度）\par

\daofasym{③}力争到二O三五年基本实现国防和军队现代化，到本世纪中叶把人民军队全面建成世界一流军队。\par

\textbf{※3、为什么维护国家安全需要我们每个人作出贡献？}\par

\daofasym{①}维护国家安全，人人可为。维护国家安全是全国各族人民根本利益所在，是我们的共同贵任。\par

\daofasym{②}我国究法和法律明确规定了公民和组织应当履行维护国家安全的义务。\par

\textbf{※4、我们应该怎样维护国家安全？}\par

\daofasym{①}我们要增强国家安全意识，树立国家安全利益高于一切的观念，自觉维护国家安全，推动全社会形成维护国家安\par

强大合力。（思想上）\par

\daofasym{②}我们可以通过各种方式为维护国家安全贡献智慧和力量。（行动上）\par

\daofasym{③}我们要认真学习有关国家安全保密工作的法律法规、规章制度，（增强维护国家安全的法治意识。（法律）\par

防范能力\par

5、我国为何要加强国防和军队现代化建设（答案：1\daofasym{②}+2\daofasym{①})\par

\subsection*{第十课建设美好祖国}
\addcontentsline{toc}{subsection}{第十课建设美好祖国}

10. 1关心国家发展\par

\textbf{1.我国取得了哪些全世界瞩目的成就？}\par

\daofasym{①}我国经济增长成就显著，已经成为世界第二大经济体。\par


\section*{PDF第39页}
\addcontentsline{toc}{section}{PDF第39页}

【漏识补录：第5题劳动素养第（1）点】\par

（1）树立正确劳动观念，尊重劳动，牢固树立劳动最光荣、劳动最崇高、劳动最伟大、劳动最美丽的思想观念。\par

有关间影丝的产中国特色安全通构国体乐\par

区年化\par

\daofasym{②}民主法治建设深入推进，人民的权利和自由得到切实保障，\par

\daofasym{③}文化建设焕发出蓬勃生机 发史\par

\daofasym{④}科技创新成就斐然。 益\par

\daofasym{③}惠民利民政策陆续出台，社会保障水平不断提高。\par

\daofasym{③}国防和军队改革取得重大突破\par

\daofasym{③}综合国力显著提升。\par

2. 我国发展中的面临的问题\daofasym{★}\daofasym{★}\daofasym{★}\daofasym{★}\daofasym{★}\par

\daofasym{①}人均GDP（国内生产总值）的世界排名还不高：\daofasym{②}发展不平衡、不协调、不可持续问题仍然突出：\daofasym{③}居民收入差距依\par

然较大，部分群众生活面临一些困难等。\par

3. 我国对发展中问题的解决措施\daofasym{★}\daofasym{★}\daofasym{★}\daofasym{★}\daofasym{★}\par

国家正在采取各种积极措施，稳增长、促改革、调结构、惠民生、防风险，着力解决各种发展中的问题，不断取得积极\par

成效，这让我们对国家的未来发展更加充满信心和期待。\par

10.2 天下兴亡匹夫有责\par

1. 劳动的作用\daofasym{★}我\daofasym{★}\daofasym{★}\daofasym{★} 道影是次权也是项务孙数足质收育至要\par

\daofasym{①}劳动是财富的源泉，也是幸福的源泉。\par

\daofasym{②}人世间的美好梦想，都是通过劳动实现的，生命里的一切辉煌，一都是通过劳动铸就的。 组成)奇作效为\par

\daofasym{③}我们国家所取得的每一项成就，都是广大人民用辛勤劳动、诚实劳动、创造性劳动换来的，中国人民用实干精神创造\par

了今天的辉煌。 确生社会卖任感使其真成长内续负之性人:\par

\daofasym{①}中小学生参加劳动有利子中小学生 幼观）美成良好的劳动习惯有利手增强中小学生的身体素质促进\par

中小学先健康成长：有利干中小学生旭求高雅的生活情趣文陶治自身的情操\par

\daofasym{①}每个人所处的岗位不同，从事不同的劳动，但都在为国家和社会发展作出贡献。\par

\daofasym{②}正是无数劳动者兢兢业业、艰苦奋斗、无私奉献，成就了我们今天的美好生活\par

\daofasym{③}无论是脑力劳动者还是体力劳动者，都是国家的建设者，都值得我们尊敬和学习。\par

3如何使中国梦变成现实？\daofasym{★}\daofasym{★}\daofasym{★}\daofasym{★}\daofasym{★}实干创造未来。 书成就今天 至同有我\par

\daofasym{①}把中国梦变成现实，创造未来的美好生活，需要一代代水理头苦≠和按力奋头，盂要每下人在多者岗位上付出更多的\par

辛勤和汗水。\par

\daofasym{②}青少年要努力学习(积极探，勇做走在时代前列的学习者、劳动者、奉献者以拱着的信念、优良的品德、丰富的\par

知识、过硬的本领，担负起历史重任，让青春绽放出绚丽的光彩。\par

4.为了国家发展，青少年怎么做？ )同炎静强购同强一化自理供,可担当\par

\daofasym{①}努力学习科学文化知识，树立远大理想。\daofasym{②}积极参加社会实践活动，养成亲社会行为，承担社会责任。有本领同易新有死这泛就向节望,月年接力张同\par

\daofasym{③}培养创新精神和创新意识，发扬实干精神。 有我·感们人以婷府起实现“\par

5、中学生应如何全面提高自身劳动素养？ 有构的美任·成为果\par

(2)具有必备的劳动能力，掌握基本的劳动知识和技能。\par

(3)培育积极的劳动精神。继承中华民族勤俭节约、敬业奉献的优良传统，弘扬开拓创新、砥砺奋进的时代精神。\par

(4)养成良好的劳动习惯和品质。形成诚实守信、吃苦耐劳的品质。珍惜劳动成果，养成良好的消费习惯，杜绝浪费等。\par

\subsection*{八下}
\addcontentsline{toc}{subsection}{八下}

2022八下道法提纲\par

\subsection*{第一课维护宪法权威}
\addcontentsline{toc}{subsection}{第一课维护宪法权威}


\section*{PDF第40页}
\addcontentsline{toc}{section}{PDF第40页}

1.1.党的主张和人民意志的统一\par

\textbf{※1.为什么要坚持中国共产党的领导？}\par

\daofasym{①}中国新民主主义革命的胜利和社会主义事业成就，是中国共产党领导中国各族人民，战胜许多艰难险阻\par

而取得的。\par

\daofasym{②}宪法规定：中华人民共和国是工人阶级领导的、以工农联盟为基础的人民民主专政的社会主义国家。\par

中国共产党领导是中国特色社会主义最本质的特征。\par

\daofasym{③}性质决定：中国共产党是中国工人阶级的先锋队，同时是中国人民和中华民族的先锋队。全心全意为\par

人民服务是党的根本宗旨，党的最高理想和最终目标是实现共产主义。\par

\daofasym{④}最大优势：中国共产党领导是中国特色社会主义最本质的特征，是中国特色社会主义制度的最大优势\par

党是最高政治领导力量。\par

******有关党的知识\par

\daofasym{①}坚持以人民为中心的发展思想；\par

\daofasym{②}党的执政理念之一是共享发展成果；\par

\daofasym{③}中国共产党是中国特色社会主义建设事业的领导核心；代表了最广大人民的根本利益\par

\daofasym{④}党的宗旨是全心全意为人民服务；\par

\daofasym{③}党的初心和使命是为中国人民谋幸福，为中华民族谋复兴；\par

\daofasym{③}党的奋斗目标是人民对美好生活的向往。\par

中国共产党是中国工人阶级的先锋队，同时是中国人民和中华民族的先锋队\par

\daofasym{③}中国共产党领导是中国特色社会主义最本质的特征，是中国特色社会主义制度的最大\par

势，是党和国家的根本所在，命脉所在，是全国各民族人民的利益所系，命运所系。党\par

最高政治领导力量\par

\textbf{※2.中国共产党与宪法和法律的关系？}\par

\daofasym{①}宪法是党的主张和人民意志的统一。\par

\daofasym{②}宪法确立党的领导地位\par

\daofasym{③}中国共产党领导人民制定宪法法律，领导人民实施宪法法律，做到党领导立法、保证执法、支持司法\par

带头守法。\par

\daofasym{④}中国共产党离履行好执政兴国的重大职责，必须在宪法和法律范围内活动，依据宪法法律治国理政。\par

※3.我国宪法的基本原则：我国是人民民主专政的社会主义国家，国家的一切权力\par

属于人民（非常非常重要）只要考试考到为什么重视人民、人权、做一些对人民\par

利的事情等问题都可以用。）\par

国家性质：中华人民共和国是工人阶级领导的、以工农联盟为基础的人民民主专政的社会主义国家。\par

※4.宪法如何保证国家权力属于人民？（归根到底是为了保证人民当家作主）\par


\section*{PDF第41页}
\addcontentsline{toc}{section}{PDF第41页}

\daofasym{①}宪法确认我国的国家性质，明确人民当家作主的地位。（性质上)\par

\daofasym{②}宪法规定的社会主义经济制度奠定了国家权力属于人民的经济基础。宪法规定，我国经济制度的基础是\par

生产资料的社会主义公有制。(经济上)\par

\daofasym{③}宪法规定的社会主义政治制度明确了人民行使国家权力的基本途径和行使。（政治上）\par

\daofasym{④}宪法规定广泛的公民基本权利，并规定实现公民基本权利和保障措施。（权利上）\par

\daofasym{③}宪法还规定国家武装力量属于人民，它的任务是巩固国防，抵抗侵略，保卫祖国，保卫人民的和平劳动,\par

参加国家建设事业，努力为人民服务。(武装归属上)\par

5.人权的实质内容和目标是：是人的自由、平等地生存和发展\par

人民幸福生活是最大的人权（教材新加的内容）\par

我国将“国家尊重和保障人权”写入宪法。\par

人权的主体：既包括我国公民，也包括外国人等。不仅保护个人，也保护群体。\par

宪法保护人权的内容：既包括平等权和人身权利、政治权利，也包括财产权、劳动权、受教育权等经济、\par

社会、文化方面的权利。\par

我国人权的特点：广泛性、公平性、真实性\par

\textbf{※6.国家应如何做到尊重和保障人权？}\par

\daofasym{①}尊重和保障人权是现代法治国家立法活动的基本要求。立法方面）\par

\daofasym{②}行政机关在执法过程中应当树立尊重和保障人权的意识，做到严格规范公正文明执法，坚持依宪施政、\par

依法行政、简政放权。（执法方面）\par

\daofasym{③}审判机关、检察机关、监察机关要依照宪法和法律的规定独立行使审判权、检察权、监察权，保护公民\par

的各项合法权益。（司法和监察方面）\par

\daofasym{④}国家加强法治宣传教育，弘扬社会主义法治精神，建设社会主义法治文化，增强全民法治观念，形成全\par

民守法的氛围和习惯，努力将人权理想变成现实。（普法方面）\par

※7.列举国家尊重和保障人权的措施（补充）\par

\daofasym{①}发展义务教育\daofasym{②}为困难地区困难学生提供营养午餐\daofasym{③}开办农家书屋\daofasym{④}开展再就业培训\daofasym{③}精准扶贫：\daofasym{③}强\par

化食品安全监管治理环境突出问题。\par

\textbf{8.如何坚持党的领导？}\par

\daofasym{①}党政军民学，东西南北中，党是领导一切的。\par

\daofasym{②}必须坚决维护党中央权威，健全总揽全局、协调各方的党的领导制度体系，把党的领导落实到国家治理\par

各领域各方面各环节。\par

\daofasym{③}中国共产党领导人民制定宪法法律，领导人民实施宪法法律，做到党领导立法、保证执法、支持司法、\par

带头守法。\par

\daofasym{④}中国共产党要履行好执政兴国的重大职责，必须在宪法和法律范围内活动，依据宪法法律治国理政。\par

9、\daofasym{★}为什么要尊重和保障人权？（P7、P9）\par

\daofasym{①}尊重和保障人权是我国的宪法原则。\par

\daofasym{②}尊重和保障人权是立法活动的基本要求。\par

\daofasym{③}我国是人民民主专政的社会主义国家，人民是国家的主人\par

几个重要句子：\par

\daofasym{①}宪法是党的主张和人民意志的统一\par

\daofasym{②}我国是人民民主专政的社会主义国家，人民是国家的主人\par

\daofasym{③}我国宪法的基本原则：我国是人民民主专政的社会主义国家，国家的一切权力属于人民。\par

\daofasym{④}宪法还有一个重要原则：尊重和保障人权\par

（注意）P11拓展空间以下句子也可以划一下线，因为简答题材料如果出现“脱贫、精准扶贫”之\par

类的材料，经常会问体现我国哪一宪法原则？要答“尊重和保障人权”。\par

4/\par


\section*{PDF第42页}
\addcontentsline{toc}{section}{PDF第42页}

划线句子贫穷是实现不权的最大障碍（选择）。中国的减贫行动是中国人权事业进步的最显著\par

志。中国的减贫行动有力促进了贫困大口发展权的实现，为全面建成小康社会打下了坚实基础。\par

注意：宪法的基本原则和重要原姒不要混在一起。\par

\daofasym{③}宪法的核心价值追求一一规范国家权力运行以保障公民权利（选题题经常出现）\par

@我国已经消除了绝对贫困，记得！！！\par

\daofasym{③}国家机关的组织和活动原则：民主集中制 h随退 粥 么.正义\par

**宪法原则：\par

\daofasym{①}（基本原则，只有一个）我国是人民民主专政的社会主义国家，国家的一切权力属于人民\par

\daofasym{②}党的领导的原则\par

\daofasym{③}尊重和保障人权\par

\daofasym{④}民主集中制原则\par

\daofasym{③}法治原则\par

\subsection*{第一课维护宪法权威}
\addcontentsline{toc}{subsection}{第一课维护宪法权威}

1.2治国安邦的总章程\par

（一旦考到宪法是治国安邦的总章程，那材料必须得出现国家机关，）以及国家机履行职贵，行使权\par

(的内容)\par

\textbf{※1、人民行使国家权力的机关是什么？}\par

我国的一切权力属于人民。人民代表大会是人民行使国家权力的机关。\par

注意：人民代表大会直接行使国家权力，人民间接行使国家权力。\par

\textbf{※2、我国的国家机构是怎样产生的？}\par

\daofasym{①}全国人民代表大会和地方各级人民代表大会都由民主选举产生，对人民负责，受人民监督：（人民私\par

民代表大会的关系）\par

\daofasym{②}国家行政机关、监察机关、审判机关、检察机关都由人民代表大会产生，对它负责，受它监督。（人\par

代表大会和其他国家机关的关系）\par

【知识点拨】我国的国家机关：\par

（1）权力机关：人民代表大会 最高权力机关：全国人民代表大会\par

（2）行政机关：人民政府 最高行政机关：（国务院\par

（3）监察机关：监察委员会 最高监察机关国家监察委员会\par

（4）审判机关：人民法院 最高审判机关：最高人民法院\par

（5）检察机关：人民检察院 最高检察机关：最高人民检察院\par

（人民法院和人民检察院统称为司法机关”\par

注意：中国共产党是领导核心，领导所有国家机构。但它本身不是国家机关\par

3、宪法与国家机构的关系是怎样的？（宪法如何保证国家权力稳定运行）\par

\daofasym{①}宪法通过设置国家机构，授予国家机构特定职权，明确国家机构的组成、任期、工作方式等内容，使\par

国家权力的运行稳定有序。\par

\daofasym{②}国家机构依据宪法行使权力，以实现和维护人民的根本利益。\par

\textbf{4、我国的国家机构实行什么原则？这一原则主要体现在哪些方面？}\par

※（1）原则：民主集中制\par

（2）我国国家机构贯彻民主集中制原则主要体现为：\par


\section*{PDF第43页}
\addcontentsline{toc}{section}{PDF第43页}

【漏识补录：民主集中制第（1）点】\par

（1）在国家机构与人民的关系方面，国家权力来自人民，由人民选举产生国家权力机关，国家权力机关在国家机构中居于主导地位。\par

家机构中居于生导地位：\par

\daofasym{②}在中央与地方国家机构的关系方面，中央和地方的国家机构职权的划分，遵循在中央的统一领导下，充\par

分发挥地方的主动性、积极性的原则；\par

\daofasym{③}在国家机关内部作出决策、决定时，实行民主集中制。\par

\textbf{注意：选择题经常展示一段材料，问你材料体现了民主集中制哪一个方面？}\par

\textbf{※5、为什么要规范权力运行？}\par

\daofasym{①}权力是把双刃剑，运用得好，可以造福于民，如果被滥用，则会滋生腐败，贻害无穷。必须加强对板力\par

运行的制约和监督让人民监督权力让权力在阳光下运行把权力关进制度的笼子。\par

\daofasym{②}规范国家权力运行以保障公民权利，这是宪法的核心价值追求。\par

\daofasym{③}只有依法规范权力运行，才能保证人民赋予的权力始终用来为人民谋利益。\par

※6、宪法是如何规范权力运行的？（=如何正确行使国家权力？）\par

\daofasym{①}国家权力必须任究法和法律限定的范围内行使。任何超越权限，滥用职权的行为均应承担法律责任。（法\par

无授权不或为)\par

\daofasym{②}人民通过宪法和法律将国家权力授予国家机关。国家机关及其工作人员必须依法行使权力、，履行职责，\par

不得懈怠、推诿。（法定职责必须为）\par

\daofasym{③}国家权力必须严格按照法定的途径和方式行使。（国家权力按照法定程序行使）\par

\daofasym{①}国家权力的行使不能任性，法定职责必须为，法无授权不可为。（国家权力的行使不能任性）\par

\textbf{7、为仕么要进行宪法宣誓？}\par

\daofasym{①}宪法是国家的根本大法，是治国安邦的总章程\par

\daofasym{②}树立宪法意识，维护宪法权威，弘扬宪法精神，履行法定职责。\par

\daofasym{③}国家的一切权力属于人民；\par

2.1坚持依宪治国\par

治国的基本方略：依法治国（法治）\par

依法治国首先是依宪治国\par

\textbf{1、我国宪法的构成、本质和内容各是什么？}\par

（1）构成：我国现行宪法除序言外，设有第一章总纲，第二章公民的基本权利和义务，第\par

三章国家机构，第四章国旗、国歌、国徽、首都，共四章一百四十三条。\par

（2）本质：我国宪法是党和人民意志的集中体现，是国家的根本法\par

（3）内容：宪法规定了我国的圈家性质（根本制度（根本任级（公民的基本权利和义务\par

国家机构的设置及职权等国家生活中的最根本、最重要的问题，涉及政治、经济、文化和\par

社会生活等各个方面。\par

（4）国家的根本任务：沿着中国特色社会主义道路，集中力量进行社会主义现代化建设。\par

\textbf{※2、国赛的地位？}\par

\daofasym{①}我国是党的主张和人民意志的统一，是治国安邦的总章程，是公民权利的保障书，\par

是公民、任何组织、国家机关的根本活动准则。\par

\daofasym{②}宪法具有至高无上的权威。\par

\daofasym{③}宪法是国家的根本法，在国家法律体系中具有最高的法律地位、法律权威和法律效力。\par

※3、为什么要维护宪法权威？(为什么以宪法为根本活动准则)\par

\daofasym{①}宪法是党和人民意志的集中体现，是国家的根本法。\par


\section*{PDF第44页}
\addcontentsline{toc}{section}{PDF第44页}

【漏识补录：第5题宪法是国家根本法】\par

5.为什么说宪法是国家的根本法？（宪法和其他法律的区别）\par

\daofasym{①}宪法所规定的内容是国家生活中带有全局性、根本性的问题，而其他法律所规定的内容通常只是国家生活中的一般性问题。\par

\daofasym{②}宪法具有至高无上的权威。宪法的权威关系国家的命运、社会的安定和人民的根本利益\par

如果宪法没有权威，法治的权威就树立不起来；如果宪法受到漠视，人民的权利和自由\par

无法保证。\par

※4、如何维护宪法权威？（如何以宪法为根本活动准则？）\par

\daofasym{①}坚持依法治国首先要坚持依宪治国，坚持依法执政首先要坚持依宪执政。\par

\daofasym{②}任何公民、社会组织和国家机关都必须以宪法法律为行动准则。\par

\daofasym{③}一切组织和个人都必须维护宪法权威，捍卫宪法尊严，保证宪法实施，依照宪法法律\par

使权利或权力7履行义务或职责，都不得有超越宪法法律的特权，一切违反宪法法律的\par

为都必须予以追究。\par

※【提示依法治国的核心是依宪治国。\par

通常只是国家生活中的一般性问题。\par

\daofasym{②}宪法具有最高的法律效力。宪法是其他法律的立法基础和立法依据，其他法律是根据\par

法制定的，不得与宪法的原则和精神相违背。否则，就会因违宪而无效。\par

\daofasym{③}宪法的制定和修改程序比其他法律更加严格。\par

\daofasym{①}宪法是国家法制统一的基础。全面依法治国，保障宪法实施，必须完善以宪法为核心\par

中国特色社会主文法律体系。\par

\textbf{6、宪法严格的制定和修改程序有何意义？}\par

一方面使得宪法的内容具有更广泛的民意基础，另一方面可保障宪法的长期稳定性，使国家长治\par

安，社会健康发展。\par

\textbf{7、为什么说“宪法是国家法制统一的基础”？}\par

8、\daofasym{①}全面依法治国，保障宪法实施，必须完善以宪法为核心的中国特色社会主义法律体系。\par

\daofasym{②}宪法的规定具有原则性的特点，各种法律制度是对宪法规定的具体落实。\par

\daofasym{③}宪法是对公民基本权利和根本确认和保障，其他法律也对公民基本权利的实现具有不可替代的\par

用。\par

2.2加强宪法监督\par

\textbf{※1.为什么要加强宪法监督？}\par

\daofasym{①}权力行使需要接受监督。监督是权力正确行使的根本保证，不受监督的权力将导致腐败\par

（权力的行使需要监督的原因）\par

\daofasym{②}为了保证国家机关严格按宪法和法律行使权力，需要建立健全完备的监督公权力行使\par

制度体系。在这一体系中，宪法监督具有基础性意义。\par

\daofasym{③}全面依法治国需要我们健全宪法实施和监督制度，不断加强宪法监督工作。\par

\textbf{2.宪法监督的主体和宪法监督的内容？}\par

主体：全国人大及其常委会行使监督宪法实施的职权\par


\section*{PDF第45页}
\addcontentsline{toc}{section}{PDF第45页}

\begin{note}[手写批注]
页面中部有手写批注；（此处有手写笔记但无法识别）。\par
\end{note}

内容：合宪性审查和监督。即审查法律、法规等法律文件的合宪性，不得与宪法相抵触【法】;\par

审查国家机关及其工作人员等的违宪行为，追究其违宪责任【人】。\par

※3. 如何加强宪法监督工作？（=如何健全宪法实施和监督机制？）\par

\daofasym{①}要完善全国人大及其常委会宪法监督制度，健全监督机制和程序。\par

\daofasym{②}要健全宪法解释程序机制，推进合宪性审查工作，加强备案审查制度和能力建设。\par

\daofasym{③}对于各种违反宪法的行为，都必须予以追究和纠正。\par

※4.为什么要增强宪法意识\par

宪法与我们息息相关，我们的一生都离不开宪法的保护。\par

\textbf{※5.如何增强宪法意识？}\par

\daofasym{①}学习宪法，领会宪法原则和精神。「知】\par

\daofasym{②}认同宪法。理解并认同宪法的价值。「情】\par

\daofasym{③}践行宪法。将宪法原则转化为自觉的行为准则，落实在实际行动上。【行】\par

6.增强宪法意识的具体举措？【开放题】\par

国家：设立国家宪法日和建立宪法宣誓制度\par

公民：\daofasym{①}参观五四宪法历史资料陈列馆、宪法主题公园、参加宪法诵读、宪法宣传等活动；\par

\daofasym{②}认真严肃地参加升国旗、唱国歌等仪式；\par

\daofasym{③}正确行使宪法规定的公民基本权利，自觉履行宪法规定的公民基本义务；\par

\daofasym{④}机智勇敢地与违反宪法和法律的行为作斗争，自觉维护宪法和法律的尊严。\par

7、实行宪法宣誓制度的意义？（P32）\par

\daofasym{①}有利于强化国家公职人员的宪法意识，让他们珍惜宪法赋予的权力，自觉规范自己的行为；\par

\daofasym{②}有利于捍卫宪法的地位，维护宪法权威，保证宪法实施；\par

\daofasym{③}有利于增强公民对宪法的了解和认识，促进公民忠于宪法、遵守宪法、维护宪法；\par

\daofasym{④}有利于实施依法治国的基本方略，营造宪法至上、人人信仰的社会氛围，进一步推进法治中国建\par

设。\par

****高度注意，选择\par

谁监督宪法实施 全国人大及常务委员会\par

谁解释宪法 全国人大常委会\par

谁保证本行政区内宪法的实施一一地方各级人大\par

\subsection*{第三课公民权利}
\addcontentsline{toc}{subsection}{第三课公民权利}

3.1公民基本权利\par


\section*{PDF第46页}
\addcontentsline{toc}{section}{PDF第46页}

\textbf{※1.我国宪法规定的公民基本权利有哪些？}\par

（1）政治权利和自由。包括：选举权和被选举权、政治自由和监督权等。\par

（2）人身自由。包括：人身自由不受侵犯，人格尊严不受侵犯，住宅不受侵犯，通信自由和通信秘密\par

法律保护。\par

（3）社会经济与文化教育权利。包括：财产权、劳动权、物质帮助权、受教育权、文化权利。\par

（4）其他权利。包括：平等权，宗教信仰自由，妇女、儿童和残疾人等特定人群的权利受到宪法和法\par

的特殊保障。\par

\textbf{2.我国公民享有哪些政治权利和自由？}\par

（1）选举权和被选举权\par

\daofasym{①}条件：我国年满十八周岁的公民，不分民族、种族、性别、职业、家庭出身、宗教信仰、教育程度、！\par

产状况、居住期限，都有选举权和被选举权；但是依照法律被剥夺政治权利的人除外。\par

※\daofasym{②}重要性：选举权和被选举权是公民的一项基本政治权利一行使这项权利是公民参与\par

理国家和管理社会的基础\par

(2政治自由\par

\daofasym{①}内容：我国公民有言论出版（集会人结社凝行、示威的息由。\par

\daofasym{②}重要性：公民享有政治自亩，有助于公民参与国家政治生活，充分表达自己的意愿。\par

(3)/※监鲁权\par

\daofasym{①}监督权的内容：公民对于任何国家机关和国家工作人员，有提出批评和建议的权利：\par

公民对于国家机关和工作人员的违法行为，有提出申诉、控告和检举的权利，但不得捏造或诬告陷害。\par

\daofasym{②}行使监督权的途径：通过信访、电子邮件、电话、找人大代表反映情况，以及通过新闻媒体、舆论监\par

等。\par

\daofasym{③}公民行使监督权的意义：公民依法通过各种途径和形式行使监督权，有助于国家机关及其工作人员依\par

行使权力，全心全意为人民服务：有利于建设社会主义法治国家，维护国家长治久安。\par

\daofasym{④}公民应如何正确行使监督权：公民行使监督权时，应实事求是，以事实为根据，如实反映情况，必要\par

出示证据和各种证明材料。投诉和举报时，不得捏造或者歪曲事实进行诬告陷害，也不能采用张贴大字\par

聚众闹事等方法。\par

\textbf{※3.公民依法行使政治权利和自由有什么意义？}\par

我国公民依法行使政治权利和自由，参与国家政治生活，通过各种途径和形式管理国家事务，管理经济\par

化事业，管理社会事务，这是人民行使当家作主权利的重要体现。\par

4.人身自由\par

（1）含义：人身自由是指公民的人身不受非法侵犯的自由。\par

（2）内容：包括人身自由不受侵犯，人格尊严不受侵犯，住宅不受侵犯，通信自由和通信秘密受法律\par

护。\par

※（3）重要性：人身自由是公民最基本、最重要的权秘只有在人身自由得到保障的前提工\par

公民才能独立、自由、有尊严地生活。\par

5.人身自由不受侵犯\par

宪法规定：任何公民，非经人民检察院批准或者决定或者人民法院决定，并由公安机关执行，不受逮捕\par

禁止非法拘禁和以其他方法非法剥夺或者限制公民的人身自由，禁止非法搜查公民的身体。\par


\section*{PDF第47页}
\addcontentsline{toc}{section}{PDF第47页}

6.人格尊严不受侵犯\par

（1）原因：公民都有自我尊重和受人尊重的需要，都应当享有受他人和社会尊重的权利。\par

（2）宪法规定：公民的人格尊严不受侵犯，禁止用任何方法对公民进行侮辱、谤和诬告陷害。\par

※（3）内容：公民的人格尊严权包括名誉权、荣誉权、肖像权、姓名权、隐私权等。\par

7. 住宅不受侵犯\par

（1）原因：住宅是公民享有安宁生活和安全感的私人空间。\par

（2）宪法规定：公民的住宅不受侵犯，禁止非法搜查或者非法侵入公民的住宅。\par

8.通信自由和通信秘密受法律保护\par

宪法规定：公民的通信自由和通信秘密受法律的保护。除因国家安全或者追查刑事犯罪的需要，由公安机\par

关或者检察机关依照法律规定的程序对通信进行检查外，任何组织或者个人不得以任何理由侵犯公民的通\par

信自由和通信秘密。\par

\textbf{9.我国公民的社会经济与文化教育权利包括哪些？}\par

(1)(财产权\par

\daofasym{①}重要性：我们的生存和发展及物质和文化生活需要的满足，都离不开财产。\par

\daofasym{②}宪法规定：公民前合法的私有财产不受侵犯。 买 投员\par

\daofasym{③}内容：公民可以通过合法方式取得财产，并依法占有和便用（获得收益和进衍处分\par

(2）劳动权（工资间题）\par

\daofasym{①}含义：一切有劳动能力的公民有劳动就业和取得劳动报酬的权利。\daofasym{②}宪法规定：公民有劳动的权利和义\par

务。\daofasym{③}重要性：这是公民赖以生存的基础。人们通过劳动，与社会生产与服务活动，获得劳动报酬和其\par

他收益，既可以保障合理的生活水平，实现自身价值，也为国家和社会作出贡献。\par

(3）物质帮助权（医保、社保这一些）\par

\daofasym{①}宪法规定：公民在年老、疾病或者丧失劳动能力的情况下，有从国家和社会获得物质帮助的权利。\daofasym{②}保\par

障措施：国家发展为公民享受这些权利所需要的社会保险、社会救济和医疗卫生事业。\par

※（4）受教育权（奖学、（助学金、(教育贷等）\par

\daofasym{①}含义：公民有按照其能力平等地从国家获得接受教育的机会，并获得相应物质保障的权利。\par

\daofasym{②}重要性：教育为个人人生幸福奠定基础，为人类文明传递薪火，成就民族和国家的未来。\par

\daofasym{③}保障措施：国家实行义务教育制度、国家制定资助政策。\par

（5）文化权利\par

\daofasym{①}宪法规定：中华人民共和国公民有进行科学研究、文学艺术创作和其他文化活动的自由。\par

\daofasym{②}保障措施：国家对于从事教育、科学、技术、文学、艺术和其他文化事业的公民的有益于人民的创造性\par

工作，给以鼓励和帮助。\par

**劳动和教育既是权利，又是义务\par

3.2依法行使权利\par

\textbf{※1、公民应如何依法行使权利？}\par

\daofasym{①}行使权利有界限。公民行使权利不能超越它本身的界限，不能滥用权利。公\par

民在行使自由和权利的时候，不得损害国家的、社会的、集体的利益和其他公\par

民的合法的自由和权利。\par

\daofasym{②}维护权利守程序。公民行使权利应依照法定程序，按照规定的活动方式、步\par

骤和过程进行。公民权利受到损害，要依照法定程序维护权利，维护权利的方\par

式包括和解、调解、仲裁和诉讼等\par


\section*{PDF第48页}
\addcontentsline{toc}{section}{PDF第48页}

※2.公民维权的方式有：\par

和解（和解是一种快速、简便的争议解决方式，当事人在自愿、互谅的基础上\par

依据法律，通过直接对话，达成协议，使纠纷得以解决。（适用范围：\par

一些常见的消费、劳动争议和交通事故纠纷）\par

调解：调解入以国家法律和政策及社会公德为依据，对纠纷双方进行疏导、对\par

说，促使他们相互谅解，自愿达成协议，解决纠纷。调解的方式主要有\par

人民调解又行政调解知司法调解\par

仲裁：当事人根据他们之间订立的仲裁协议，自愿将争议提交仲裁，并受仲表\par

裁判约束。（合同纠纷和财产权益净议）\par

诉讼：公民通过诉讼向人民法院起诉，依法维护自身权益。包括民事诉讼、开\par

事诉讼、行政诉讼。\par

注意：（公民受到人身关系或财产关系的争议（问题比较轻）一－民事诉讼\par

公民对于某些侵犯自己人身、财产权利的行为（问题比较重）一一刑事自诉)\par

\subsection*{第四课公民义务}
\addcontentsline{toc}{subsection}{第四课公民义务}

4.1公民基本义务\par

※1.公民的基本义务：\par

遵守宪法法律、维护国家利益、依法服兵役、依法纳税、劳动的义务、受教育的义务、\par

妻双方实行计划生育的义务、父母抚养教育未成年子女的义务和成年子女赡养辅助父母\par

义务等。\par

\textbf{2.公民要遵守宪法和法律的原因、表现及做法分别是什么？}\par

（1）原因：我国宪法和法律是全国各族人民意志和利益的集中体现，维护宪法和法律的\par

严是公民对国家和社会应尽的职责。\par

（2）表现：保护国家秘密，爱护公共财产，遵守劳动纪律，遵守公共秩序，尊重社会\par

※（3）做法：\daofasym{①}忠于宪法，维护宪法尊严，保障宪法实施。\par

\daofasym{②}自觉做到学法尊法守法用法，共同营造守法光荣、违法可耻的社会氛围。\par

\daofasym{③}自觉学习法律知识，了解法律程序规定，以法律来指导和约束自己的行为，做到依法\par

事。\par

※3.维护国家利益包括维护国家统一和全国各民族团结，维护国家安全、荣誉和利益。\par


\section*{PDF第49页}
\addcontentsline{toc}{section}{PDF第49页}

\textbf{4.维护国家安全、荣誉和利益各自包括哪些方面?}\par

\daofasym{①}维护国家安全包括维护国家的主权、领土完整不受侵犯，国家秘密不被窃取、泄露和出\par

卖，社会秩序不被破坏等。\par

\daofasym{②}维护国家荣誉包括维护国家的尊严不受侵犯，国家的荣誉不受站污。\par

\daofasym{③}维护国家利益包括维护国家的政治、经济和安全等各方面的利益。\par

巡5.我国的民族关系是：平等、团结、互助、和谐。 （已经形成)\par

6.依法服兵役\par

（1）原因：\daofasym{①}保卫祖国、抵抗侵略是公民的神圣职责。\daofasym{②}依照法律服兵役和参加民兵组织\par

是公民的光荣义务。\par

(2)我国的兵役制度：我国兵役法规定，我国实行志愿兵役为主体的志愿兵役和义务兵役\par

相结合的兵役制度。（教材改动）\par

（3）兵役分类：现役和预备役\par

7.依法纳税\par

（1）重要性：\daofasym{①}税收是国家财政收入的主要来源，依法纳税是公民的一项基本义务。\par

（2）违反纳税义务的行为：偷税、欠税、骗税、抗税\par

（3）任何偷税、欠税、骗税、抗税的行为都是违法行为，情节严重、构成犯罪的要依法追\par

究刑事责任。\par

4.2依法履行义务\par

\textbf{※1.权剩与必雾的关系是怎样的？}\par

\daofasym{①}公民的权利与义务相互依存、相互促进。权利的实现需要义务的履行，义务的履行促进\par

权利的实现。\par

\daofasym{②}公民既是合法权利的享有者，又是法定义务的承担者。\par

\daofasym{③}公民的某些极利同时也是义务。\par

【特别提醒】劳动和受教育既是公民的基本权利，也是公民的基本义务。\par

\textbf{※2.如何变确处理权利与义务的关系？}\par

\daofasym{①}坚持权利和义务相统一，任何公民既不能只享受权利而不承担义务，也不应只承担义务\par

而不享受权利。阳阳老师\par

\daofasym{②}我们不仅要增强权利意识，依法行使权利，而且要增强义务观念，自觉履行法定的义务。\par

3.公民行使权利和履行义务有什么意义\par

（1）行使权利：公民权利的充分实现，可以激发公民的主人翁意识，调动其履行义务的积\par

极性和主动性，自觉承担对国家和社会的责任。\par

（2）履行义务：公民自觉履行义务，促进国家发展和社会进步，又为其权利的实现提供和\par

创造更好的条件。\par

\textbf{4.为什么必须履行法定义务？}\par

\daofasym{①}法定义务是由我国宪法和法律规定的，具有强制性。\par


\section*{PDF第50页}
\addcontentsline{toc}{section}{PDF第50页}

\daofasym{②}自觉履行法定义务，是公民不可推卸的责任。\daofasym{③}建设富强、民主、文明、和谐、美丽的\par

社会主义现代化强国需要我们每个中国人的奋斗，需要我们自觉履行法定义务。\par

\textbf{※5.如何履行法定义务？}\par

\daofasym{①}法律要求做的必须去做。\par

\daofasym{②}法律禁止做的坚决不做。\par

\textbf{6.什么是违反法定义务的行为？}\par

公民实施了法律所禁止的行为，或者没有实施法律要求做的行为。\par

\textbf{7.违反法定义务必须承担怎样的责任？}\par

违反法定义务，必须依法承担相应的法律责任。\par

\daofasym{①}违反民事法律，应当依法承担民事责任；\par

\daofasym{②}违反行政法律，应当依法承担行政责任；\par

\daofasym{③}违反刑事法律，构成犯罪的，应当依法承担刑事责任。\par

9、生活中有些事情（如义工等）虽然不是法律要求我们必须履行的义务，但这是法律房\par

鼓励的，他是社会主义道德的要求，是我们应该履行的道德义务。\par

重要签题框架： 从义间度床\par

\textbf{材料中只有涉及到谁谁谁违法了，侵权了，受到法律制裁体现了什么？}\par

\daofasym{①}法定义务是我国宪法和法律规定的，具有强制性。谁谁谁有 x权利，但也有履行 x 的\par

义务。\par

\daofasym{②}我们要自觉履行法定义务。法律要求的必须去做，法律禁止的坚决不做。\par

\daofasym{③}违反法定义务就要承担相应的法律责任。材料中，谁谁违反了民事法律，应当依法承担\par

民事责任/违反行政法律，应当依法承担行政责任/违反刑事法律，构成犯罪的，应当依法\par

承担刑事责任。\par

\subsection*{第五课 我国的政治和经济制度}
\addcontentsline{toc}{subsection}{第五课 我国的政治和经济制度}

5. 1 基本经济制度\par

\textbf{※1、我国的基本经济制度是什么？确立依据是？}\par

（1）内容：公有制为主体、多种所有制经济共同发展。\daofasym{②}按劳分配为主体、多种分配方式并存。\daofasym{③}\par

会主义市场经济体制。\par

（2）依据：我国正处于社会主义初级阶段。（我国的基本国情）\par

\textbf{※2、公有制经济包括哪些？}\par

（1)国有经济。\daofasym{①}地位：国有经济是国民经济的主导力量\daofasym{②}含义：国有经济的生产资料属于全体人\par


\section*{PDF第51页}
\addcontentsline{toc}{section}{PDF第51页}

法同所有。\daofasym{③}作用：国有经济主要分布在关系国家安全、国民经济命脉的重要行业和关键领域、重点提\par

供公共服务、发展重要前略性战略产业、保护生态环境、支持科技进步，保障国家安全。\par

（2）集体经济。含义：集体经济的生产资料属于一部分劳动者共同所有、（有利于实现共同富裕）\par

（3）温合所有制经济中的国有成分和集体成分。\par

【说明】公有制是主体，国有经济是主导，选择题经常考，注意区分\par

\textbf{然3、非公有制经济包括哪些？}\par

个体经济、私营经济外资经济等。\par

（重点）非公有制经济的作用有利手增加就业，增加人民收入，增加国家税收一促进经济社会发展，\par

推动大众创业万众创新。\par

\textbf{x4、为开么要坚持公有制为主体、多种所有制经济共同发展？}\par

\daofasym{①}在我国，公有制经济与非公有制经济都是社会主义市场经济的重要组成部分，都是我国经济社会发展的\par

重要基础。（地位）\par

\daofasym{②}坚持公有制为主体、多种所有制经济共同发展，促进了生产力的发展、综合国力的增强和人民生活水平\par

的提高、为人民当家作主奠定了坚实的物质基础。（作用）\par

\textbf{※5、对于公有制和罪公有制经济，我国的国家政策是怎样的？}\par

必须毫不动摇巩固和发展公有制经济，毫不动摇鼓励、支持、引导非公有制经济发\par

展。（选择题要注意区分）非hC业)[的动度\par

\textbf{6、按劳分配的依据、基本内容和要求分别什么？}\par

（1）依据：按劳分配是由生产资料公有制快定的。\par

（2）基本内容和要求：有劳动能力的会成员必须参加劳动：在作了必要的扣除后，以劳动者所提供的\par

劳动(包括劳动数量和质量)为尺度对个人进行分配，多劳多得，少劳少得。\par

注意：按劳分配一定是存在在公有制经济中。\par

\textbf{7、除了按劳分配外，我国还存在哪些分配方式？}\par

\daofasym{①}让劳动、资本土也、（知识技管理数授等生产要素参与收入分配，由市场评价贡献，按贡献决\par

定报配。（优点：提高效率，调动积极性。缺点：可能会拉大贫富差距。）\par

\daofasym{②}在一定条件下，人们还可以取得社会保障收入，获得社会公益事业的帮助。\par

公有制企业（国有企业和集体企业）获得的工资、奖金、福利、津贴都属于按劳分配\par

非公有制企业获得工资、奖金、福利、津贴属于按生产要素里面的劳动要素分配\par

投资等属于按生产要素的资本分配\par

劳动收入=公有制里面的按劳分配和非公有制里面的按劳动要素分配（只要通过劳动获得的都算）\par

\textbf{8、什么是市场经济体制？}\par

市场经济体制是市场在资源配置中起决定性作用的经济体制。\par

※9、市场经济中，市场是如何配置资源的？（无形的手）市场在资源配置中起决定性作用\par

\daofasym{①}在市场经济中，生产什么、如何生产、生产多少，主要是通过价格、供求、竞争等机制来调节。\par


\section*{PDF第52页}
\addcontentsline{toc}{section}{PDF第52页}

\daofasym{②}市场机制就像一只“看不见的手，在资源配置中起决定性作用。\par

\textbf{※10、我国的社会主义市场经济体制有何优越性？}\par

我国社会主义市场经济体制把社会主义制度和市场经济有机结合起来，充分发挥市场在资源配置中的决定\par

性作用。更好发挥政府作用，进行科学宏观调控，激发各类市场主体的活力，为人民对美好生活的需求提\par

供保障。\par

几个易错点：\par

国有经济是主导，公有制经济是主体\par

2、国民经济是整全国家的经济公+非公\par

、国民经济的支村国有经流\par

4、社会主义市场经济二公+非公\par

（公有制经济和非公有制经济都是社会主义市场经济的重要组成部分。高度注意：不能写成社会主义经济\par

、市场在资源配置中起决定性作用\par

6、我国经济制度的基础=公有制经济\par

作该治因而机统\par

\subsection*{第五课 我国的政治和经济制度}
\addcontentsline{toc}{subsection}{第五课 我国的政治和经济制度}

5. 2 根本政治制度\par

※1.我国的根本政格制度是：（人民代表大会制度。\par

※2.人民代表大会制度的基本内容包括：\par

\daofasym{①}国家的一切权力属于人民\par

\daofasym{②}人民通过民主选举出代表，组成各级人民代表大会作为国家权力机关。\par

\daofasym{③}由人民代表大会产生国家行政机关、监察机关、审判机关、检察机关，这些国家机关位\par

法行使各自的职权，并对人民代表大会负责，受人民代表大会监督。\par

\daofasym{④}实行民主集中制，重大问题经人民代表大会充分讨论，遵循少数服从多数原则，民主社\par

定。\par

※3.人民代表的职权：有权依法审议各项议案和报告、表决各项决定、提出议案和质询案\par

（审议权、表决权、提案权、质询权）\par

4.人大代表如何对人民负责？（义务）\par

\daofasym{①}人大代表必须与尺民群众保持密切联系，听取和反映人民群众的意见和要求。\par

\daofasym{②}努力为人民服务！对人民负责，并接受人民监督。\par

5.为什么要坚持和完善人民代表大会制度？（优越性）\par

\daofasym{①}实践充分证明，人民代表大会制度是符合中国国情和实际、体现社会主义国家性质、\par

证人民当家作主、保障实现中华民族伟大复兴的好制度。\par

\daofasym{②}人民代表大会制度是坚持党的领导、人民当家作主、依法治国有机统一的根本政治制\par

安排。\par

\daofasym{③}活力人民代表大会制度建立以来，不断得到巩固和发展，展现出蓬勃生机。\par

\textbf{※6.如何坚持和完善人民代表大会制度？}\par

\daofasym{①}必须毫不动摇坚持中国共产党的领导。必须毫不动摇地坚持中国共产党的领导，通过>\par

民代表大会制度，使党的主张通过法定程序成为国家意志。\par

\daofasym{②}必须保证和发展人民当家作主。必须保障和发展人民当家作主，支持和保证人民通过>\par


\section*{PDF第53页}
\addcontentsline{toc}{section}{PDF第53页}

民代表大会行使国家权力，扩大人民民主，健全民主制度，丰富民主形式，拓宽民主渠道。\par

\daofasym{③}必须全面推进依法治国。通过人民代表大会制度，弘扬社会主义法治精神，实现国家各\par

项工作法治化。\par

\daofasym{①}必须坚持民主集中制。人民代表大会统一行使国家权力，国家机关既有合理分工又有相\par

互协调，保证国家统一高效推进各项事业。\par

\subsection*{第五课 我国的政治和经济制度}
\addcontentsline{toc}{subsection}{第五课 我国的政治和经济制度}

5. 3 基本政治制度\par

\textbf{然1.我国的基本政治制度有哪些？}\par

\daofasym{①}中国共产觉领显的多党合作和政治协商制度（我国的政党制度）\daofasym{②}民族区域自治制度\daofasym{③}基层群众自治制度\par

※2. 多党合作的基本方针长期共存、互相监督、肝胆相照、荣辱与共\par

x3.中国共产党：\par

\daofasym{①}地位：中国特色社会主义事业的领导核心\par

\daofasym{②}性质：中国共产党是工人阶级领导的先锋队，同时是中国人民和中华民族的先锋队\par

\daofasym{③}宗旨：全心全意为人民服务\par

\daofasym{①}党的初心和使命：为中国人民谋幸福，为中华民族谋复兴。\par

※4.多党合作的重要机构：中国人民政治协商会议（人民政协）\par

※5．中国人民政治协商会议（简称：人民政协）：\par

\daofasym{①}性质：多党合作和政治协商的重要机构，爱国统一战线组织\par

\daofasym{②}主题，团和民主\par

\daofasym{③}职能/政治协商民主监复参政议购\par

\textbf{6.我国恶成子忘样的爱国统子战线？}\par

在长期的革命、建设、改革过程中，已经结成由中国共产党领导的，有各民主党派和各人民团体参加的，包括\par

全体社会主义劳动者、社会主义事业的建设者、拥护社会主义的爱国者、拥护祖国统一和致力于中华民族伟大复兴\par

的爱国者的广泛的爱国统一战线。\par

※Z坚持中国共产党领导的多党令作和政治协商制度的意义？（我国政党制度的优越性？)\par

\daofasym{①}是发扬社会主义民主的重要形式，有利于反映民意，集中民智，促进科学民主决策\par

\daofasym{②}有利于协调关系，化解矛盾，维护社会稳定和谐\par

\daofasym{③}有利于凝聚人心，反对分裂推进祖国和平统一大业\par

8.民族区域自治制度的内容！\par

\daofasym{①}民族自治地方的分级：我国民族自治地方分为自治区、自治州、自治县三级。\par

※\daofasym{②}民族自治机关：民族自治地方的人民代表大会和人民政府是自治机关。\par

\daofasym{③}自治机关的职权：在行使一般地方国家机关职权的同时，依法行使自治权，即根据本地方、本民族政治、经济、\par

社会、文化等方面的特点，自主管理本地方、本民族的内部事务。\par

\daofasym{④}民族自治地方与中央的关系：我国民族区域自治是在国家统一领导下的自治，各民族自治地方是国家不可分割的\par

组成部分，民族自治机关必须服从中央的领导。\par

【说明1】民族区域自治不能说成“高度自治”！【说明2】社会主义民族关系：平等团结互助和谐（已经形成）！\par

9.实行民族区域自治的意：\par

\daofasym{①}利于把国家的集中、统一与各民族的自主、平等结合起来\par

\daofasym{②}有利于把国家的法律、政策与民族自治地方的具体实际、特殊情况结合起来\par

\daofasym{③}有利于把各民族人民热爱祖国的感情与热爱自己民族的感情结合起来\par

证明民族区域直治度行合我国国情，维护国家统一 加强民族团纪促进民认地口发展，\par

消强中华民族凝聚力等方面起到重要们用，（新寺\par

※10.基层群众自治制度：\par

\daofasym{①}含义：由居民或村民分别选举产生居民委员会或村民委员会，实现自我管理、自我教育、自我服务、自我监督。\par

※\daofasym{②}基层群众性自治组织（形式）：城市的居民委员会和农村的村民委员会【不是基层国\par


\section*{PDF第54页}
\addcontentsline{toc}{section}{PDF第54页}

家政权机关！还有党，政协都不是国家机关！！！\par

\daofasym{③}居民委员会成员由居民直接投票选举产生，重要事务提请居民委员会讨论决定，居民委员会实行办事公开制度，\par

接受居民的监督。\par

※11.基层群众自治制度的意义：\par

直利于人民群众直接行使民主权利）管理基层公共事务和公益事业，推动社会主义民主建设，促进社会和谐稳\par

定。【注意】这个地方可以用“直接”！\par

我国的制度优势（国家制度）\par

\daofasym{①}社会主义制度是中国人民共和国的根本制度，本质是保障人民当家作主，能够集中力量办大事，维护\par

人民根本利益\par

\daofasym{②}中国共产党的领导是中国特色社会主义制度的最大优势和最本质特征，中国共产党是最高自治领导力\par

量。\par

\daofasym{③}社会主义基本经济制度（三个一一公有制为主体，多种所有制经济共同发展；按劳分配为主体，多种分\par

配方式并存：社会主义市场经济体制，考试可能三个都涉及到，具体得看材料）\par

\daofasym{④}社会主义基本政治制度（三个一中国共产党领导的多党合作和政治协商制度；民族区域自治制度；基\par

层群众自治制度，考试可能都涉及到，具体得看材料）\par

\daofasym{③}人民代表大会制度是我国的根本政治制度（具体展开）\par

\subsection*{第六课 我国国家机构}
\addcontentsline{toc}{subsection}{第六课 我国国家机构}

6.1国家权力机关\par

※1.人民行使国家权力的机关是：全国人民代表大会和地方各级人民代表大会。\par

2.全国人民代表大会：\par

※\daofasym{①}性质：最高国家权力机关\par

\daofasym{②}地位：在我国国家机构中居于最高地位\par

\daofasym{③}与其它中央国家机关的关系：其他中央国家机关都由它产生，对它负责，并受它监督。\par

地方各级人民代表大会：\par

※\daofasym{①}性质：地女国家权力机关\par

\daofasym{②}地位：它是本行政区域内人民行使国家权力的机关\par

\daofasym{③}与本级其他国家机关的关系：地方国家行政机关、监察机关、审判机关和检察机关都由本级人民代表大\par

会产生，对它负责，受它监督。\par

※3.人民代表大会的职权有：\par

（）立法权：\daofasym{①}关键词：通过…法律、制定……法律、修改…·法律等。制定、修改、废除)\par

\daofasym{②}全国人大及其常委会一一行使国家立法权（修改宪法、基本法律）\par

\daofasym{③}县级以上各级人大及其常委会一一行使地方立法权（制定地方性法规）\par

（2）决定权一一关键词：决定国家和社会或本行政区域内重大事项的权力\par

（3）任免权一一关键词：选举、决定、免的权力\par

（4）监督权一一关键词：（全国人大及常委会）监督宪法和法律的实施，\par

监督“一府一委两院”工作的权力\par


\section*{PDF第55页}
\addcontentsline{toc}{section}{PDF第55页}

【结构整理：我国国家机构】\par

国家机构包括：国家权力机关、中华人民共和国主席、行政机关、监察机关、司法机关等，是一个严密的组织体系。\par

【结构整理：中央国家机构组织系统简表】\par

全国人民代表大会居于最高地位；国家主席、全国人大常委会、国务院、中央军事委员会、国家监察委员会、最高人民法院、最高人民检察院等由其产生或与其相连。\par

国务院是最高行政机关（中央人民政府）；中央军事委员会是最高军事机关；国家监察委员会是最高监察机关；最高人民法院是最高审判机关；最高人民检察院是最高检察机关。\par

其他中央国家机关由全国人民代表大会产生，对它负责，受它监督。\par

【结构整理：我国人民如何行使国家权力】\par

人民通过民主选举人大代表，组成人民代表大会；人民代表大会作为国家权力机关统一行使国家权力；由人民代表大会产生行政机关、监察机关、审判机关、检察机关等国家机关，具体管理国家和社会事务。\par

我国国家机构\par

国家权力 中华人民共、行政 监察 司法\par

机关 机关 机关 机关\par

国家机构是为实现国家职能而建立起来的国家机关的统称，具体包括\par

国家权力机关、国家主席、行政机关、监察机关、司法机关等，是\par

个严密的组织体系。既是人民意志的执行者，又是人民利益的捍卫者。\par

中央国家机构组织系统简表\par

全国人民代表大会\par

国家主席 全国人大常委会\par

国家监察委员会最高人民检察院最高人民法院国务院 中央军事委员会\par

中央人民政府\par

最高监察机关 最高检察机关 最高审判机关 最高行政机关 最高军事机关\par

各部、委员会\par

（其他中央国家机关由它产生、对它负责、受它监督）\par

\textbf{我国人民是如何行使国家权力的？}\par

人民 间接行使\par

选举\par

人大代表 直接行使\par

组成\par

人民代表大会 统一行使\par

(国家权力机关)\par

产\par

生\par

行政机关、监察机关、 管理国家\par

审判机关、检察机关 具体行使 和社会的权力\par


\section*{PDF第56页}
\addcontentsline{toc}{section}{PDF第56页}

【结构整理：人民代表大会和人大代表的职权】\par

人民代表大会的职权：立法权、决定权、任免权、监督权。\par

人大代表的职权：审议权、表决权、提案权、质询权。\par

不要混淆：全国人大及其常委会依法行使最高立法权、最高决定权、最高任免权、最高监督权；人大代表个人行使审议权、表决权、提案权、质询权。\par

3.中华人民共和国主席、副主席的任职条件：\daofasym{①}有选举权和被选举权；\daofasym{②}年满四十五周岁；\daofasym{③}中华人民共和国公民。\par

国家主席的职权：\daofasym{①}公布法律、发布命令；\daofasym{②}任免权；\daofasym{③}外事权；\daofasym{④}授予荣誉权。\par

人民代表大会和人大代表的职权\par

人民代表大会职权立法权 审议权\par

决定权不要混淆表决权人大代表的职权\par

任免权 提案权\par

监督权 质询权\par

注意：最高立法权、最高决定权、最高监督权、最高任免权一一全国人大及常\par

一，）\par

6.2中华人民共和国主席\par

\textbf{1.国家主席是怎样产生的？}\par

国家主席由全国人民代表大会选举产生\par

\textbf{※2.国家主席的性质是什么？}\par

\daofasym{①}国家主席代表中华人民共和国的国家机关\par

\daofasym{②}是中华人民共和国的国家元首\par

\daofasym{①}有选举权和被选举权\par

\daofasym{②}年满四十五周岁\daofasym{③}中华人民共和国公民 外发\par

4.中华人民共和国主席、副主席的任期：\par

中华人民共和国主席、副主席每届任期同全国人民代表大会每届任期相同。（全国人民代表大会五年一届\par

※5.国家主席的职权：\par

\daofasym{①}公布法律、发布命令\daofasym{②}任免权外事权、授予荣誉权（新教材）\par

※第六课我国国家机构\par


\section*{PDF第57页}
\addcontentsline{toc}{section}{PDF第57页}

【漏识补录：第2题行政机关外部监督要求】\par

为了防止工作人员出现履职不力、监管缺失、失职渎职、徇私枉法等问题，必须加强对行政权的监督和制约。\par

6.3国家行政机关\par

1.行政机关（国家权力机关的执行机关）：\par

\daofasym{①}含义：行政机关是依据宪法设立的，依法行使国家行政职权，组织和管理国家行政事务\par

的国家机关。\par

\daofasym{②}组成：我国行政机关由国务院及其领导的地方各级人民政府组成。\par

\daofasym{③}县级以上人民政府设置若干部门，这些部门分工不同，分别在各自职权范围内为人民群\par

众提供服务，他们相互配合，共同致力于提高为经济发展服务、为人民服务的能力和水平。\par

2. 职权\par

(1) 内容\par

\daofasym{①}根据宪法，我国县级以上人民政府主要管理经济、教育、科学、文化、卫生、体育事业、\par

城乡建设事业和财政、民政、公安、民族事务、司法行政、计划生育等行政工作。\par

\daofasym{②}行政机关层级不同，职权也不同。国务院，即中央人民政府，是我国最高行政机关，统\par

一领导地方各级人民政。\par

例：中央人民政府一广东省人民政府一一深圳市人民政府\par

（2）要求（行政机关如何依法行政)\par

\daofasym{①}（从行政机关自身来说）行政机关必须依法行政。行政机关享有哪些职权，具体如何行\par

使职权，都必须严格遵照宪法和法律。行政机关要坚持法定职责必须为，法无授权不可为，\par

勇于负责、敢于担当，坚决纠正不作为、乱作为，坚决克服懒政、怠政。\par

\daofasym{②}（从外部监督来说）行政机关的职权通常由行政机关工作人员具体实施。为了防止工作\par

制约，切实做到有权必有责，用权受监督，权责要对等，失责要追究，侵权要赔偿。\par

\textbf{3.行政机关与权力机关的关系？}\par

\daofasym{①}国家行政机关是国家权力机关的执行机关，负责贯彻执行国家权力机关通过的有关法律、\par

决议和决定。\par

\daofasym{②}各级人民政府对本级人民代表大会负责并报告工作。\par

\daofasym{③}县级以上地方各级人民政府在本级人民代表大会闭会期间，对本级人大常委会负责并报\par

告工作。\par

\textbf{4.建设什么样的政府？}\par

\daofasym{①}建设服务型政府。行政机关必须全心全意为人民服务，努力建设人民满意的服务型政府。\par

\daofasym{②}建设依法行政的政府。行政机关享有哪些职权，具体如何行使职权，都必须严格遵照宪\par

法和法律。行政机关要坚持法定职责必须为，法无授权不可为，勇于负责、敢于担当，坚\par

决纠正不作为、乱作为，坚决克服懒赖政、怠政。\par

\daofasym{③}建设敢于当的政府。\par

\daofasym{④}建设清正廉洁的政府\par

\daofasym{③}建设阳光型政府，做到政务公开，自觉接受人民监督。\par

（为什么要建设这样的政府？因为人民是国家的主人，行政机关的权力来自人民，是人民\par

授予的。）\par


\section*{PDF第58页}
\addcontentsline{toc}{section}{PDF第58页}

\textbf{5.行政机关与人民的关系?}\par

\daofasym{①}人民是国家的主人，行政机关的权力来自人民的授予\par

\daofasym{②}行政机关必须全心全意为人民服务，努力建设人民满意的服务型政府。（国家性质决定）\par

\subsection*{第六课 我国国家机构}
\addcontentsline{toc}{subsection}{第六课 我国国家机构}

6.4.国家监察机关\par

国家监察机关=国家监察委员会（最高级别的）+地方各级监察委员会\par

1.监察委员会\par

※\daofasym{①}性质：中华人民共和国各级监察委员会是国家的监察机关。\par

\daofasym{②}监察范围：监察委员会对所有行使公权力的公职人员进行监察，实现国家监察全面覆盖。\par

※\daofasym{③}组成：国家监察委员会（最高监察机关）和地方监察委员会\par

※\daofasym{④}产生：国家监察委员会由全国人民代表大会产生，对它负责，并受它监督。\par

地方监察委员会由地方各级人民代表大会产生，对它负责，并受它监督。同时还对上一级\par

监察委员会负责，受它监督。\par

\textbf{※2.监察委员会如何行使职权？}\par

\daofasym{①}监察委员会依照法律规定独立行使监察权，不受行政机关、社会团体和个人的干涉。\par

\daofasym{②}监察机关办理职务违法和职务犯罪案件，应当与审判机关、检察机关、执法部门互相配\par

合、互相制约。\par

\daofasym{③}监察机关在工作中需要协助的，有关机关和单位应当根据监察机关的要求依法予以协助。\par

\textbf{※3.监察委员会职责？}\par

\daofasym{①}监督：监察委员会首要职责。监督所有公职人员行使公权力，确保权力不被滥用，确\par

保权力在阳光下运行\par

\daofasym{②}调查经常性工作。调查公职人员涉嫌职务违法和职务犯罪，能够减少和遏制腐败行为\par

的发生，保持公权力行使的廉洁性。\par

\daofasym{③}处置：对违法的公职人员依法作出政务处分决定；对履行职责不力、失职失责的领导\par

人员进行问责；对涉嫌职务犯罪的，将调查结果移送人民检察院依法审查、提起公诉；向\par

监察对象所在单位提出监察建议。\par

\textbf{4、为何要进行监察体制的改革？}\par

\daofasym{①}有利于规范权力运行\par

\daofasym{②}有利于构建监督体系，反腐倡廉。\par


\section*{PDF第59页}
\addcontentsline{toc}{section}{PDF第59页}

\daofasym{③}有利于推进党和国家治理体系及治理能力现代化\par

6.5国家司法机关——人民法院和人民检察院\par

1.人民法院（审判机关)\par

（1）性质：我国宪法规定：“中华人民共和国人民法院是国家的审判机关。”\par

（2）组织体系：我国设立最高人民法院、地方各级人民法院（分为基层、中级、高\par

级人民法院）和专门人民法院（分为军事、海事人民法院）。\par

(3）基本职权：审理刑事、民事和行政案件。\par

遂（4）作用：人民法院通过行使国家审判权，惩办犯罪分子，解决民事和行政争议，\par

维护社会秩序，引导公民自觉遵守宪法和法律。\par

（5）人民法院在司法活动中的要求(=人民法院应该如何行使职权？）\par

人民法院在司法活动中，必须坚持以事实为依据，以法律为准绳，依法独立公\par

正行使审判权，不受行政机关、社会团体和个人的干涉。人民法院对公民权利提供\par

有效救济和保障，捍卫社会公平正义。\par

2.人民检察院（法律监督机关)\par

※（1）性质：我国宪法规定：“中华人民共和国人民检察院是国家法律监督机关。”\par

（2）组织体系：我国设立最高人民检察院、地方各级人民检察院和军事检察院等专\par

门人民检察院。\par

※（3）职权：\par

人民检察院通过行使检察权，追诉犯罪，维护国家安全和社会秩序，维护个人\par

和组织的合法权益，维护国家利益和社会公共利益，保障法律正确实施，维护社会\par

公平正义，维护国家法制统一、尊严和权威，保障中国特色社会主义建设的顺利进\par

行。\par

※(4)人民检察院在司法活动中的要求（=人民检察院应该如何依法行使职权？）\par

人民检察院依照法律规定独立行使检察权，不受行政机关、社会团体和个人的\par

干涉。各级人民检察院的工作人员，必须忠实于事实真相，忠实于法律，忠实于社\par

会主义事业，全心全意为人民服务。\par

\textbf{3、如何做到公正司法？}\par

（1）人民法院独立公正行使审判权。人民法院在司法活动中，必须坚持以事实为依\par

据，以法律为准绳，依法独立公正行使审判权，不受行政机关、社会团体和个人的\par

干涉。人民法院对公民权利提供有效救济和保障，捍卫社会公平正义。\par

（2）人民检察院独立行使检察权。人民检察院依照法律规定独立行使检察权，不\par

受行政机关、社会团体和个人的干涉。各级人民检察院的工作人员，必须忠实于事\par

实真相，忠实于法律，忠实于社会主义事业，全心全意为人民服务。\par


\section*{PDF第60页}
\addcontentsline{toc}{section}{PDF第60页}

\textbf{4、法院和检察院都是司法机关，两者有什么区别呢？}\par

法院是审判机关。\par

检察院是检察机关，同时也是公诉案件的审查起诉机关。依法对法院审判过程进\par

行监督，并且有权对法院的违法审判提出纠正意见和提出抗诉。\par

重要提醒：\par

法院是审判机关，行使审判权！！\par

人民检察院是法律监督机关！！！行使检察权！！！！\par

\subsection*{第七课尊重自由平等}
\addcontentsline{toc}{subsection}{第七课尊重自由平等}

7.1 自由平等的真谛\par

1. 自由:\par

\daofasym{①}含义：主要指人们在法律规定的范围内，依照自己意志活动的权利。（很重要。）\par

\daofasym{②}意义（价值）：拥有自由，能增强个人的幸福感，还能激发每个人的活力。推动社会的进步与繁荣。\par

2.为什么要对自由限制？.（为什么自由是相对的，而不是绝对的）\par

\daofasym{①}自由不是为所欲为，它是有限制的、相对的。\par

\daofasym{②}必要的限制是对自由的保护。无限制的自由只会走向自由的反面，导致混乱与伤害。\par

\daofasym{③}无论是现实世界还是网络空间，自由都是法律之内的自由。\par

※3. 法治与自由的关系：\par

\daofasym{①}法治与自由相互联系，不可分割。\par

\daofasym{②}法治标定了自由的界限，自由的实现不能触碰法律的红线，违反法律可能付出失去自由的代价\par

\daofasym{③}法治是自由的保障，人们合法的自由和权利不受非法干涉和侵害。\par

\daofasym{④}法治既规范自由又保障自由。\par

4.平等的含义：一是同等情况同等对待；二是不同情况差别对待。\par

5.平等的意义：法律面前人人平等，是社会进步文明的标志，也是社会主义法治的基本原则之一。\par

※6.法律面前人人平等是如何体现的？（表现)\par

\daofasym{①}任何公民，都一律平等地享有宪法和法律规定的各项权利，同时必须平等地履行宪法和法律规定的各项\par

义务。\par

\daofasym{②}我国公民的合法权益一律平等地受到法律保护，违法或犯罪行为一律平等地依法予以追究，任何组织或\par

个人都不得有超越宪法和法律的特权。\par

提醒：类似在网络上肆意发表言论等的行为，可以从三个角度分析：\par

\daofasym{①}权利与义务的关系\par

\daofasym{②}法治与自由的关系\par


\section*{PDF第61页}
\addcontentsline{toc}{section}{PDF第61页}

\daofasym{③}规则与自由的关系\par

注意：生活中一些不平等的现象：地域、性别歧视等。这些现象的危害：损害\par

了公民的人格尊严，违背了法律面前人人平等。\par

7. 2自由平等的追求\par

\textbf{※1.如何珍视自由？}\par

\daofasym{①}思想上：珍视自由，就要珍惜宪法和法律赋予我们的权利。具体要求：我们要知晓自己的权利，正确认\par

识权利的价值，积极行使和维护自己的正当权利。\par

\daofasym{②}行动上：珍视自由。必须依法行使权利。具体要求：作为公民，应自觉守法、遇事找法、解决问题靠法,\par

树立守法光荣、违法可耻的法制意识。\par

\textbf{※2.如何践行平等？}\par

\daofasym{①}践行平等，就要反对特权。法律的尊严和权威不容侵犯，任何践踏法律的行为必然受到制裁和惩罚。每\par

个人都不得享存不受法律约束的特权。\par

\daofasym{②}践径就要平等对待他人的合法权利。我们要以法律为基本的行为准则，平等地对待所有成员，尊\par

重他人的合法权益\par

\daofasym{③}践行平等，就要敢于抵制不平等的衍为。面对这些不平等现象，我们不能听之任之，应据理力争，必要\par

时依法维权。\par

\daofasym{①}我们要增强平等意识，努力践行平等，共同构建平等有序的社会制度。\par

\subsection*{第八课维护公平正义}
\addcontentsline{toc}{subsection}{第八课维护公平正义}

8.1公平正义的价值\par

1. 公平\par

（1）含义：公平通常指人们基于一定标准或原则，处理事情合情合理、不偏不倚的\par

态度或行为方式。\par

（2）内涵：公平有着丰富的内涵，包括权利公平、规则公平、机会公平。\par

\daofasym{①}权利公平：要求每个人依法平等参与社会活动；\par

\daofasym{②}规则公平：要求每个人都受到行为规范的约束；\par

\daofasym{③}机会公平：要求社会为每个人提供同等的发展机会和条件。\par

※(3)意义\par

\daofasym{①}公平是个人生存和发展的重要保障。公平不仅能保证个人应得的利益，使个人获\par

得生存和发展的物质条件，而且能让人感受到尊严，从而激发自身潜能，提高工作\par

6/\par


\section*{PDF第62页}
\addcontentsline{toc}{section}{PDF第62页}

效率\par

@公平是社会稳定和进步的重要基础公平有利于协调社会各方面的利益关系，\par

和社会矛盾，减少社会沁突，维护社会秩序，保证社会的长治久安。公平有利于营\par

造更好的竞争环境，创造更多的社会财富，推动社会持续发展。\par

2. 正义\par

（1）含义和内涵：正义是社会文明的尺度，体现了人们对美好社会的期待和追求。\par

正义有着极其丰富的内涵，一般而言，正义行为都是有利于促进社会进步、维护公\par

共利益的行为。\par

※（3）要求\par

\daofasym{①}要求依法保障人们的正当权利，使受害者得到救济、违法者受到惩罚：（受害者)\par

\daofasym{②}要求人们分辨是非，惩恶扬善，维护社会公共利益；（全社会）\par

\daofasym{③}要求人们对弱者给予必要的扶助，以保证其有尊严地生存。（弱者）\par

※(4) 重要性\par

\daofasym{①}正义是社会文明的尺度D是法治追求的基本价值目标之之\par

\daofasym{②}正义是社会制度的重要价值。（正义作为人类社会的理想追求，往往被当作评价社\par

会制度的价值标准制度的生命在于正义，制度的效用在于保障正义。有了正义的\par

制度，即使是社会弱势群体，也能获得基本的社会保障\par

\daofasym{③}正义是社会和谐的基本条件。切实维护和实现社会正义，有利手恰当地调整和处\par

理人与人之间的关系，充分发挥人们的主动性和创造性：有利于营造和谐、稳定\par

安宁的社会环境，为社会发展注入不竭的动力。\par

\daofasym{①}正义是法治追的基本价值日标之\par

\subsection*{第八课 维护公平正义}
\addcontentsline{toc}{subsection}{第八课 维护公平正义}

8.2 公平正义的守护\par

\textbf{※1.如何坚守/维护社会公平？}\par

（1）个人角度：\par

\daofasym{①}面对利益冲突，我们要站在公平的立场，学会担当，以公平之心为人处世。\par

\daofasym{②}遇到不公平的行为时，我们要坚守原则立场，敢于对不公平说“不”，采用合理合法的方式和手段，谋\par

求最大限度的公平，努力营造一个公平的环境。\par

\daofasym{③}树立公平意识，同破坏公平的行为作斗争，对受害者伸出援助之手。\par

（2）制度角度：\par

\daofasym{①}对于立法而言，在规定权利义务、分配社会资源时，要公平地对待每个人，保障每个人得到他应得的。\par

\daofasym{②}对于司法而言，在解决纠纷、化解矛盾时，要公平地对待当事人，切实维护其合法权益。\par

（3）国家角度（补）：\par

\daofasym{①}要逐步建立以权利公平、机会公平、规则公平为主要内容的社会保障体系，保障公民平等的生存和发展\par

权利。\par

\daofasym{②}实现在法律、制度面前人人平等，让每一位社会成员平等地享有权利，平等地履行义务。\par

\daofasym{③}加紧建设对保障社会公平正义具有重大作用的制度，进一步完善民主权利保障制度及社会保障体系，保\par

证人民当家作主。\par

b2\par


\section*{PDF第63页}
\addcontentsline{toc}{section}{PDF第63页}

\daofasym{①}加快建立覆盖全国城乡的基本公共服务体系，关注弱势群体，解决好人民最关心、最直接、最现实的利\par

益问题，努力使公共服务成果更好惠及广大人民群众，实现学有所教、劳有所得、病有所医、老有所养、\par

住有所居。\par

\textbf{2、为什么要守护正义？}\par

\daofasym{①}追求正义是社会生活的一个重要主题。\par

\daofasym{②}正义感是公民的基本德行。\par

\daofasym{③}期盼正义、实现正义、维护正义，是我们的共同心声。\par

\textbf{※3.如何守护正义?}\par

(1)个人守护正义。守护正义需要勇气和智慧，面对非正义行为，一方面要敢于斗争，相信正义必定战胜\par

那恶；另一方面要讲究策略，寻找有效的方法，做到见义“智”为。\par

[2)司法维护正义。\daofasym{①}司法是捍卫社会公平正义的最后一道防线。司法机关必须坚持以事实为根据，以法\par

建为准绳，确保司法过程和结果合法、公正。\daofasym{②}国家积极推进以司法公正为核心的司法改革，要求司法机\par

关依法独立公正行使司法权，力让人民群众在每一个司法案件中感受到公平正义。\par

\textbf{、为什么要维护社给公平正义？}\par

\daofasym{①}公平正义是人类追求的永恒目标，是法治社会的核\par

心价值。\par

2实现会平正义，是国家、社会和全体公民的共同责任\par

重点知识整合（如何追求或者实现社会公平正义）：注意，本题把公平正义合并一起答题\par

人角度：（面对利益冲突，我们要站在公平的立场，学会担当，以公平之心为人处世。\par

2遇到不公平的行为时，我们要坚守原则立场，敢于对不公平说“不”，采用合理合法的方式和手段，谋\par

求最大限度的公平，努力营造一个公平的环境。\par

\daofasym{②}树立公平意识，同破坏公平的行为作斗争，对受害者伸出援助之手。（1-3可以单独用来回答怎么维护\par

公平）\daofasym{④}守护正义需要勇气和智慧，面对非正义行为，一方面要敢于斗争，相信正义必定战胜邪恶：另一\par

方面要讲究策略，寻找有效的方法，做到见义“智”为。（第四可以单独用来回答怎么维护正义）\par

国家角度：\par

在立法方面，在规定权利义务、分配社会资源时，要公平地对待每个人，保障每个人得到他应得的国\par

家应不断制定、完善一系列法律法规、法律援助制度等司法救济制度，帮助社会弱势群体。\par

\daofasym{②}在司法方面，第一，司法是提卫社会公平正义的最后道防线。司法机关必须坚持以事实为根据，以\par

法律为准绳，确保司法过程和继果会法一公吐。在解决纠纷、化解矛盾时，要公平地对待当事人，切实维\par

护其合法权益。第一国家积极推进司法公为核心的司法改革，要求司法机关依法独立公正行使司法权，\par

努力让人民群众在每一个司法案件中感受到公平正义。\par

实现在法律、制度面前人大平，让每一位社会成员平等地享有权利，平等地履行义务。（比如依法打\par

未各种违法犯罪行为，包括高官犯罪）\par

4加紧建设对保障社会公平正义具有重大作用的制度进一步完善民主权利保障制度、社会保障体系、基\par

本公共服务体系，关注弱势群体，保证人民当家作主，让发展成果更好地惠及人民。\par


\section*{PDF第64页}
\addcontentsline{toc}{section}{PDF第64页}

【漏识补录：公平正义辨析题题干】\par

小明说：“真正的公平和正义是不存在的，我们没有必要去追求公平正义。”请你运用所学知识，对小明的观点进行评析。\par

典型题目：\par

学知识，对小明的观点进行评析。\par

(1)小明的说法是片面的。\par

(2)公平和正义总是受到各种条件的限制，任何社会都存在不公平和非正义的行为。建立\par

一个公平正义的社会，是人类永恒的追求和共同的理想，维护公平正义，是社会主义法治\par

的价值追求\par

(3)实现公平正义，离不开每个人的积极参与和不懈努力。我们要在生活中追求公平、捍\par

卫正义，要采用合理合法的方式，谋求最大限度的公平，努力营造公平的环境，做有正义\par

感的人。\par

拓展：(挑几个例子背诵就可以)\par

为维护社会公平正义，国家采取的举措（八下P113)\par

(1)法规方面：制定、修改一系列法律法规、法律援助制度等司法救济制度，帮助社会弱\par

势群体。\par

(2)收入方面：\daofasym{①}实施最低生活保障制度；\daofasym{②}提高个税起征点。\par

(3)医疗方面：\daofasym{①}跨省异地就医住院医疗费用直接结算；\daofasym{②}编织了世界上最大的全民医疗\par

保障网，建立了基本医疗保障体系；\daofasym{③}加快农村新型医疗制度改革；\daofasym{④}全面推行大病保险\par

政策。\par

(4)养老方面：\daofasym{①}全国大部分省份建立了80周岁以上高龄老人津贴制度；\daofasym{②}关爱老年群体，\par

实施“互联网+养老”行动，推动智慧养老，降低老年人意外风险；\daofasym{③}建立健全定期巡访\par

独居、空巢、留守老年人工作机制。\par

(5)文化教育方面：\daofasym{①}加强新时代乡村教师队伍建设，提高乡村教师待遇；\daofasym{②}城乡免费义\par

务教育全面实现；\daofasym{③}全国所有公共图书馆、文化馆、美术馆实现免费开放。\par

页 5 64\par


\section*{PDF第65页}
\addcontentsline{toc}{section}{PDF第65页}

科技是第 生剂新发质\par

九上 发理上\par

1，论议现允化患设\par

\subsection*{第一课踏上强国之路第一框：坚持改革开放}
\addcontentsline{toc}{subsection}{第一课踏上强国之路第一框：坚持改革开放}

如何认识我国推行的改革？ (What)\par

(1）改革的领域：我国推行的改革是一场全面而深刻的社会变革，包括经济、政治、文化、社会、生态、\par

国防和军队以及党的建设等领域的改革。\par

(2）全面深化改革的总目标完善和发展中国特色社会主义制度，推进国家治理体系和治理能力现代化\par

、为什么要改革开放？(Why)\par

(1)改革开放的意义?(重点)\par

中国共产党团结带领中国人民进行改革开放的伟大革命，极大激发广大人民群众的创造性，极大解放\par

和发展社会生产力，极大增强社会发展的活力，人民生活显著改善，综合国力显著增强，国际地位显著提\par

高，实现了中国人民从此站起来到富起来的伟大飞跃。（三个极大、三个显著）-m天铺一发展\par

2)改革开放的作用/地位？（重点）\par

D改革开放是强国之路。\par

2改革开放解放和发展了生产力。\par

)改革开放是新时代最鲜明的时代特征，它是决定当代中国命运的最关键的一招，也是决定中华民族伟大\par

复兴的关键一招。是决定当代中国命运的关键执择。\par

3.改革开放是怎样促发展的？(How）\par

D我国逐步确立了公有制为主体、多种所有制经济共同发展、按劳分配为主体，多种分配方式并存、社会\par

主义市场经济体制的基本经济制度。\par

\daofasym{②}社会主义基本经济制度既体现了社会主义制度优越性，又同我国社会主义初级阶段社会生产力发展水平\par

相适应。\par

\daofasym{③}改革开放使广大人民群众参与社会劳动、创造社会财富的积极性和主动性空前高涨。尊重劳动、尊重知\par

只、尊重人才、尊重创造已成为社会共识。\par

1.改革开放40年来，中国腾飞主要表现在哪些方面？（中国成就）\par

D极大解放和发展了社会生产力：中国已经成为世界第二大经济体、制造业第一大国，货物贸易第一大国，\par

商品消费第二大国，外资流入第二大国。第二大工业国、外汇储备连续多年位居世界第一。科技、教育、\par

文化等各项事业蓬勃发展。\par

\daofasym{②}中国人民通过改革开放过上了幸福生活,改革开放以来,我国城乡就业规模持续扩大，人民收入较快增长，\par

家庭财产稳步增加,城乡社会保障制度逐步建立和完善,扶贫工作取得举世瞩目的成就。\par

\daofasym{②}改革开放深刻影响着世界，中国成为影响世界的重要力量。中国在全球治理体系中贡献了中国智慧，\par

冕现了大国担当。从“引进来”到“走出去”从加入世界贸易组织到共建“一带一路”从应对亚洲金融危机和\par

国际金融危机到成为世界经济增长的主要稳定器和动力源。\par

\textbf{党的十九大召开有什么重要意义？}\par

党的十九大宣告中国特色社会主义进入新时代。\par

\daofasym{②}为中国人民谋幸福、为中华民族谋复兴，中国共产党不忘初心，牢记使命，带领中国人民走向富强民主\par

文明和谐美丽的社会主义现代化强国之路，\par

第二框：走向共同富裕\par

.进入新时代，我国为什么要全面深化改革？ （重点）\par

\daofasym{①}进入新时代,我国社会主要矛盾已经转化为人民日益增长的美好生活需要和不平衡不充分的发展之间的\par

矛盾。（我国社会主要矛盾）\par

\daofasym{③}我国经济发展还面临区域发展不平衡、城镇化水平不高、城乡发展不平衡不协调等现实挑战。\par

\daofasym{①}改革开放是当代中国最鲜明的特色。改革只有进行时，没有完成时，只有不断弘扬改革创新精神，才能\par


\section*{PDF第66页}
\addcontentsline{toc}{section}{PDF第66页}

(战乡)\par

奏响走向繁荣富强的最强音。\par

\textbf{2.怎样全面深化改革？}\par

\daofasym{①}将改革进行到底，开启全面深化改革新征程。\par

\daofasym{③}促进区域协调发度，坚持中国特色新型城镇化道路，推动城乡发展一体化。\par

\daofasym{③}不断弘扬与时俱进、锐意进取、勤于探索、勇于实践的改革创新精神。\par

\textbf{3.为什么要坚持共享发展成果？}\par

\daofasym{①}衡量一个社会的文明程度,不仅要看经济发展,面且要看发展成果是否惠及全体人民,人民的合法权益是\par

否得到切实保障。 :艺同\par

\daofasym{②}人民对美好生活的向往,就是党的奋斗目标。（选择题经常用到）\par

\daofasym{③}发展的根本且的就是增进民生福证。（选择题出现的频率非常高）\par

\daofasym{④}实现发展成果由人民共享是共产党执政的发展理念之一。\par

\daofasym{③}坚持走共同富裕道路。维护社会公平正义\par

4.如何共享发展成果？（党和政府为了使发展的成果惠及全体人民做出了哪些努力？）\par

\daofasym{①}党和政府抓住人民最关心最直接最现实的利益问题。\par

\daofasym{②}党和政府坚持以大民为中心的发展思想（选择题经常用到），强调人人参与、人人尽力、人人享有，让\par

人民群众共享发展成果。\par

\daofasym{③}提高就业质量和人民收入水平，加强社会保障体系建设，坚决打赢脱贫攻坚战，实施健康中国战略，\par

打造共建共治共享的社会治理格局\par

\daofasym{①}不断满足人民日益增长的美好生活需要，使人民的获得感、幸福感、安全感更加充实、更有保障\par

\textbf{5、近年来，我国不断促进城乡、区域协调发展/保障和改善民生有何重要意义？}\par

(1）有利于缩小贫富差距，实现城乡、区域协调发展，实现共同富裕。\par

(2）（2)有利于提高贫困地区人民的生活水平和获得感，满足人民日益增长的美好生活需要。\par

（3）（3）有利于大民群众坚定中国特色社会主义道路自信、理论自信和制度自信。\par

(4)有利于使改革发展成果更多更公平地惠及奎体大民，促进社会公平正义。\par

\subsection*{第二课创新驱动发展舒功的较量}
\addcontentsline{toc}{subsection}{第二课创新驱动发展舒功的较量}

1、我国目前的科技发展现状是怎样的？（What)\par

（1Y科技创新能力已经成为综合国力竞争的决定性因素\par

（2）成就.在尖端技术的掌握和创新方面打下了坚实基础在一些重要领域走在世界前列。\par

（3）不足：众整体上看,仍然面临科技创新能力不强、科技发展总体水平不高、科技对经济社会的支撑能力\par

不足、科技对经济增长贡献率低于发达国家水平等问题。\par

（4）中国科技创新之路任重道远,需要加快建设创新型国家。是一种活方式\par

2、为什么需要创新？(Why)\par

（1）（个人)创新给我们带来惊喜，让我们获得成就感，改变我们的思维方式和行为方式，激发潜能，超\par

越自我。\par

（2）（社会）技术的创新促进生产力发展，增加社会财富，制度的创新促进公平正义，推动社会进步。\par

创新让生活更美好。我们的生活将因创新变得更加便捷、舒适和丰富多彩。\par

（3）（国家文\daofasym{①}创新驱动是国家命运所系，实施创新驱动发展战略，推动以科技创新为核心的全面创新，\par

让创新成为推动发展的第一动力，是适应和引领我国经济发展新常态的现实需要\par

\daofasym{②}创新是改革开放的生命，改革在创新中提升发展品质。因家用改革之手激活创新引擎，释放更多创新活\par

力。改革创新需要时代精神，改革创新推动中国走向富强。\daofasym{③}创新是引领发展的第一动力（选择题经常经\par

常用到这句话），是民族进步的灵魂，是一个国家兴旺发达的不竭源泉，也是中华民族最鲜明的民族赋。\par

（4）（世界）时代发展呼唤创新（创新已经成为世界主要国家发展战略的重心，在激烈的国际竞争中，\par

唯创新者进,唯创新者强唯创新者胜\par

（5）（人类社会）创新是推动人类社会向前发展的重要力量\par


\section*{PDF第67页}
\addcontentsline{toc}{section}{PDF第67页}

同，推功议尺\par

穿在药动，才4价(个社一\par

3、怎样建设创新型国家? (How)\par

(1）（国家层面分析)\par

\daofasym{①}必须落实科教兴国战略、人才强国战略,将科技和教育摆在经济社会发展的重要位置,把经济建设重心转\par

移到依靠科技进步和提高劳动者素质的轨道上来,加速实现国家的繁荣昌盛。\par

要增强自主创新能方,坚定不移地走中国特色自主创新道路)\par

\daofasym{③}必须加茯形成有利手创新的落理格局和协机制,搭建有利于创新的活动平台和融资平台,营造有利于创\par

新的舆论氛围和法治环境。\par

\daofasym{①}重视素质教育，培养学生创新精神与实践能力，培养创新型人才。\par

\daofasym{③}鼓励支持引导企业自主创新，敢于突破\par

\daofasym{③}鼓励大众创业，万众创新。\par

(2）（社会层面)\par

社会大力宣传创新的重要性，营造大众创业，万众创新的良好社会氛围\par

(3）（企业层面）不显青力车\par

企亚夏社会创新的重要力量，提升创新能力，是企业持续发展之基，市场制胜之道（创新）。大国重器\par

要掌握在自己手里，核心技术不是别大赐孕的，企业不能只是跟着别人走，而必须自强奋斗、敢于突破。\par

企业要制定和完善创新机制对有创新创意的员工进行奖励，激发企业内部良性竞争，形成良好的创新氛\par

围。\par

(4）（青少年层面)\par

D培养创新意识，敢于打破常规，大胆质疑，要有挑战权威的批判精神。\par

\daofasym{②}努力学习科学文化知识，提高科学素养，为创新能力奠定扎实的知识基础。中月天眼攻入伙月\par

\daofasym{③}注重实践，勤动脑多动手，积极参加小发明、小制作等创新活动。\par

0尊重他人的知识产权，学会运用法律武器维护自己的知识产权。\par

、改革与创新的关系？ 我月可拼纵孕是于计博\par

（1）创新是改革开放的生命，改革在不断创新中提升发展品质\par

（2）国家用改革之手激活创新引擎，释放更多创新活力， 组冲中之学\par

我国为什么大力发展教育事业？（教育的重要性） (中月气间站-有状期缩(\par

D一个民族创新能力的提高离不开创新人才的培养。百年大计,教育为本。祠「三 到来\par

\daofasym{②}教育是民族振兴、社会进步的基石（非常非常非常喜欢考选择），是提高国民素质、培养创新型人才、\par

足进人的全面发展的根本途径。\par

\daofasym{③}教育为人们的幸福生活奠基，教育寄托着亿万家庭对美好生活的期盼。\par

\textbf{创新精神有哪些表现？}\par

D表现为敢为人先敢于冒险的勇气和自信。\par

表现为探索新知的好奇心和挑战权威的批判精神。\par

\daofasym{③}表现为承受挫折的坚强意志和沟通合作的团队精神\par

\daofasym{①}表现为舍我其谁的责任担当和造福人类的济世情怀。\par

我国为什么要鼓励万众创新？【注意：创业与创新两个词不同】\par

\daofasym{①}每个人都是创新者,都向往在创新中实现自我价值;每个人都是创业者,都可以通过辛勤劳动为国家和人\par

民作出贡献。\par

\daofasym{②}企业是社会创新的重要力量。\par

\daofasym{③}时代需要弘扬创新精神。\par

\daofasym{①}创新的目的是增进人类福社（选择题），让生活更美好。创新让我们获得更多的尊重和认可,让我们过上\par

体面而有尊严的生活。\par

生意：\daofasym{①}国家创新的三大战略：创新驱动发展战略+科教兴国战略+人才强国战略\par

\daofasym{②}国家创新的道路：中国特色自主创新道路\par

\daofasym{③}答题一定要学会分清主体：国家、个人、企业、社会\par


\section*{PDF第68页}
\addcontentsline{toc}{section}{PDF第68页}

【结构整理：社会主义民主的属性】\par

产生：新中国成立后，随着各级人民民主政权的建立和社会主义改造的完成，社会主义民主在中国大地上得以真正确立；它从中国社会土壤中长出来，在实践中不断得到验证，是有生命力的。\par

本质：人民当家作主。\par

特点：我国社会主义民主是维护人民根本利益的最广泛、最真实、最管用的民主。\par

作用：有助于推动经济社会持续健康发展，实现人民安居乐业、社会和谐稳定、国家繁荣富强。\par

真谛：有事好商量，众人的事情由众人商量。\par

形式：选举民主和协商民主；协商民主是我国社会主义民主政治的特有形式和独特优势。\par

目的：保障最广大人民的利益。\par

制度保障：根本政治制度是人民代表大会制度；基本政治制度包括中国共产党领导的多党合作和政治协商制度、民族区域自治制度、基层群众自治制度。\par

【结构整理：公民参与民主生活的途径】\par

民主选举：人民实现民主权利的一种重要形式；形式包括直接选举和间接选举、等额选举和差额选举；要遵循公开、公平和公正的原则，公民要积极、主动、理性地参与。\par

民主决策：形式包括社情民意反映制度、专家咨询制度、重大事项社会公示制度、社会听证制度；意义在于保障人民利益充分实现，促进决策科学化、民主化。\par

民主监督：是公民参与民主生活、行使公民监督权的具体体现；有利于国家机关和国家工作人员改进工作、防止滥用权力、预防腐败，也有助于增强公民参与意识。\par

度。该\par

中间到)分 尾国力天\par

(同多关) 第三课追求民主价值\par

1.民主价值的实现方式：民主价值的实现要靠民主形式和民主制度的建立。\par

2. 新型的民主一—社会主义民主的属性? (重点）\par

新中国成立后，随着各级人民民主政权的建立和社会主义改造的完成，社会主义民主在中国大地上得\par

以真正确立。\par

产生它是从中国的社会土壤中长出来的，在实践中不断得到验证，是有生命力的。\par

本质人民当家作主。\par

特点我国社会主义民主是维护人民根本利益的最广泛人最真实最管用的民主。\par

作用有助于推动经济社会持续健康发展，实现人民安居乐业、社会和谐稳定、国家\par

繁荣富强。\par

真谛有事好商量，众人的事情由众人商量。\par

选举民主和协商民主\par

形式\par

（协商民主是我国社会主义民主政治的特有形式和独特优势）\par

目的保障最广大人民的利益。\par

制度根本政治制度：人民代表大会制度。基本政治制度：中国共产党领导的多党合\par

保障作和政治协商制度；民族区域自治制度；基层群众自治制度。\par

3.公民参与民主生活的途径（三个）=民主选举+民主决策+民主监督 （重点）\par

注意：如果考到民主管理，一定要涉及到基层管理，居委，村委这一些最管用民主\par

（1）地位：民主选举是人民实现民主权利的一种重要形式\par

民主\par

（2）形式：直接选举和间接选举、等额选举和差额选举\par

选举\par

（3）要求：\daofasym{①}不管采用哪种形式，民主选举都要遵循公开、公平和公正的原则\par

（选择）\par

\daofasym{②}公民要积极、主动、理性地参与民主选举。\par

（1）形式：社情民意反映制度、专家咨询制度、重大事项社会公示制度、社会\par

听证制度。\par

（2）意义：\daofasym{①}(个人)民主决策是保障人民利益得到充分实现的有效方式，有助\par

民主\par

于增强公民的参与感、责任感和主人翁意识，民主意识。有利于保障公民的合\par

决策\par

法权益，有利于维护人民的根本利益\daofasym{②}（对国家）有利于决策方认真听取各方\par

意见，集中民智，促进决策的科学化和民主化。\daofasym{③}有利于社会形成民主参与政治\par

生活的氛围，维护良好的社会秩序（对社会）\par

（1）地位：民主监督是公民参与民主生活、行使公民监督权的具体体现。\par

（2）意义：\daofasym{①}（国家机关）实行民主监督，有利于国家机关和国家工作人员改\par

民主\par

进工作，提高工作效率，克服官僚主义，防止滥用权力，预防腐败。\daofasym{②}（个人）\par

监督\par

实行民主监督，有助于增强公民的参与意识，激发公民的参与热情。有利于保\par

障公民的合法权益，有利于维护人民的根本利益\par

收为人员务，使[民满专服务n5b 公\par

5、\par

4.公民参与民主生活需要具备哪些能力（或者如何参与民主生活）\par

具有社会责任感和主人翁意识；\par

\daofasym{②}要以理性、公正、客观的态度全面、深刻、辩证地看问题；\par

\daofasym{③}立场正确、逻辑清晰地表达观点和意见\par

\daofasym{④}逐步提高依法有序参与民主生活的能力。\par

\textbf{5.为什么要增强民主意识？}\par

\daofasym{①}在现代社会，民主应该成为公民的一种生活方式和生活态度。一个国家和社会民主生活的质量和水平，\par

与公民的民主意识密切相关。（必要性）\par

\daofasym{②}一个国家拥有什么素养的公民，就会有什么样的未来。增强我国公民的民主意识，有利于完善中国特色\par


\section*{PDF第69页}
\addcontentsline{toc}{section}{PDF第69页}

【漏识补录：第5题增强民主意识续句】\par

增强民主意识是社会主义民主的要求，也是社会主义制度永葆生命力的重要保证。\par

（重要性）\par

6.如何增强民主意识？ （重点）\par

\daofasym{①}公民要自觉遵守宪法，始终按照宪法原则和精神参与民主生活。\par

\daofasym{②}公民要不断积累民主知识，形成尊重、宽容、批判和协商的民主态度。\par

\daofasym{③}公民要通过参与公共事务和民主生活，在实践中逐步增强民主意识。\par

\daofasym{①}不断增强民主意识，使民主思想和法治精神成为公民的自觉信仰，体现在日常言行中。\par

\subsection*{第四课建设法治中国 率质果束:全间依法治}
\addcontentsline{toc}{subsection}{第四课建设法治中国 率质果束:全间依法治}

1.法治要求是什么？(What）（重要)\par

\daofasym{①}法治要求实行良法之治。（何为良法？）良法应当反映最广大人民群众的意志和利益，反映社会的发展规律，维护公\par

民的基本权利，符合公平正义的要求，促进人与社会的共同发展。（科学立法)\par

\daofasym{②}有法律制度不等于就有法治，法治还要求实行善治。（何为善治？）法治建立在民主的基础上，通过赋予公民更多的\par

参与公共活动的机会和权利，实现公共利益的最大化。（严格执法+公正司法+全民守法）\par

2（1）为什么选择法治道路？（Why）（或选择法治道路意义、作用）（重要）\par

の（个人）法治能够为人们提供良好的生活秩序，让人们能够建立起一个基本、稳定、持续的生活预期，保障人们在社\par

会各个领域依法享有广泛权利和自由，使人们安全、有尊严地生活。\par

\daofasym{②}（社会）法治是现代政治文明的核心，是发展市场经济、实现强国富民的基本保障，是解决社会矛盾、维护社会稳定、\par

实现社会正义的有效方式。\par

\daofasym{③}（国家）走法治道路是实现中华民族伟大复兴的必然选择）法治是人们生活的共同准则。是国家治理现代化的重要标\par

志，助推中国梦的实现，是实现政治清明、社会公平、民心稳定、国家长治久安的必由之路。\par

（人类社会）法治是人类社会进入现代文明的重要标志。\par

(2）全面依法治国的重要性？（Why）（重要）\par

D全面依法治国的地位（重要性）：全面依法治国是中国特色社会主义的本质要求和重要保障。\par

\daofasym{②}全面推进依法治国总目标：建设中国特色社会主义法治体系，建设社会主义法治国家。\par

\daofasym{③}全面依法治国要求：实现科学立法、严格执法、公正司法、全民守法（十六字方针），促进国家治理体系和治理能力\par

现代化。\par

（怎样厉行法治/怎样建设法治中国？（How）（重要）\par

D国家：要努力使每一项立法都得到人民群众的普遍拥护（科学立法），使每一部法律法规都得到严格执行（严格执法），\par

使每一个司法案件都体现公平正义（公正司法），使每一位公民都成为法治的忠实崇尚者、自觉遵守者和坚定捍卫者（全\par

民守法）；坚定不移地走中国特色社会主义法治道路，必须坚持党的领导、人民当家作主、依法治国有机统一。\par

公民：要增强尊法学法守法用法意识，弘扬法治精神，强化规则意识，树立正确的权利义务观念。\par

3政府：政府及其工作人员带头尊法学法守法用法，提高运用法治思维和法治方式深化改革、推动发展、化解矛盾、维\par

户稳定、应对风险的能力。\par

D社会：要加强法治宣传，弘扬法治精神，共同营造良好的法治文化环境，在全社会鲜明地树立起“守法光荣、违法可\par

让”的法治文化导向。\par

国家和社会治理需要法律和道德共同发挥作用，既重视发挥法律的规范作用，又重视发挥道德的教化作用。\par

坚定不移地走中国特色社会主义法治道路，必须坚持党的领导、人民当家作主、依法治国有机统一。\par

\textbf{政府作用、宗旨、核心、准则？}\par

D一方面，人们的社会生活需要政府管理；另一方面，人们又享受着政府提供的公共服务。\par

\daofasym{②}政府的宗旨是为人民服务，政府的工作要对人民负责，为人民谋利益。\par

依法行政的核心是规范政府的行政权。\par

D依法行政是现代法治政府行使权力普遍奉行的基本准则。\par

，如何实现依法行政？（怎么建设法治政府）（重要）\par

府角度：\daofasym{①}政府及其工作人员应依据宪法和法律的规定正确行使权力。防范行政权力的滥用，维护广大人民群众的合\par

去权益，提高政府公信力，从而推进民主法治建设进程。\par

\daofasym{②}全面推进政务公开，保障公民的知情权、参与权、表达权和监督权，促进政府决策的科学化和民主化。\par

公民角度：\daofasym{③}公民也要积极参与，献计献策，主动监督。\par

道德与法治关系是什么？（重要）\par

高度注意，此题容易出现辨析题，答题思路：\daofasym{①}先写一下道德或者法律的作用。结合材料分析\daofasym{②}写一下法律或者道德\par


\section*{PDF第70页}
\addcontentsline{toc}{section}{PDF第70页}

任件，泻比义，\par

的作用，结合材料分析\daofasym{③}最后写二者的关系】\par

\daofasym{①}国家和社会治理需要法律和德共同发挥作用，既重视发挥法律的规范作用，又重视发挥德的教化作用\par

\daofasym{②}以法治承找道德理念，强化法律对道德建设的促进作用：以道德滋养法治精神，强化道德对法治文化的支撑作用。\par

\daofasym{③}法律与道德相辅相成，法治与德治相得益彰。\par

\subsection*{第五课守望精神家园第一框延续文化血脉}
\addcontentsline{toc}{subsection}{第五课守望精神家园第一框延续文化血脉}

1. 中华文化的特点、组成、作用(重要)\par

（1）特点：源远流长，博大精深，薪火相传，历久弥新。\par

具有应对挑战、与时俱进的创造力和海纳百川、有容乃大的包容力。\par

（2）组成：独具特色的语言文字，浩如烟海的文化典籍，名扬世界的科技工艺，异彩纷呈的文学艺术组\par

成了源远流长、博大精深的中华文化。\par

作用：\daofasym{①}文化是一个国家、一个民族的灵魂。文化兴国运兴，文化强民族强。\par

\daofasym{②}中华文化积淀着中华民族最深层的精神追求，代表着中国民族独特的精神标识【单选】，为中华民族的\par

伟大复兴提供精神动力。\par

\daofasym{③}中华文化独一无二的理念、智慧、气度、神韵，增添了中国人民和中华民族内心深处的自信和自豪。\par

2.中国特色社会主义文化的构成、来源\par

构成：（1）中华优秀传统文化：（2）革命文化；（3）社会主义先进文化。\par

来源：中国特色社会主义文化源自于中华优秀传统文化，熔铸于革命文化和社会主义先进文化，植根于中\par

国特色社会主义的伟大实践，积淀着中华民族最深层的精神追求。【选题）\par

3.文化自信的含义、原因、措施\par

（1）含义：自信是对自身文化价值的充分肯定，是对自身文化生命力的坚定信念，是一个国家、一个民\par

族发展中更基本、更深沉、更持久的力量。\par

（2）原因：\par

\daofasym{①}没有高度的文化自信，没有文化的繁荣兴盛，就没有中华民族的伟大复兴。\par

\daofasym{②}坚定文化自信，事关国运兴衰、文化安全和民族精神的传承发展。\par

（3）（如何坚定文化自信？=如何发展中国特色社会主义文化？）\par

要坚持以马克思主义为指导，推动中华优秀传统文化创造性转化、创新性发展，继承革命文化，发展社会\par

主义先进文化，不忘本来，吸收外来，面向未来，不断铸就中华文化新辉煌。\par

4.中华传统美德内容、原因、措施？（重要）\par

（1）内容：中华传统美德内涵丰富，博大精深。其中，有忧国忧民、道济天下的爱国情怀，有勤劳勇敢、\par

自强不息的奋进品格，有自尊互敬、助人为乐的和乐风范，有诚信守法、见利思义的高尚情操，有孝敬父\par

母、尊敬师长的道德规范，有律已宽人、扬善抑恶的处世准则，等等。\par

（2）中华传统美德的重要性（作用及影响/弘扬传统美德的原因)\par

\daofasym{①}中华传统美德是中华文化的精髓，蕴含着丰富的道德资源，是建设富强民主文明和谐美丽的社会主义现\par

代化强国的精神力量。\par

\daofasym{②}经过长期的历史积淀，中华传统美德已经融入中华民族的思维方式、价值观念，行为方式和风俗习惯，\par

成为一种文化基因。\par

\textbf{（3）青少年如何践行中华传统美德？}\par

\daofasym{①}推进社会公德、职业道德、家庭美德、个人品德建设，青少年责无旁贷。\par

\daofasym{②}倡导向上向善、孝老爱亲、忠于祖国、忠于人民，青少年必须身体力行\par

第二框凝聚价值追求\par

1.中华民族精神的内涵、品格、本质、原因、措施 （重要）\par

（1）内涵：中华民族精神是指以爱国主义为核心的团结统一爱好和平)勤劳勇敢、自强不总的伟大民\par

族精神。\par

品格：中华民族精神具有与时俱进的品格。它在不同的历史时期有着不同的表现，并随着时代进步不断丰\par

富和发展。\par

本质：当代中国，爱国主义的本质就是坚持爱国和爱党，爱社会主义的高度统\par

（2）为什么要传承和弘扬中华民族精神（重要性、价值）？ (重要)\par


\section*{PDF第71页}
\addcontentsline{toc}{section}{PDF第71页}

理；绿色，科绽，幼请是积山\par

0一个民族要生存和发展，就要有昂扬向上的民族精神。否则，就会失去凝聚力和生命力，就难以屹立于\par

世界民族之林。\par

\daofasym{②}伟大民族精神始终是中华民族生生不息、发展壮大的强大精神支柱，是维系我国各族人民世世代代团结\par

套斗的牢固精神纽带，是激励中华儿女为实现中国梦而奋斗的不竭精神动力。\par

(3）传承和弘扬民族精神的具体表现是什么？（措施） （重要）【可以用来回答如何弘扬民族精神\par

或者其他精神】\par

の国家：表现在国家危难、民族危亡的紧要关头能够挺身面出、舍生忘死、前仆后继。\par

\daofasym{②}他人：表现在他人生命、财产遇到危险的关键时刻能够见义勇为、扶危济困、无私奉献。\par

\daofasym{③}个人：表现在日常学习工作中的勤勤愿愿、任劳任怨、创业创优，要从自己做起、从现在做起、从小事\par

做起，自觉高扬民族精神\par

\textbf{2 社会主义核心价值观是怎样形成的、内容、重要性、措施？}\par

(1)怎样形成:\par

D中国独特的文化传统、独特的历史命运、独特的基本国情，注定我们必然坚守植根于中华文化沃土又具\par

有当代中国特色的价值观。\par

\daofasym{②}社会主义核心价值观不仅是中国人民在共同生活中形成的价值共识，而且吸收了世界文明的有益成果，\par

是当代中国精神的集中表现。\par

(2）内容: （重要）【选择】\par

富强、民主、文明、和谐是国家层面的价值目标（建设什么样的国家），自由、平等、公正、法治是社会\par

面的价值取向（建设什么样的社会），爱国、敬业、诚信、友善是公民个人层面（建设什么样的公民）\par

的价值准则。\par

(3）倡导社会主义核心价值观的重要性\par

0价值观是文化最深层的内核。\par

社会主义核心价值观是当代中国人评判是非曲直的价值标准，是当代中国精神的集中体现。【选择】\par

社会主义核心价值观凝结着全体人民共同的价值追求，是坚持和发展中国特色社会主义的价值导向，又\par

实现中华民族伟大复兴的价值引领，\par

)社会主文核心价值观促进人的全面发展，引领社会全面进步。\par

\textbf{【3）青少年如何培育和践行社会主义核心价值观？}\par

)培育和践行社会主义核心价值观，要与日常生活紧密联系起来，做到落细、落小、落实。\par

自觉做到勤于学习、敏于思考，注重修养、勇于实践，明辨是非、善于选择，认真做事、踏实做人。\par

充题目：\par

你从民族精神的角度，谈谈“××精神”的时代价值。【如何弘扬XX精神】\par

中华民族精神在不同的历史时期有着不同的表现，“××精神”是中华民族精神在新的时代条件下的集中体现\par

它是中华民族生生不息发展壮大的强大精神支柱，是团结奋斗的精神纽带，是实现中国梦的不竭精神动力：它有利\par

激发爱国主义热情，构筑中国精神、中国价值、中国力量。\par

五课几个关键词：中华文化、中华民族精神、文化自信、社会主义核心价值观\par

华传统美德、中国特色社会主义文化\par

抗瘦弱钾。生甲,工,孕因间小食生乱身生科学,年泌与其\par

5敏以功,劳模。遍\par


\section*{PDF第72页}
\addcontentsline{toc}{section}{PDF第72页}

不同时期的民族精神：\par

\daofasym{①}革命时期：井冈山精神、长征精神、延安精神、西柏坡精神\par

\daofasym{②}社会主义革命和建设时期：抗美援朝精神，雷锋精神，两弹一星精神\par

\daofasym{③}改革开放事情的精神：航天精神、抗疫精神、脱贫攻坚精神、奥运精神等\par

补充题目：\par

1、如何传承和弘扬中华优秀传统文化（多主体：国家十个人+社会）\par

（1）国家：\daofasym{①}制定并完善保护传统文化的法律法规；\daofasym{②}加大保护传统文化的资金投入；\daofasym{③}开发与保护相\par

结合，大力发展传统文化旅游产业；\daofasym{④}开展中华优秀传统文化进校园活动，举办经典诵读、国学讲堂活动；\par

\daofasym{③}开展全民阅读活动；\daofasym{③}宣传和提升节日的文化内涵，宣传保护传统文化的重要性；培养和发掘有传统\par

技艺的人才；\daofasym{③}坚持创新，探索保护传统文化的新途径；\daofasym{③}增强与世界文明交流、对话的意识，不断丰富\par

和发展中华文化，推动中华文化走向世\par

(2)社会：\daofasym{①}加强宣传教育，增强人们的文化遗产保护意识；\daofasym{②}鼓励继承、发扬民间艺术。\par

(3)个人：\daofasym{①}增强民族自尊心、自信心、自豪感，积极主动学习优秀传统文化；\daofasym{②}保护本民族文化，珍爱\par

自己的精神家园；\daofasym{③}做优秀传统文化的传播者、弘扬者、建设者；\daofasym{④}宣传优秀传统文化，让世界了解源远\par

流长、博大精深的中华文化；\daofasym{③}取其精华，弃其糟，紧跟时代发展要求，不断为之增添新的富有生命力\par

的内容。\par

2、有人认为，有些文化遗产在现代生活中没有实用价值，没有必要保护。请你运用文化相关知识对这一\par

观点进行评析。(九上P59)(让国宝回家，修复文物，保护文化遗产，申遗等的作用）\par

(1)此观点错误。\par

(2)文化遗产是民族智慧的结晶、民族精神的象征，是民族生命力和创造力的重要体现，也是人类文明的\par

瑰宝。\par

(3)保护文化遗产，有利于增强中华民族的文化认同感，有助于中华民族守望精神家园，铸魂强魄，凝心\par

聚力。\par

(4)因此，我们应该积极保护文化遗产，让文化遗产通过现代化手段发挥其重要价值。\par

3、中国人民具有的伟大精神\par

(1)中国人民是具有伟大创造精神的人民表现为敢为人先。\par

(2)中国人民是具有伟大奋斗精神的人民）表现为自强不息。\par

(3)中国人民是具有伟大团窈精神的人民，表现为同舟共济。\par

(4)中国人民是具有伟大梦想精神的人民，表现为志存高远。\par

\textbf{4、如何传承和弘扬民族精神？}\par

(1)国家：\daofasym{①}立足于发展中国特色社会主义的伟大实践；\daofasym{②}结合时代和社会发展要求，不断为之增添新的\par

富有生命力的内容；\daofasym{③}加强法治建设，为传承和弘扬民族精神提供法治保障；等。\par

(2)学校：举办形式多样的民族精神宣传活动和主题教育活动。如：\daofasym{①}组织学生阅读优秀历史文化读物，\par

感悟民族精神；\daofasym{②}组织学生参加社会实践活动，感悟民族精神；\daofasym{③}通过各种形式、途径加强民族精神的宣\par

传教育；等。\par


\section*{PDF第73页}
\addcontentsline{toc}{section}{PDF第73页}

3请少年：\daofasym{①}从自己做起、从现在做起、从小事做起，自觉高扬民族精神；\daofasym{②}继承和发扬中华民族的优\par

传统，培育和践行社会主义核心价值观；\daofasym{③}坚定理想信念，努力学习，自觉肩负起实现中华民族伟大复\par

的历史使命；\daofasym{④}自觉接受民族精神教育，培养爱国情怀和奉献精神；\daofasym{③}勇于面对并克服学习、生活中的\par

难，艰苦奋斗，自强不息；\daofasym{③}学习和宣传英雄模范人物的优秀品质，做新时代的好少年。\par

弘扬英雄人物先进事迹有什么意义！[2019.31(2]\par

有利于在全社会倡导和弘扬英雄精神，形成崇敬英雄、学习英雄的良好社会氛围。\par

)有利于传承中华传统美德，培育社会主义核心价值观，传播社会正能量。\par

)有利于引导人们结合国家和社会的需要，在平凡的岗位中创造和实现生命的价值。\par

)有利于引导青少年树立正确的世界观、人生观、价值观。\par

\subsection*{第六课建设美丽中国第一框正视发展挑战}
\addcontentsline{toc}{subsection}{第六课建设美丽中国第一框正视发展挑战}

几句重要的话：\daofasym{①}走绿色发展道路，建设生态文明，建设美丽中国\par

\daofasym{②}3个基本国策：针对人口问题一一计划生育基本国策；针对环境问题一一保护环境基本国策；针对资源问题-\par

节约资源基本国策\par

\daofasym{③}坚持创新、协调、绿色、开放、共享的发展理念，建设资源节约型、环境友好型社会。\par

\daofasym{④}生态文明与脱贫攻坚、美丽中国、转变经济发展方式（经济由高速发展转向高质量发展）、坚持走中国特色城\par

镇化道路都有关系。\par

\textbf{我们为什么要重视发展中的人口问题？}\par

人口问题已经成为一个日益严峻的全球性问题,成为人类社会面临的重大挑战之一。\par

人口问题加重了资源和环境的压力,也带来了严重的社会问题。\par

我国人口现状的特点有哪些？(重要)\par

口国情：我国是世界上人口最多的国家\par

本特点：人口基数大，人口素质低\par

的特点（总人口增速趋缓\daofasym{②}总和生育率明显低于更替水平出生人口男女性别比偏高老龄化加\par

\daofasym{③}大量的大口流动。\par

我国为什么要长期坚持计划生育基本国策？（实行计划生育的必要性） (重要)\par

人口问题始终是我国面临的全局性、长期性、战略性问题。在未来相当长的时期内我国人口众多的基\par

国情不会根本改变，人口对经济社会发展的压力不会根本改变，人口与资源环境的紧张关系不会根本改变。\par

国家推行计划生育，使人口的增长同经济和社会发展计划相适应，有效地缓解了人口对资源，环境的压力，\par

力地促进了经济发展与社会进步。（计划生育的作用）\par

\textbf{我国为什么更新生育政策？}\par

计划生育要随着人口和经济社会发展形势的变化不断完善。\par

我国实施“全面两孩”政策，就是为了更好地适应经济发展的需要，促进家庭幸福和社会和谐\par

人口长期均衡发展。\par

我国的资源现状、形成原因、危害分别是什么？ (重要)\par

状：我国自然资源丰富，总量大，种类多，但人均资源占有量少,开发难度大，总体上资源紧缺。\par

成原因：长期以来,我国资源开发利用不尽合理、不够科学,依靠消耗大量资源换取经济发展的现象突出，\par

此造成的浪费、损失、污染和破坏都很严重。\par

害：必然会导致资源的枯竭和对生态环境的破坏，严重影响经济的可持续发展，经济发展的空间和后劲\par


\section*{PDF第74页}
\addcontentsline{toc}{section}{PDF第74页}

【漏识补录：第6题关爱和保护环境第\daofasym{①}点】\par

\daofasym{①}环境恶化加剧自然灾害的发生，严重破坏生态平衡，威胁着人民的生命安全和身体健康。\par

会越来越小。\par

\textbf{6.为什么要关爱和保护环境？}\par

\daofasym{②}生态环境没有替代品,用之不觉,失之难存。\par

\daofasym{③}人类关爱和保护环境就是走向重生，漠视和破坏环境就是走向自我毁灭。\par

7.面对严峻的人口、资源、环境问题，我们如何应对/出路/必然选择？ (重要)\par

通过转变发展方式才能得到解决，坚持绿色发展，走生产发展、生活富裕、生态良好的文明发展道路，\par

是我们的必然选择。\par

8.为什么要保护环境，节约资源？（如何构建资源节约型和环境友好型社会？为什么要坚持可持续发展\par

道路？）\par

个人：有利于提高人们的环保和节能意识，人与自然和谐共生。\par

社会：有利于我国社会的可持续发展，建设环境友好型社会和资源节约型社会：\par

有利于践行绿色发展理念，走绿色发展道路；建设生态文明，实现可持续发展：有利于建设和谐社\par

会\par

国家：有利于缓解我国严峻的环境、资源形势，提高资源利用率。\par

有利于推动整个社会走生产发展、生活富裕、生态良好的文明发展道路。\par

有利于贯彻落实保护环境、节约资源的基本国策 1民民业\par

9．我国的基本国策：对外开放、计划生育、保护环境、节约资源。\par

10.我国的发展战略：可持续发展战略、科教兴国战略、人才强国战略、创新驱动发展战略、西部木开发\par

4战略、健康中国战略、乡村振兴战略。牡吉习拍绿发限，一传科生百-环境]\par

人质物以生而m相质慧础\par

第二框：共筑生命家园\par

\textbf{1.为什么要坚持人与自然和谐共生？}\par

\daofasym{①}追求人与自然和谐共生，是人类面对生态危机作出的智慧选择 h\par

\daofasym{②}自然为人类的生存与发展提供滋养和必要条件；人类作为自然的一部翁他有责任避免自然受到不必要\par

的伤害,同时要为开发和利用自然作出必要的补偿和修复。\par

\daofasym{③}人类开发和利用自然,但不能肆意凌驾于自然之上,必须遵循自然规律。否则一味地索取，必然受到它的\par

的惩罚。\par

\textbf{2.怎样建设生态文明？（重要）=怎样走绿色发展道路，建设美丽中国？}\par

（1）国家：\par

\daofasym{①}以资源环境承载能力为基础,以自然规律为准则,以可持续发展、人与自然和谐共生为目标。\par

\daofasym{②}坚持节约资源和保护环境的基本国策，建设资源节约型、环境友好型社会。\par

\daofasym{③}坚持创新、协调、绿色、开放、共享的发展理念。\par

\daofasym{④}正确处理好经济发展与生态环境保护的关系，坚持绿色富国、绿色惠民。\par

\daofasym{③}大力倡导节能、环保、低碳、文明的绿色生产生活方式，让绿色发展理念渗透到人们日常生活细节中，\par

成为每个社会成员的自觉行动。\par

\daofasym{③}建设生态文明必须严守资源消耗上限、环境质量底线、生态保护红线。\par

转变经济发展方式，坚持绿色发展，走生产发展、生活富裕、生态良好的文明发展道路。\par

（2）企业：【补充】\par

\daofasym{①}树立科学发展观，增强社会责任感，加快科技创新，提高节能环保技术；\par

\daofasym{②}发展循环经济，降低能耗，提高资源利用率；\par

\daofasym{③}推行绿色生产，减少污染物排放，保护环境。\par

（3）个人：青少年在生活中如何建设美丽中国？ (重要)\par

\daofasym{①}学习环保知识，增强保护环境和节约资源的意识。\par

\daofasym{②}积极宜传生态环境知识，倡导低碳绿色生活方式。\par


\section*{PDF第75页}
\addcontentsline{toc}{section}{PDF第75页}

【结构整理：民族相关概念表】\par

民族分布特点：大散居、小聚居、交错杂居。\par

处理民族关系的基本原则：民族平等、民族团结和各民族共同繁荣。\par

基本政治制度：民族区域自治制度。\par

社会主义新型民族关系：平等、团结、互助、和谐。\par

依法保护环境，善于同破坏环境的行为作斗争。\par

正确行使公民的监督权，为改善空气质量献计献策。\par

从身边的点滴小事做起，培养保护环境、节约资源的行为习惯。例如：节约用水用电；做好垃圾分类；\par

绿色出行，少坐私家车等环保行动。\par

如何正确处理经济发展与生态环境保护的关系？（如何处理好绿水青山就是金金山银山的关系）\par

保护生态环境就是保护生产力，改善生态环境就是发展生产力，决不能以牺牲环境、浪费资源为代价，\par

取一时的经济增长。\par

我们既要绿水青山，也要金金山银山。如果金金山银山要以失去绿水青山为代价，那么宁要绿水青山，不要\par

金山银山。绿水青山就是金金山银山。\par

\subsection*{第七课 中华一家亲第一框促进民族团结}
\addcontentsline{toc}{subsection}{第七课 中华一家亲第一框促进民族团结}

民族相关概念\par

概念 解释\par

族分布特点 大散居、小聚居、交错杂居。\par

理民族关系的基本原则 民族平等、民族团结和各民族共同繁荣\par

族基本政治制度 坚持民族区域自治制度。\par

型民族关系 平等、团结、互助、和谐。\par

\textbf{为什么要加强和巩固民族团结？}\par

加强和巩固民族团结维护祖国统一，是中华民族的最高利益，（选择题）\par

我国各民族在数千年的交往中孕育了团结友爱的宝贵传统。\par

我国各民族始终同呼吸、共命运、心连心，克服种种困难和艰险，顶住压力，直面挑战,追求共同发展、共同富裕、共\par

繁荣。\par

维护和促进民族团结，是每个公民的神圣职责和光荣义务。\par

铸牢民族共同体意识，手足相亲、守望相助、齐心奋斗，伟大的祖国才能繁荣发展。\par

\textbf{为什么要加快民族地区经济社会文化发展？}\par

加快民族地区经济社会文化发展，逐步缩小发展差距，促进民族地区共同繁荣，是增进民族团结、发展社会主义民族关\par

的必由之路。\par

\textbf{党和政府为促进民族地区共同繁荣的措施？}\par

经济社会方面：加大对欠发达民族地区的扶贫支持力度，推动西部大开发战略，实施兴边富民行动，通过输入技术、\par

理、人才等方式增强民族地区自我发展能力\par

民生方面：支持民族地区发展教育，实施积极的就业政策，初步建立基本医疗保障制度。\par

文化方面：大力扶持少数民族文化的保护、继承、创新和发展工作，积极促进各民族之间的文化交流，使少数民族\par

文化获得前所未有的发展。\par

政治方面：坚持民族区域自治制度，坚持处理民族关系的基本准则，坚持推进新型民族关系。\par

\textbf{青少年应如何维护民族团结？}\par

尊重各民族的语言文字、风俗习惯、宗教信仰；\par

拥护和宣传党和国家的民族政策；\par

关心、爱护少数民族同学；多说有利于民族团结的话，多做有利于民族团结的事。\par

与破坏民族团结的言行作斗争。\par

补充题目：党和政府重视西藏地区发展有何意义\par

\daofasym{①}有利于完善民族区域自治制度和巩固社会主义新型民族关系，维护民族团结，坚持民族平等民族\par

团结和各民族共同繁荣。\daofasym{②}有利于促进区域协调发展，缩小发展差距。\daofasym{③}有利于巩固国防维护边疆\par

稳定。\daofasym{④}有利于提高西藏人民生活水平，实现中华民族伟大复兴\par

第二框维护祖国统一\par

注意几个重点：\par


\section*{PDF第76页}
\addcontentsline{toc}{section}{PDF第76页}

\daofasym{①}香港澳门问题-----”一国两制"\daofasym{②}台湾问题-页——”—国两制，和平统一\par

\daofasym{③}两岸关系的政治基础 一个中国原则\par

\textbf{1. 怎样反对分裂?}\par

\daofasym{①}反对分裂，就要维护国家统一、国家主权和领土完整。\par

\daofasym{②}反对分裂，就要反对一切形式的民族分裂活动，尤其要坚决反对借民族和宗教之名搞暴力恐怖活动。\par

\daofasym{③}反对分裂，就要维护国家安全。\par

2.为什么要维护国家统一、反对分裂？（反对分裂的原因）\par

\daofasym{①}分裂会导致社会动荡，经济发展停滞不前，各族人民就会遭殃。\par

\daofasym{②}任何企图搞民族分裂的人都是历史罪人，一切破坏民族团结、制造民族分裂的行为都将受到法律的制裁。\par

\daofasym{③}维护国家统一，反对分裂是爱国主义精神的具体体现，是每个公民义不容辞的责任。\par

\textbf{3.“一国两制”的内容是什么？}\par

“一个国家、两种制度”，简称为“一国两制”。在祖国统一的前提下，国家的主体坚持社会主义制度，同时在台\par

湾、香港、澳门保持原有的资本主义制度和生活方式长期不变，享有高度的自治权。“一国”是实行“两制”的前提和\par

基础，“两制”从属和派生于“一国”并统一于“一国”之内。\par

\textbf{4、为什么要坚持一国两制?}\par

\daofasym{①}香港、澳门回归祖国以来，“一国两制”实践取得了举世公认的成功。\par

\daofasym{②}香港、澳门回归以来，与祖国内地优势系补，共同发展，共担民族复兴的历史责任，共享祖国繁荣富强的伟大荣光。\par

5、怎样保持香港、澳门长期繁荣稳定？多元出 用面已经形战\par

（1）必须全面准确贯彻“一国两制之“港人治港”“澳人治澳”、高度自治的方针\par

（2）坚持依法治港治澳，维护宪法和基本法确定的特别行政区宪制秩序。\par

（3）落实中央对特别行政区全面管治权，落实特别行政区维护国家安全的法律制度和执行机制，维护国家主权、安全、\par

发展利益和特别行政区社会大局稳定\par

（4）坚决防范和遏制外部势力干预港澳事务，支持港澳巩固提升竞争优势，更好融入国家发展大局。\par

【4+5用来回答香港澳门问题】\par

6.为什么祖国完全统一必定会实现2\par

（1）解决台湾问题，实现祖围完全统一）是全体中华儿女的共同愿望，是中华民族的根本利益所在\par

国中国国国国)\par

（3）两岸同胞同根同源、同文同种，是命运与共的骨肉兄弟，是血浓于水的一家人。\par

74）中华文化是两岸同胞共同的精神财富，也是两岸同胞血脉相连的精神纽带。\par

\textbf{怎样解决台湾问题，实现祖国的完全统一？}\par

“和平统一、一国两制”是解决台湾问题的基本方针，也是实现国家统一的最佳方式。\par

一个中国原则是两岸关系的政治基础，必须坚持“九二共识”、反对“台独"。\par

\daofasym{③}深化两岸融合发展，推进两岸经济合作制度化，打造两岸共同市场，夯实和平统一的基础。\par

\daofasym{①}两岸人民多走动、多交流、多沟通，增进理解、信任。\par

\daofasym{③}两岸同肥向根回源，同文回种，是命运与共的骨内元受，是血浓于血的一家大。要共同弘扬中华文化，实现心灵契合\par

增进对和平统一的认同。\par

【6+7用来回答台湾问题】\par

8.内实现祖国统一，青少年应怎么做？（开放题）\par

\daofasym{①}思想：树立远大理想，热爱祖国，勤奋学习，努力掌握建设祖国和保卫祖国的本领。\par

\daofasym{②}宣传：积极宣传我国解决台湾问题的方针、原则和立场。\par

\daofasym{③}行动：坚决同一切分裂祖国、破坏祖国统一的言行作斗争。\par


\section*{PDF第77页}
\addcontentsline{toc}{section}{PDF第77页}

\subsection*{第八课中国人中国梦 第一框我们的梦想}
\addcontentsline{toc}{subsection}{第八课中国人中国梦 第一框我们的梦想}

中国梦的内涵\par

规中华民族伟大复兴，就是中华民族近代以来最伟大的梦想。\par

现中国梦，就是要实现国家富强、民族振兴、人民幸福。\par

\textbf{为什么要实现中国梦？}\par

反映了近代以来一代又一代中国人的美好愿，揭示了中华民族的历史命运和当代中国的发展走向，指明了全国各族\par

民共同的奋斗目标。\par

本现子中华民族和中国人民的整体利益是国家的梦、民族的梦，也是每个中国人的梦。\par

一种题目：\par

\textbf{十么青少年要以实现中国梦为己任？}\par

因为中华民族伟大复兴的中国梦终将在一代代青年的接力奋斗中变为现实。\daofasym{②}\par

手一代有理想、有本领、有担当，国家就有前途，民族就有希望。\daofasym{③}青少年承载着国家和民\par

的未来命运。青少年的品格影响着国家未来发展。\daofasym{④}中国梦反映了近代以来一代又一代中\par

人的美好风原，揭示了中华民族的历史命运和当代中国的发展走向，指明了全国各族人民\par

同的奋斗目标。实现中件民族伟大复兴，体现了中华民族和中国人民的整体利益，是国家的\par

民族的梦，也是每个中国人的梦\par

中国特色社会主义进入新时代\par

1）新的历史方位：\par

经过长期努力，中国特色社会主义进入了新时代，这是我国发展新的历史方位。\par

2）指导思想：马克思列宁主义、毛泽东思想、邓小平理论、“三个代表”重要思想、科学发展观、习近平新时代中\par

特色社会主义思想（最新理论成果）。\par

3）新的战略目标一一“两个一百年”奋斗目标：\par

第一个百年：到建党一百年时（2020年），全面建成小康社会（全面小康）。\par

第二个百年：到新中国成立一百年时（2049年）（全面建成社会主义现代化强国。\par

实现第二个百年奋斗耳标一参键在党。党要继续弘扬坚持真理、坚守理想，践行初心、担当使命，不怕牺牲、英勇斗\par

对党忠诚、不负人民的伟大建党精神。\par

向第二个百年奋斗目标进军的两个阶段：\par

从2020年到2035年，在全面建成小康社会的基础上，基本实现社会主义现代化。\par

从2035年到本世纪中叶，在基本实现现代化的基础上，把我国建成富强民主文明和谐美丽的社会主义现代化强国。\par

4）中国特色社会主义进入新时代的历史意义（三个意味）：\par

意味着中华民族迎来了从站起来、富起来到强起来的伟大飞跃，迎来了实现中华民族伟大复兴的光明前景；\par

意味着科学社会主义在21世纪的中国焕发出强大生机活力，在世界上高高举起了中国特色社会主义伟大旗帜；\par

意味着中国特色社会主义道路、理论、制度、文化不断发展，为解决人类问题贡献了中国智慧和中国方案。\par

党提出的“三步走”战略目标（一般了解）\par

一步：从1981年到1990年，实现国民生产总值翻一番，解决人民的温饱问题。\par

二步：从1991年到20世纪末，国民生产总值再翻一番，人民生活达到小康水平（总体小康）。\par

三步：到21世纪中叶，人均国民生产总值达到中等发达国家水平，人民生活比较富裕，基本实现现代化。\par

第二框 共圆中国梦\par

\textbf{如何实现中华民族伟大复兴的中国梦？}\par

家层面：\daofasym{①}必须坚持党的领导，贯彻创新、协调、绿色、开放、共享的理念（五大发展理念）统筹推进经济建设、\par

治建设、文化建设、社会建设、生态文明建设（“五位一体”建设），协调推进全面建设社会主义现代化国家、全面\par

化改革、全面依法治国、全面从严治党（四个全面）。\par

必须走中国道路。中国道路就是中国特色社会主义道路。\par

必须弘扬中国精神。中国精神就是以爱国主义为核心的民族精神和以改革创新为核心的时代精神。\par

必须凝聚中国力量。中国力量就是全国各族人民大团结的力量。\par

人层面：中国梦的实现离不开每个人的奋斗。\par

中国自信的根本原因\par


\section*{PDF第78页}
\addcontentsline{toc}{section}{PDF第78页}

【结构整理：三步走战略和两个百年目标】\par

三步走：第一步（1980-1990）解决温饱；第二步（1990-2000）达到总体小康；第三步（2000-21世纪中叶）基本实现社会主义现代化，进而实现中华民族伟大复兴。\par

第一个百年目标（建党100年，2021）：全面建成小康社会。\par

第二个百年目标（建国100年，2049）：全面建成社会主义现代化强国。\par

第二个百年目标分两个阶段：2020-2035年基本实现社会主义现代化；2035年到本世纪中叶，把我国建设成富强民主文明和谐美丽的社会主义现代化强国。\par

\daofasym{①}中国共产党领导中国人民开辟了中国特色社会主义道路，形成了中国特色社会主义理论体系，确立了中国特色社会主\par

义制度，发展了中国特色社会主义文化，这是中国自信，民族自信的根本所在。\par

\daofasym{②}在党的领导下，中国特色社会主义伟大事业不断取得新的成就，国家富强、民族振兴让中国人更加自信。\par

3.自信的中国人的表现\par

\daofasym{①}对国家有认同。 \daofasym{②}对文化有底气。 \daofasym{③}对发展有信心。\par

\textbf{4.如何做自信的中国人？}\par

\daofasym{①}增强国家认同感，自觉维护国家利益和国家尊严。\par

\daofasym{②}坚守中华文化立场，讲好中国故事。\par

\daofasym{③}要坚持中国特色社会主义道路自信、理论自信、制度自信、文化自信。\par

不妄自尊大，也不故步自封,要培育理性平和、不卑不亢、开放包容的心态。\par

\daofasym{⑤}既是梦想家又是实干家既胸怀理想又求真务实既满怀激情又锲而不舍。\par

\textbf{5、个人梦与中国梦的关系?}\par

实现中华民族伟大复兴，体现了中华民族和中国人民的整体利益，是国家的梦、民\par

族的梦，也是每个中国人的梦。实现伟大梦想，需要我们凝心聚力，坚韧不拔、锲而不舍，奋力\par

开启时代新征程\par

温饱总体小康 基本实现社会\par

三步走 主义现代化 中华民族伟大复兴\par

第三步\par

第一步第二步\par

19801990 2000 2021 2035 2049\par

1（建党100年） （建国100年）\par

两个100年目标 全面建成 全面建成社会主\par

小康社会 义现代化强国\par

第一阶段 第二阶段\par

第二个百年目标 基本实现社会 把我国建设成富\par

主义现代化 强民主文明和谐\par

美丽的社会主义\par

现代化强国\par

【党的知识】\par

\daofasym{①}坚持以人民为中心的发展思想；\par

\daofasym{②}党的执政理念之一是共享发展成果；\par

\daofasym{③}中国共产党是中国特色社会主义建设事业的领导核心；代表了最广大人民的根本利益；\daofasym{④}党的宗旨是\par

全心全意为人民服务；\par

\daofasym{③}党的初心和使命是为中国人民谋幸福，为中华民族谋复兴；\par


\section*{PDF第79页}
\addcontentsline{toc}{section}{PDF第79页}

党的奋斗目标是人民对美好生活的向往。\par

中国共产党是中国工人阶级的先锋队，同时是中国人民和中华民族的先锋队\par

【民生框架】\par

有利于缩小贫富差距，实现城乡、区域协调发展，实现共同富裕。\par

有利于提高贫困地区人民的生活水平和获得感，满足人民日益增长的美好生活需要。\daofasym{③}有利于人民群\par

坚定中国特色社会主义道路自信、理论自信和制度自信。\daofasym{④}有利于使改革发展成果更多更公平地惠及\par

全体大民，促进社会公平正义。\daofasym{③}有利于巩固全面小康成果，构建社会主义和谐社会\par

发展的根本目的是增进民生福让党的知识\par

【成就框架、自信原因】\par

)坚持中国共产党的正确领导\daofasym{②}发挥社会主义制度集中力量办大事的优势\par

党和政府坚持以人民为中心的发展思想\par

党领导人民开辟了社会主义道路，形成了中国特色社会主义理论体系，确立了中国特色社会主义制度，\par

展了中国特色社会主义文化一（【根本原因\par

坚持了改革开放“\daofasym{③}坚持以经济建设为中心\par

弘扬中国精神，凝聚中国力量，全国人民团结奋斗\par

D坚持了（创新驱动发展战略、人才强国战略、等七个战略中的一些）\par

D坚持共享发展成果等等弘扬中华优秀传统文化，大力发展先进文化\par

【为什么要弘扬某种精神】\par

华民族精神在不同的历史时期有着不同的表现，“某某精神”是中华民族精神在新的时代条件下的集\par

体现。它是中华民族生生不息发展壮大的强大精神支柱，是团结奋斗的精神纽带，是实现中国梦的不\par

精神动力；它有利于激发爱国主义热情，构筑中国精神、中国价值、中国力量。新时代，我们要实现\par

会主义现代化和中华民族的伟大复兴，更要弘扬和培育“某某精神”。\par

【人物品质类】\par

物精神品质类框架\par

精神：民族精神，时代精神，创新精神，实干精神，敬业精神，工匠精神，劳模精神，劳动精神。脱\par

攻坚精神，航天精神，法治精神，伟大建党精神，伟大抗疫精神，人类命运共同体理念，社会主义核心\par

价值观，中华传统文化美德，\par

品质：高度的责任感，使命感，责任意识，爱国主义情怀。自信自强，行已有耻，止于至善。艰苦奋\par

，勇于实践。诚实守信，文明有礼，尊重他人。个人命运与国家命运，个人利益与国家利益。亲社会，\par

参与社会实践，积极参加社会劳动。自觉承担责任，不言代价与回报。爱岗敬业。关爱他人，服务社会，\par

无私奉献。敬畏生命，平凡中创造伟大，人生真正价值在于对社会责任与贡献，实现人生价值\par

旦遇到以下关键词，可以写以下知识点：\par

D遇到“总书记，中共中央”马上写中国共产党是中国特色社会主义事业的领导核心，坚持以人民为中心\par

的发展思想。\par

\daofasym{②}遇到“法”，马上写我国坚持依法治国的基本方略\par

\daofasym{③}遇到国家能解决某个问题（如疫情防控、抗震救灾等），可以写发挥社会主义制度集中力量办大事的优\par

势\par


\section*{PDF第80页}
\addcontentsline{toc}{section}{PDF第80页}

一旦遇到和人民有关的，都可以写人民是国家的主人，党和政府坚持以人民为中心的发展思想\par

（1）***关于第一单元《改革与创新》常用的句子：\par

1．坚持改革开放，全面深化改革，共享发展成果勃发房盼段：开会面连皮礼议机化化用有新\par

2.2.党和政府坚持以人民为中心的发展思想\par

3.3.创新是引领发展的第一动力\par

4.4.改革开放是强国之路\par

5．5.全面实施创新驱动发展战略，人才强国战略，科教兴国战略\par

6.6.全面推进大众创业，万众创新反吊裂法(s 年同年回有疑两次-作)\par

7．7.大力弘扬创新精神 1日统一七业\par

8．8.推进社会公平正义，建设社会主义和谐社会，实现共同富裕，全面小康\par

9.9.坚持党的正确领导，发挥社会主义制度集中力量办大事的优势\par

培图网改工作\par

（2）****关于第二单元《民主与法治》常用的句子：\par

1.社会主义民主是维护人民利益最广泛，最真实，最管用的民主\par

2.协商民主是社会主义民主政治的特有形式和独特优势 是\daofasym{②}有均\par

3.社会主义民主的本质是人民当家做主\par

4.民主决策是保障人民利益得到充分实现的有效方式\par

5.坚定不移的走中国特色社会主义法治国家道路，坚持当的领导，人民当家做主和依法治国的有机统一\par

6.建设社会主义法治国家，要坚持科学立法，公正司 -执法，全民守法(D 呈付均中期\par

（3）***关于第三单元《文明与家园》常用句子：\par

1．大力继承、弘扬和发展中华民族文化，坚定文化自信，坚定发展中国特色社会主义文化。净以币\par

2．2.大力弘扬民族精神，践行社会主义核心价值观 礼会！张物核研值观\par

3. 33.走绿色发展道路，建设生态文明，建设美丽中国。\par

4.坚持创新，协调，绿色，开放，共享发展理念，建设资源节约型，环境友好型社会。\par

5.有关环境发展的理念：绿色发展理念，可持续发展理念，绿水青山就是金山银山理念，人与自然和谐共生理念\par

6．我国资源日益短缺、环境污染严重，经济发展与资源、环境之间的矛盾日益突出，已经成为我国经济社会发展必须面\par

对的严峻挑战。\par

（4）***关于第四单元《和谐与梦想》常用句子：\par

1、坚持民族平等、民族团结、各民族共同繁荣的原则，巩固平等、团结、互助、和谐的新型民族关系，坚持民族区域自\par

治制度，巩固边疆，实现中华民族伟大复兴\par

2、加强和巩固民族团结，维护国家统一，是中华民族的最高利益\par

3、台湾问题的基本方针是“和平统一，一国两制”；港澳问题的基本方针是“一国两制”。国家统一的政治基础是一个\par

中国原则\par

4、中国梦一一实现中华民族的伟大复兴，实现国家富强、民族振兴、人民幸福\par

5、两个百年：第一个百年：到建党一百年时，全面建成小康社会；策二个百年：到新中国成立一百年时，全面建成社会\par

主义现代化强国。 年:0为可一轻统胖子可理\daofasym{②}理医系网络促媒获级I\par

两个阶段：从2020年到2035年，在全面建成小康社会的基础上，基本实现社会主义现代化：从2035年到本世纪中叶，\par

6、\daofasym{①}五大发展理念：创新、协调、绿色、开放、共享\daofasym{②}五位一体：经济建设、政治建设、文化建设、社会建设、生态建\par

设\daofasym{③}四个全面：全面建设社会主义现代化国家，全面深化改革，全面依法治国，全面从严治党\daofasym{④}个息信：道路自信、\par

理论自信、制度自信、文化自信\par

了大战增 图我吸秀效\par

学统许同的例\par


\section*{PDF第81页}
\addcontentsline{toc}{section}{PDF第81页}

【漏识补录：民生框架第\daofasym{④}\daofasym{⑤}点】\par

\daofasym{④}有利于发展成果更多更公平惠及全体人民，促进社会公平正义。\par

\daofasym{⑤}有利于巩固全面小康成果，构建社会主义和谐社会。\par

新发欣理气\par

19绿小青山足壹银山受史效：为中国中\par

展理袭丸上重点知识有奶继草卡精为全串皮全纹拟\par

党的知识】\par

坚持以人民为中心的发展思想；\par

党的执政理念之一是共享发展成果；\par

中国共产党是中国特色社会主义建设事业的领导核心；代表了最广大人民的根本利益；\daofasym{④}党的宗旨是\par

心全意为人民服务；\par

)党的初心和使命是为中国人民谋幸福，为中华民族谋复兴；\par

党的奋斗目标是人民对美好生活的向往。\par

中国共产党是中国工人阶级的先锋队，同时是中国人民和中华民族的先锋队(2中国共泛度9,我协明方作用日中因艺1周特包\par

【民生框架】\par

有利于缩小贫富差距，实现城乡、区域协调发展，实现共同富裕。物心昂很\par

D有利于提高贫困地区人民的生活水平和获得感，满足人民日益增长的美好生活需要。\daofasym{③}有利于人民群\par

坚定中国特色社会主义道路自信、理论自信和制度自信。\daofasym{④}有利于使改革发展成果更多更公平地惠及\par

)发展的根本目的是增进民生福让\daofasym{①}党的知识买月度\par

前利造或负指理1国择由本肾慧D 已生国成新\par

【成就框架、自信原因；同时可以用来回答怎么实现中国梦】\par

坚持中国共产党的正确领导\daofasym{②}发挥社会主义制度集中力量办大事的优势\par

党和政府坚持以人民为中心的发展思想\par

党领导人民开辟了社会主义道路，形成了中国特色社会主义理论体系，确立了中国特色社会主义制度，\par

展了中国特色社会主义文化一—【根本原因】\par

坚持了改革开放\daofasym{③}坚持以经济建设为中心\par

の弘扬中国精神，凝聚中国力量，全国人民团结奋斗\par

坚持了（创新驱动发展战略、人才强国战略、等七个战略中的一些）\par

坚持共享发展成果等等\daofasym{①}弘扬中华优秀传统文化，大力发展先进文化\par

1、中国人民的力量，中国人民是具有伟大创造精神、奋斗精神、团结精神、梦想精神的人民\par

【为什么要弘扬某种精神】\par

中华民族精神在不同的历史时期有着不同的表现；“某某精神”是中华民族精神在新的时代条件下的集\par

中体现。它是中华民族生生不息发展壮大的强大精神支柱，是团结奋斗的精神纽带，是实现中国梦的不\par

竭精神动力：它有利于激发爱国主义热情，构筑中国精神、中国价值、中国力量。新时代，我们要实现\par

社会主义现代化和中华民族的伟大复兴，更要弘扬和培育“某某精神”。\par

【人物品质类】\par

精神：民族精神，时代精神，创新精神，实干精神，敬业精神，工匠精神，劳模精神，劳动精神。脱\par

贫攻坚精神，航天精神，法治精神，伟大建党精神，伟大抗疫精神，人类命运共同体理念，社会主义核心\par


\section*{PDF第82页}
\addcontentsline{toc}{section}{PDF第82页}

价值观，中华传统文化美德，\par

2．品质：高度的责任感，使命感，责任意识，爱国主义情怀。自信自强，行已有耻，止于至善。艰苦奋\par

斗，勇于实。诚实守信，文明有礼，尊重他人。个人命运与国家命运，个人利益与国家利益。亲社会，\par

参与社会戏，积极参加社会剪动，自觉承担责任，不言代价与回报。爱岗敬业。关爱他人，服务社会，\par

无私奉献。敬畏生命，平凡中创造伟大，人生真正价值在于对社会责任与贡献，实现人生价值\par

斗m特生/a长精\par

【弘扬优秀传统文化的意义、保护文化遗产等的意义】\par

\daofasym{①}有利于传承源远流长(博大精深的中华文化，器现性文化一无二的理6细慧气度神期\par

\daofasym{②}有利于功强文化自信，动中华优秀传统文化实现创造性转化和创新性发展建设文化强用购\par

\daofasym{③}有利玉世强民族文化认同感之地强逸国情感一功强民证生力\par

\daofasym{①}有利于提开中年文化的国际影响力提高我国的文化软实力\par

\daofasym{③}有利于维护世界文化的多样性，促进世界文明交流互鉴与繁荣发展。\par

\daofasym{③}有利于弘扬中华民饶传统美德和民族精神，培育社会主义核心价值观。\par

【开展感动中国人物评选/学习英雄模范人物有什么意义？】\par

\daofasym{①}有利王物中华民游传统美德，弘扬伟大的民族精神，增强民族凝聚力：\par

\daofasym{②}有利于发挥模范人物的校样作用，传播正能量\par

\daofasym{③}道利于提高公民的愿想道德素励，践行社会主义核心价值观；\par

\daofasym{①}有利于形成良好的会风气，构建和谐社会\par

\daofasym{③}有利于为社会主义现代化建设提供强大的精神动力\par

【推行垃圾分类/建设“无废城市”的重要意义。强环境保护，有\par

1]样出大因寻但扰但等味计出代们,也任任件随价回授权沉办又第-致\par

什么意义？调，均发寻有\par

\daofasym{①}有利于落实节约资源，杀护环境的基木国策，实施可持续发展战略\par

\daofasym{②}有利于走绿色发展逆路，加快达设资源节约型、环境友好型社会；\par

\daofasym{③}有利于建设生态文明，实现人与自然的和谐共生\par

\daofasym{①}有利于处理好经济发展与环境保护的关系，建设美丽中国，实现中华民族的永续发展。\par

【了保护环境，国家、企业、青少年应该怎么做？】\par

（1）国家：政策、法律、科技、宜传\par

\daofasym{①}坚持节约资源和保护环境的基本国策，实施可持续发展战略，坚持走绿色发展道路：\par

\daofasym{②}坚持依法治国基本方略，严历打击各种浪费资源、破坏环境的违法行为；\par

（2）企业：道德、法律、科技\par

\daofasym{①}增强责任意识、环保意识，减少污染物的排放动\par

\daofasym{②}严格逆守国家的法律法规，依法规范自身的衍划\par

\daofasym{③}加快科技创新，改进节能环保技术，提高资源秘率； 磁中和因称\par

（3）青少年：思想上、宣传、行动上、作斗争、小事\par

\daofasym{①}学习环保知识，增强节约资源和保护环境的意识。\par

82 Z\par


\section*{PDF第83页}
\addcontentsline{toc}{section}{PDF第83页}

积极参加节约资源、保护环境的宣传活动，倡导低碳绿色生活方式；\par

积极同破坏环境的行为作斗争；\par

从身边的小事做起，顺手关水龙头、一水多用、人走灯灭、不乱扔垃圾等。\par

污哈[月:\par

【青少年怎么做类】学习+思想+行动\par

少年：\daofasym{①}树立崇高远太理想，发扬艰苦奋斗精神：\par

努力学习科学文化知识，培养创新能力\par

积极参加公益活动和社会实践活动，服务和奉献社会，增强社会责任感；至以因象，不41以展中图特\par

热爱祖国、报效祖国，勇敢面对学习和生活中的挫折，磨砺坚强意志；\par

弘扬和培育民族精神和时代精神，继承和引扬中华优秀传统文化，\par

自觉践行社会主义核心价值观，传递社会正能量： 网终订t境女孔礼会联样做\par

增强法治观念，尊法学法守法用法，依法维护国家利益，依法规范自身行为：\par

0积极向政府相关部门建言献策，为国家发展贡献力量；\par

从身边的小事做起，自觉做到节约资源、保护环境： 不引值\par

自觉承担实现中华民族伟大复兴的历史使命，为实现理想而努力奋斗，等等。\par

图以:侵如他右法女差导身V雄发展\par

【台湾问题】 不均年人正h 人 值观 世\par

\textbf{为什么祖国完全统一必定会实现？}\par

1）解决台湾问题，实现祖国完全统一，是全体中华儿女的共同愿望，是中华民族的根本利益所在。\par

（2）我国宪法规定，台湾是中华人民共和国的神圣领土的一部分。世界上只有一个中国，大陆和台湾同属一个中国。\par

【3）两岸同胞同根同源、同文同种，是命运与共的骨肉兄弟，是血浓于水的一家人。\par

4）中华文化是两岸同胞共同的精神财富，也是两岸同胞血脉相连的精神纽带。\par

\textbf{范样解决台湾问题，实现祖国的完全统一？}\par

）“和平统一、国两制”是解决台湾问题的基本方针，也是实现国家统一的最佳方式。\par

)一个中国原则是两岸关系的政治基础，必须坚持“九二共识”、反对台独"。\par

D深化两岸融合发展，推进两岸经济合作制度化，打造两岸共同市场，夯实和平统一的基础。\par

两岸人民多走动、多交流、多沟通，增进理解、信任。\par

两岸同胞同根同源，同文同种，是命运与共的骨肉兄弟，是血浓于血的一家人。要共同弘扬中华文化，实现心灵契合，\par

进对和平统一的认同。\par

习近计么了形调我现中年瓦滨伟大象关 公泌体与中同人既自 2 团线\par

0 人 是目多人 ,是b史包选看 左梦想6美坝有 L死醇 同者有力型,良后\par

目时是斗头中用多而是乏人民再中牛\par

九下[日若功实现多想 弟中许犹炸鸡\par

中日人)月头子骑群煌\par

\daofasym{①}音小年亢娱伟大复兴须疑梦中国\par

\subsection*{第一单元我们共同的世界（国与国） ，本的}
\addcontentsline{toc}{subsection}{第一单元我们共同的世界（国与国） ，本的}

经济全球化（趋势），世界多极化（趋势），文化多样性\par

当今世界是怎样的世界？（P3-4)\par


\section*{PDF第84页}
\addcontentsline{toc}{section}{PDF第84页}

【漏识补录：第7题文明交流互鉴第（4）点】\par

（4）学习和借鉴人类文明的一切优秀成果，不能只满足于欣赏物件的精美，更应该领略其中蕴含的人文精神。\par

K 一起向味来为中人美争运国件\par

\daofasym{①}这是一个开放的世界。国家间相互开放的程度不断加深，在政治、经济、文化各领域的开放范国也在不断扩展。\par

\daofasym{②}这是一个发展的世界。世界正经历着新一轮大发展大变革大调整。与此同时，贫富差距、发展不平衡等问题依然困扰\par

着人类社会。\par

\daofasym{③}这是一个紧密联系的世界。现代交通、通信、贸易把全球各地的国家、人们联系在一起，彼此影响，休戚相关。\par

2、经济全球化的表现：\daofasym{①}商品生产在全球范围内完成。\daofasym{②}商品贸易在全球范围内进行。\par

3.经济全球化的影响是什么？（怎样理解经济全球化是一把“双刃剑”？）（P5-7）\par

(1)积极影响：\daofasym{①}经济全球化改变了我们的生活：\par

\daofasym{②}经济全球化促进商品、资本和劳动力在全球流动，有利于在世界范围内配置资源，促进资源利用更加合理有效。\par

\daofasym{③}经济全球化也使各国经济相互联系、相互依赖的程度不断加深。为经济发展提供了新的机会。\par

(2)消极影响经济全球化也使风险与危机跨国界传递。\par

4、面对经济全球化，我们应该怎么做？（P7）\par

\daofasym{①}顺应历史潮流，保持积极、开放的心态，主动参与竞争；\par

\daofasym{②}居安思危，增强风险意识，注重国家经济安全，为应对各种困难和挑战做好充分准备。\par

5.文化多样性的原因、表现是什么？（P7）\par

(1)原因：各民族文化都有其产生和发展的渊源，表达了人们对世界特有的理解与情感。(2)表现：各民族文化共同构成媲\par

紫嫣红、生机盎然的世界文化大花园。\par

6.为什么要尊重文化多样性？（文化多样性的作用、意义）（P7-8）\par

\daofasym{①}文化多样性是人类社会的基本特征，是世界文化充满活力的表现，也是人类文明进步的重要动力。\par

\daofasym{②}多样的文化提供了更为广泛的选择，让人们的生活更加丰富多彩，让世界变得更加绚丽多姿。\par

\daofasym{③}文化多样性是实现文化创新与发展的前提和基础。不同特质的文化相互交融，能够为彼此增添新的元素，激发新的活\par

力。\par

\daofasym{③}每个民族的友化都是独特的，都有其存在的价值，都有值得尊重的经验和智慧。\par

\daofasym{③}各民族文化共同构成姹紫嫣红、生机盎然的文化大花园。\par

如何对待文化多样性，如何对待文化差异性又如何进行文明交流互鉴\par

用开放包容的心态，学习和借鉴人类文明的一切优秀成果，坚持以我为主，兼收并蓄。\par

积极王动地与世芥各国交往，从不同文明中寻求智慧、汲取营养。\par

（3且互尊重，通过平等交流、对话，达成彼此的理解与包容\par

8、国家间关系的特点（P12-14）【注意：世界多极化：影响国际关系的因素一国家利益+国家力量。任何国家都是以本\par

国的国家利益为出发点】\par

\daofasym{①}随着科学技术的进步，人类的脚步越走越远，国与国之间的交往越来越频繁，当今世界，越来越多国家在汲取历史经\par

验的基础上，不断探索新的互动方式。\par

当参国与国之间，既有竞争，也存合作。在经济全球化与世界多极化的进程中，各国在经济、政治、安全、文化等领\par

域深化合作，寻求发展。（合作发展）。\par

\daofasym{③}为今世界，国际的实质是以经流和科技买力为基础的综合国力的轻意，每个国家都希望通过自身努力在新一轮国\par

原关系的调整中获得发展机遇。国际竞争需要各方关同遵循一定的国际规则竞争）【\daofasym{②}和\daofasym{③}用来回答如何理解国家间\par

的竞争与合作】\par

\daofasym{③}尽管国际形势缓和的大趋势没有改变，但各种矛盾相互交织，各种力量重新分化纽合，各类冲突事件不断发生，考验\par

著国与国之间的关系，整个世界复杂多变。\par

9、我国怎样应对复杂多变的国际关系？我国高举和平、发展、合作、共离的旗帜，为推动建设相互奠重\textasciicircum{}公平正义\par

作共的新型国际关系作着不解的努力。\par

10、当今时代的主题是：和平与发展\par

当今时代潮流；和平、发展、合作、共赢\par

注意：世界整体上维持和平的态势，但战争的阴影从未远离。\par

\textbf{11、当令威胁世界和平的因素有哪些？}\par

局部战争与冲突从未间断，霸权主义、民族问题、宗教冲突、领土争端和恐怖主义。恐怖主义日益成为威\par

肋世界和平的重要因素。\par

\textbf{12维护和平的做法?}\par

\daofasym{①}世界：以和平谈判代替武力解决争端。成立联合国，不断建立健全维护和平的机制：派驻维和部队：签\par

署核不扩散条约··P18\par

我国：我国坚持对话协商、以和平手段解决国际争端和热点难题永远做维护世界和平不断进行新的探索。\par

办8\par


\section*{PDF第85页}
\addcontentsline{toc}{section}{PDF第85页}

我们要看到和平来之不易，珍借今天的和平，积极表达爱好和平的愿望，主动承担维护世界和平\par

为维护世界和平作出责献。P19\par

导致贫因的因素：\par

的战争；恶劣的自然环境：落后的教育。\par

\textbf{为什么要消除贫因？}\par

除资困、是当今人类必须共同面对的发展课题，PTS\par

上界经济发展不平衡现象仍然存在、发达国家与发展中国家间的差距仍然很大。P20\par

发的战争、恶劣的自然环境、落后的教育等诸多因素使许多国家仍难以摆脱贫困。P20\par

怎样消除贫因？【注意：我国已经消除了绝对贫国】\par

类社会对于消除贫困、促进发展的追求没有止境、在不同的发展阶段设立了递进式的减贫目标。\daofasym{②}各\par

结合自身的实际情况，不探索消除贫困的有效途径。\daofasym{③}面对问题，世界各国积极行动，为消除一切形\par

内贫困、改善教育和医疗卫生条件\par

人类命运共同体的基本内涵（内容）：【非常非常重要】\par

人和平、普遍安全、共同繁荣、开放包容、清洁美丽的世界。\par

\textbf{中国首倡构建人类命运共同体有何意义？}\par

国积极开展中国特色大国外交，推动建设新型国际关系，推动构建人类命运共同体，弘扬和平、发展\par

F、正义、民私、自由的全人类共同价值，引领人类进步潮流。\par

为什么集协建大类命运共同体？【构建人类命运共同理念的意义】\par

当今世界，客国相互联系、相互依存的程度空前加深。\par

人类面临许多共同挑战，需要解决许多全球性问题。\par

有哪个国家能多独自应对人类面临的各种挑战，采取共同行动、承担共同责任，构建人类命运共同体，\par

成为各国解决全球性问题的必然选择。\par

\textbf{怎样构津人类命运共同体？}\par

家角度：\daofasym{①}国要努力扩大利益交汇点，谁求开放创所、包容互惠的发展前景。\daofasym{②}客国坚持对话协商（政\par

、建设一个持人和平的世界，、坚持共建共商（安全）建设一个音普通安全的世界、、坚持合作共赢\par

经济）建设一个共同紧案的世界，通过文流互鉴（文化)建设一个开放包容的世界、坚持绿色低碳（生\par

建设一个美丽清洁的世界。。\par

民角度：\daofasym{①}相互信任、守望相助和共同担当。@关纤生命、尊重生命的价蓝不仅萃待自己一而且把关\par

的目光投向世界。\daofasym{③}关注他人命运，理解他人的室茎与快乐，看到生拿中比同的遇双，增通包容与合少\par

心年角度：\daofasym{①}肥把构建人类命运共同体的理念冠实际行动联系起来：\daofasym{②}关注世界的发展，关注大类的命运：\par

心系祖国，缺飞梦想，）在为人民利益的不态平中书写人生华置\par

我国为构建人类命运共同体做出不解努力，举例子\par

动一带一路建设；积极推动区域发展。\par

\subsection*{第二单元世界舞台上的中国（中国与世界）}
\addcontentsline{toc}{subsection}{第二单元世界舞台上的中国（中国与世界）}

定一定要背下的包子：\par

中国着眼于时代发展大势？遵循共商共建共享原则，为全球治理提出中国方家\par

献中国智慧，展现大国担当\par

我国坚持并倡导构建人类命运共同理念，弘扬和平、发展、公平正义自由民主的\par

人类共同价值\par

作为一个负责伊的大国，中国努力提高自身在国际上的影响力、（感召力和塑造力D\par

力于成为世界和平的建设者、全球发展的贡献者、国际秩序的维护者\par

(0 例外区会斑动难题中\par

【遇到中国在国际上做的事情都可以用这几句话】 同如终乳圾主的承但贵任，即\par

因归可性高因展机;用、甚影发展成关\daofasym{①}糊积极开展中月特色月\par

18 交,推功津技新型间阶天南,\par


\section*{PDF第86页}
\addcontentsline{toc}{section}{PDF第86页}

【漏识补录：第2题第（5）点开头】\par

（5）中国在推动国际秩序朝着更加公正合理的方向发展，更好维护我国和广大发展中国家共同利益的同时，坚持以经济建设为中心。\par

第一框：中国担当\par

\textbf{1.我国的基本国情是什么？面对国情国家应该怎么做？}\par

我国仍处于并将长期处于社会主义初级阶段的基本国情没有变，这需要我们全面深化改革，统筹国际国内\par

两个大局，走和平发展道路，奉行互利共赢的开放战略。\par

2．中国在国际事务中的中国智慧、中国担当、中国方案\par

(1)面对各种区域性和全球性的危机与难题，中国不推，不逃避，也不依赖他人，积极主动地承担起相\par

应的责任。\par

随着国际参与程度的不断加深，中国为国际社会各种难题与危机的化解作出了巨大贡献。中国担当向世界\par

展现了大国风范，显示了中国智慧。\par

(2)中国的脱贫攻坚取得全面胜利，创造了减贫治理的中国样本，为全球减贫事业做出重大贡献。\par

（3）中国着眼于时代发展大势，遵循共商共建共享原则，为全球治理提出中国方案，贡献中国智慧。\par

(4）中国广泛参与国际事务，在承担责任中不断积累经验，提升能力，增长智慧。同时，在解决人类面临\par

的各种问题的过程中积极探索、有效行动，发挥负责任大国作用，促进人类社会共同发展。\par

坚持以经济建设为中心，集中力量办好自己的事情，不断增强我们在国际上说话办事的实力及国际影响力。\par

我们积极参与全球治理，主动承担国际责任，既尽力而为，又量力而行。\par

（6）作为一个负责任的大国，中国努力提高自身在国际上的影响力，感召力和塑造力，中国始终是世界\par

和平的建设者、全球发展的贡献者、国际秩序的维护者。\par

（7）我国坚持人类命运共同体理念，不仅关注国内人民，而且把关切目光投向世界，关心人类共同命运。\par

\textbf{3.中国对世界的影响具体有哪些方面？}\par

(1)随着中国的发展，中国文化对世界的影响越来越大。在中国与其他国家的深度互动中，外\par

国人有更多的机会、更强烈的意愿，接触、体验和认识中国文化，感受中国文化的魅力，接受\par

中国文化的熏陶。\par

(2)中国正为世界经济增长注入新的活力。中国日益成为世界经济发展的引擎与稳定器\par

(3)中国秉承“和而不同”的思想，以及共商共建共享的全球治理观，倡导国际关系民主化、\par

构建国际新秩序、使世界向覆公平公正、多元共治、包容有序的局面发展。\par

(4)中国关于构建全球治理体系的探索与实践，为人类思考与建设未来提供了新的路径得到\par

广泛认同，越来越多的国家参与其中，共同行动，这将对世界的和平与发展产生深远的影响，\par

4,文明交流互鉴有什么意义？（文明兼收并蓄、交流互鉴的原因）（P38-39)\par

（中华文明在交流互鉴中发展。中国积极主动地与世界各国交往，从不同文明中寻求智慧、取营养，\par

不仅有助于自身文明的发展，而且能够推动世界文明的进步，与其他文明携手解决人类共同面临的各种问\par

题。\par

(2)文明因交流而多彩，文明因互鉴而丰富。\par

（3）通过精神的交流互鉴，为人类社会发展提供精神支撑和心灵慰藉。\par

5.我们应该如何对待其他文明？（怎样进行文明交流互鉴？）（P37-39)\par

（1）用开放包容的心态，学习和借鉴人类文明的一切优秀成果，坚持以我为主，兼收并蓄。\par

（2）积极主动地与世界各国交往，从不同文明中寻求智慧、汲取营养。\par

（3）相互尊重，通过平等交流、对话，达成彼此的理解与包容\par

（4）对其他文明的学习，我们不能只满足于欣赏物件的精美，更应该领略其中蕴含的人文精神。\par


\section*{PDF第87页}
\addcontentsline{toc}{section}{PDF第87页}

\textbf{6、在文化交流中，我们青少年应该如何做一名友好往来的使者？}\par

(1)要注意学习，提高自身的文化素养和文化交流的本领。\par

(2)热情友善，以礼相待，不卑不亢。\par

[3)尊重外国友人的生活方式和风俗习惯。\par

[4)积极宜传、弘扬我们民族的优秀文化。\par

7、我国面临哪些发展机遇？（P41-42）一抓住机退\par

内部环境：（1）我国已转向高质量发展阶段，制度优势显著，治理效能提升，经济长期向好，物质基础雄厚，人力资源\par

丰富，市场空间广阔，发展韧性强劲，社会大局稳定，经济发展具有多方面的优势和条件。\par

（2）经过几十年的发展，中国已在资金、人才、技术、管理经验、基础设施等领域具备良好的积累，为推动中国制\par

造向中国创造转变，中国速度向中国质量转变，中国产品向中国品牌转变莫定了良好的基础。\par

外部环境：（3）和平、发展、合作、共赢的时代潮流越来越强劲，为我国的发展提供了良好的外部环境。\par

（4）中国的国际影响力不断提升，引领作用越来越大。中国在国际合作各个领域获得更大发展空间，更加有所作为。\par

B、我国面临哪些风险和挑战？（P43-45）\par

1)近年来，受全球经济大环境的影响，中国经济面临一定的下行压力和不少困难，几十年高速发展所积累的一些矛盾与\par

风险也逐渐暴露出来，急需得到稳妥处理与解决。\par

[2)企业劳动力成本上升，“中国制造”在新的历史条件下需要转型升级。\par

[3)复杂多变的国际形势，给中国“走出去”发展战略带来发度机遇的同时，也带来了各种挑战。\par

\textbf{9、面对机遇与挑战，我们应该怎么做？}\par

增强机遇意识和风险意识，立足社会主义初级阶段基本国情，认识和把握发展规律，发扬斗争精神，树立\par

底线思维，科学应变、主动求变，抓住机遇，应对挑战，趋利避害，奋勇前进。\par

10、为什么要积极谋求发展？（谋求发展的原因）（必要性）\par

\daofasym{①}当今世界，科学技术日新月异，全球竞争不断升级。\daofasym{②}时不我待，中国要把握世界的发展趋势，积极谋求自身发展，\par

提高国际竞争力。 六质果发展：发展是进造业拒功能统方丝代[针级\par

\textbf{11、我国应该采取哪些措施积极谋求发展？}\par

(1）促进发展，要把提升发展质量放在首位。\par

(2)促进发展，要积极寻求新的经济增长点：培育壮大经济发展新动能。\par

(3）促进发展，要以更加开放的态度积极参与全球规则制定\par

(4)促进发展，要加快构建以国内大循环为主体、国内国际双循环相互促进的新发展格局。\par

\textbf{12、我国积极参与全球规则制定有什么意义？}\par

\daofasym{①}有利于我国在国际竞争中掌握话语权，保障对外经贸利益，进一步拓展国际市场。\par

\daofasym{②}我国正积极表达、多方参与，通过全球规则的制定和修改推动建立国际经济新秩序。\par

13、中国是怎样与世界各国共享发展机遇的？/中国的发展与世界发展的关系怎样？\par

（1)在谋求自身发展的同时，坚持合作共赢的理念，让发展的成果更多更公平地惠及各个国家。\par

(2)重视与相关国家和地区的合作，致力于共同建设一个繁荣的世界。\par

(3)中国的发展为世界各国提供了更广阔的市场、更充足的资本、更丰富的产品、更宝贵的合作契机。中国以“-一带一路\par

“为契机，把我国同沿线国家发展结合起来，开展跨国互联互通，实现优势互补，互利共赢。\par

4)中国与世界各国分享发展机遇，共享发展成果。中国人民愿同各国人民一道，共同开辟人类更加繁荣，更加安宁美好\par

的未来。\par

（5）中国坚持胸怀天下，始终以世界眼光关注人类前途命运，从人类发展大潮流、世界变化大格局、中国发展大历史，\par

正确认识和处理同外部世界的关系。\par

\subsection*{第三单元走向未来的少年（少年与世界）}
\addcontentsline{toc}{subsection}{第三单元走向未来的少年（少年与世界）}

，我们为什么要为世界添光彩？（个人与世界的关系）\par

D我们与丰富多彩的世界紧密相连，始终与这个世界彼此互动，同呼吸，共命运。\par

\daofasym{②}每个人都是世界的一员，所做的事都有可能对世界发展产生影响。\par

为什么说少年强、中国强？（青少年为什么要有担当？）\par

0个人的命运与国家的命运息息相关，个人的未来与民族的未来紧密相连。\par

\daofasym{②}青年兴则国家兴，青年强则国家强。青年一代有理想、有本领、有担当，国家就有前途，民族就有希望。\par

\daofasym{③}青少年承载若国家和民族的未来命运。青少年的品格影响着国家未来发展，\par

.青少年应如何承担时代赋予的贵任，走向世界？（我们应有怎样的情怀和抱负？也可以用来回答怎么做自信中国人）\par

D我们要树立全球观念，树立国际意识，关注世界发展，关心人类命运。\par


\section*{PDF第88页}
\addcontentsline{toc}{section}{PDF第88页}

【结构整理：四大模块之心理模块】\par

青春：青春生理变化；认识自我；接纳欣赏自我（做更好的自己）；青春心理矛盾（三种表现：反抗与依赖、闭锁与开放、勇敢与怯懦）；青春情绪情感；青春的三种思维；青春飞扬（自信、自强）；青春有格（行己有耻、止于至善）。\par

梦想：梦想的作用；少年与梦；如何实现梦想。\par

学习：学习的特点、学习的重要性、如何学习。（七上）\par

\daofasym{②}坚定理想信念，志存高远，脚踏实地，勇做时代的弄潮儿，努力在实现中国梦的伟大实践中建功立业，创造精彩的人\par

生。\par

\daofasym{③}要传承、弘扬中华优秀传统文化，增强爱国情感，弘扬民族精神和时代精神，践行社会主义核心价值观，并将其转化\par

为自己的情感认同和行为习惯，做有自信、懂自尊、能自强的中国人，成为中华民族的栋梁。\par

\daofasym{③}要了解人类文明进程，积极关切人类问题和世界局势，掌握相应的知识，在与世界各国青少年交流中提高我们的影响。\par

同时，要有忧患意识、居安思危，在世界复杂多变的环境中提高改变世界的素质和能力。\par

\daofasym{③}要尊重差异、理解不同、包容多样文化，向国际社会讲好中国故事、传播好中国声音、阐发中国精神、展现中国风貌，\par

承担起推动人类共同发展的责任。\par

\textbf{4. 为什么要在实践中学习？}\par

\daofasym{①}实践出真知。在实践中，提高认识世界、改造世界的能力。\par

\daofasym{②}在实践中，我们锤炼自己，丰富人生经历，完善自我，提升自身素质。\par

\daofasym{③}学习是一个长期的过程，学校学习只是其中的一个阶段，不可急功近利，也不能错失良机。\par

\daofasym{③}青少年随着生活范围的扩大会遇到越来越多的困惑，要在实践中逐步解决，以促进自身的发展和提升\par

\daofasym{①}我们要重视实践，积极参加社会调查、志愿服务、科学实验等各类社会实践活动，增强问题意识，培养研究能力，努\par

力做到知行合一。\par

\daofasym{②}必须树立终身学习的理念，养成主动学习、不断探索的习惯，增强自我更新、学以致用的能力。\par

\textbf{6.我们要怎样做出正确的选择？}\par

（1）要理性分析主客观条件，思考并决定自己的目标和方向。（学会合理选择方法：\daofasym{①}弄清自己的真实需要。\daofasym{②}多方面\par

收集信息。\daofasym{③}掌握选择的策略。）\par

（2）我们要系好人生的第一粒扣子，自觉践行社会主义核心价值观，走好这重要的一步，用勤劳和汗水开辟人生和事业\par

的美好前程。\par

四大模块\par

心理模块\par

（青春\par

\daofasym{①}青春生理变化 认识自我 接纳欣赏自我（做更好的自己\par

\daofasym{②}青春心理矛盾向三,自季,同爱\par

\daofasym{③}青春情绪情感\par

\daofasym{①}青春的三种思维\par

\daofasym{③}青春飞扬自信自强 t\par

\daofasym{③}青春有格（行己有耻，止于至善）\par

(2)梦想<梦想作用，少年与梦，如何实现梦想）\par

（3）学习学习的特点，重要性、如何学习）七上\par


\section*{PDF第89页}
\addcontentsline{toc}{section}{PDF第89页}

【结构整理：心理模块续】\par

生命问题。\par

【结构整理：道德模块】\par

人际交往中的道德品质：与同学（友谊问题、异性交往、网络交友）；与老师（师生交往）；与家人（孝亲敬长、和谐家庭）；与集体（与集体共成长、建设美好集体等）。\par

与社会：现实社会（亲社会行为、个人与社会）；虚拟社会（网络）。\par

社会交往需要具备的道德：文明有礼、诚实守信、尊重他人、做负责任的人、不言代价与回报、关爱他人、服务社会、奉献精神等。\par

【结构整理：法律模块】\par

\daofasym{①}规则与秩序；\daofasym{②}法律的特征和作用；\daofasym{③}对未成年人的法律保护、六大保护、依法办事；\daofasym{④}做守法公民（违法行为、犯罪、善用法律）；\daofasym{⑤}宪法；\daofasym{⑥}权利与义务；\daofasym{⑦}法治精神（自由、平等、公平、正义）；\daofasym{⑧}民主与法治。\par

4生命问题\par

道德模块人际交往中的道德品质）\par

与同学：友谊问题异性交往，网络交友\par

与老师：师生交往\par

与家人孝亲敬长、和谐家庭\par

\daofasym{①}与集体：与集体共成长、建设美好集体等\par

与社会\par

（现实社会一一亲社会行为、个人与社会；Tsm 800 枚\par

虚拟社会——网络；\par

社会交往需要具备的道德文明有礼诚实守信，尊重他人、做负责任的人\par

不言代价与回报、关爱他人、服务社会奉献精神(因服第人 m多周糖列\par

法律模块\par

同韩川,女为物，类坡后学鼠精摊\par

\daofasym{①}规则与秩序\par

\daofasym{②}法律的特征作用\par

\daofasym{③}对未成年人的法律保护、六大保护、依法办事\par

\daofasym{④}做守法公民（违法行为、犯罪、善用法律）天精神·特别能艺音\par

特剂骼好斗\par

宪法\par

\daofasym{③}权利与义务\par

\daofasym{⑦}法治精神（自由平等公平正义）\par

\daofasym{③}民主与法治\par


\section*{PDF第90页}
\addcontentsline{toc}{section}{PDF第90页}

\begin{note}[手写批注]
页面中下部有大段手写整理；（此处有手写笔记但无法识别）。\par
\end{note}

【结构整理：四大模块之国情模块】\par

国情模块（经济、政治、文化、社会、生态、国际）。\par

\daofasym{①}经济：富强与创新（改革开放、共享发展成果、创新）；经济制度；关心国家发展（我国的成就、劳动成就今天，实干创造未来）。（八上）\par

\daofasym{②}政治：政治制度；国家机构。\par

\daofasym{③}文化：守望精神家园（中华文化、传统美德、社会主义核心价值观、中国精神）。\par

\daofasym{④}社会：和谐与梦想（国家统一、民族团结、中国特色社会主义新时代、中国梦）。\par

\daofasym{⑤}生态：美丽中国；人口、环境、资源。\par

\daofasym{⑥}世界：维护国家利益（安全）；构建人类命运共同体；中国担当、中国智慧、中国方案；机遇与挑战；走向未来的少。\par


\end{document}
